%!TEX root = std.tex
\rSec0[quantities]{Quantities and units library}

\rSec1[quantities.summary]{Summary}

\pnum
This Clause describes components for dealing with quantities and units,
as summarized in \tref{quantities.summary}.

\begin{modularlibsumtab}{Quantities and units library summary}{quantities.summary}
\ref{qty.utils}         & Utilities                              & \tcode{mp_units.core} \\
\ref{qty.ref}           & Reference                              & \\
\ref{qty.rep}           & Representation                         & \\
\ref{qty}               & Quantity                               & \\ \rowsep
\ref{qty.systems}       & Systems                                & \tcode{mp_units.systems} \\
\ref{qty.chrono}        & \tcode{std::chrono} interoperability   & \\
\end{modularlibsumtab}

\ednote{
Following the SG16 recommendation at \url{https://lists.isocpp.org/sg16/2024/10/4490.php},
the \fakegrammarterm{universal-character-name}s should be replaced by their UTF-8 code points.
}

\rSec1[mp.units.syns]{mp-units module synopses}

\rSec2[mp.units.syn]{Module \tcode{mp_units} synopsis}

\indexmodule{mp_units}%
\begin{codeblock}
export module mp_units;

export import mp_units.core;
export import mp_units.systems;
\end{codeblock}

\rSec2[mp.units.core.syn]{Module \tcode{mp_units.core} synopsis}
\indexmodule{mp_units.core}%
\begin{codeblock}
// mostly freestanding
export module mp_units.core;

import std;

export namespace std {

// \ref{qty.utils}, utilities

// \ref{qty.sym.txt}, symbol text

enum class @\libglobal{character_set}@ { unicode, basic };

template<std::size_t N, std::size_t M>
class symbol_text;

// \ref{qty.sym.expr}, symbolic expressions

// \ref{qty.sym.expr.types}, types

template<typename T, typename... Ts>
struct @\exposidnc{per}@;  // \expos

template<typename F, int Num, int... Den>
  requires @\seebelownc@
struct @\exposidnc{power}@;  // \expos

// \ref{qty.ref}, reference

// \ref{qty.dim}, dimension

// \ref{qty.dim.concepts}, concepts

template<typename T>
concept dimension = @\seebelownc@;

// \ref{qty.dim.types}, types

template<symbol_text Symbol>
struct base_dimension;

template<@\exposconceptnc{SymbolicConstant}@... Expr>
struct derived_dimension;

struct dimension_one;
inline constexpr dimension_one @\libglobal{dimension_one}@{};

// \ref{qty.dim.ops}, operations

consteval @\libconcept{dimension}@ auto inverse(@\libconcept{dimension}@ auto d);

template<std::intmax_t Num, std::intmax_t Den = 1, @\libconcept{dimension}@ D>
  requires(Den != 0)
consteval @\libconcept{dimension}@ auto pow(D d);
consteval @\libconcept{dimension}@ auto sqrt(@\libconcept{dimension}@ auto d);
consteval @\libconcept{dimension}@ auto cbrt(@\libconcept{dimension}@ auto d);

// \ref{qty.dim.sym.fmt}, symbol formatting

struct @\libglobal{dimension_symbol_formatting}@ {
  character_set char_set = character_set::unicode;
};

template<typename CharT = char, std::@\stdconcept{output_iterator}@<CharT> Out, @\libconcept{dimension}@ D>
constexpr Out dimension_symbol_to(Out out, D d, const dimension_symbol_formatting& fmt = {});

template<dimension_symbol_formatting fmt = {}, typename CharT = char, @\libconcept{dimension}@ D>
consteval std::string_view dimension_symbol(D);

// \ref{qty.spec}, quantity specification

// \ref{qty.spec.concepts}, concepts

template<typename T>
concept quantity_spec = @\seebelownc@;

template<typename T, auto QS>
concept quantity_spec_of = @\seebelownc@;

// \ref{qty.spec.types}, types

// \ref{named.qty}, named

struct subkind;
inline constexpr subkind @\libglobal{subkind}@{};

template<auto...>
struct named_quantity_spec;  // \notdef

template<@\exposconceptnc{BaseDimension}@ auto Dim, @\exposconceptnc{QSProperty}@ auto... Args>
struct named_quantity_spec<Dim, Args...>;

template<@\exposconceptnc{DerivedQuantitySpec}@ auto Eq, @\exposconceptnc{QSProperty}@ auto... Args>
struct named_quantity_spec<Eq, Args...>;

template<@\exposconceptnc{NamedQuantitySpec}@ auto QS, @\exposconceptnc{QSProperty}@ auto... Args>
struct named_quantity_spec<QS, Args...>;

template<@\exposconceptnc{NamedQuantitySpec}@ auto QS, @\exposconceptnc{DerivedQuantitySpec}@ auto Eq, @\exposconceptnc{QSProperty}@ auto... Args>
struct named_quantity_spec<QS, Eq, Args...>;

// \ref{derived.qty}, derived

template<@\exposconceptnc{SymbolicConstant}@... Expr>
struct derived_quantity_spec;

// \ref{dimless.qty}, base quantity of dimension one

struct dimensionless;
inline constexpr dimensionless @\libglobal{dimensionless}@{};

// \ref{kind.of.qty}, kind of

template<@\libconcept{quantity_spec}@ Q>
  requires @\seebelownc@
struct kind_of_;
template<@\libconcept{quantity_spec}@ auto Q>
  requires requires { typename kind_of_<decltype(Q)>; }
inline constexpr kind_of_<decltype(Q)> @\libglobal{kind_of}@{};

// \ref{delta.qty}, delta

template<@\libconcept{quantity_spec}@ Q>
  requires @\seebelownc@
struct @\exposidnc{delta-spec-t}@;  // \expos

// \ref{point.qty}, point

template<typename...>
struct point_;  // \notdef

template<@\libconcept{quantity_spec}@ Q>
  requires @\seebelownc@
struct point_<Q>;

template<@\libconcept{quantity_spec}@ Q, @\exposconceptnc{PointOrigin}@ PO>
  requires @\seebelownc@
struct point_<Q, PO>;

// \ref{magnitude.qty}, magnitude

template<@\libconcept{quantity_spec}@ Q>
  requires @\seebelownc@
struct @\exposidnc{magnitude-t}@;  // \expos

// \ref{qty.spec.ops}, operations

template<@\libconcept{quantity_spec}@ Q>
  requires(!@\exposconceptnc{PointQuantitySpec}@<Q>) && @\exposconceptnc{ScalarOperand}@<Q>
consteval @\libconcept{quantity_spec}@ auto inverse(Q q);

template<std::intmax_t Num, std::intmax_t Den = 1, @\libconcept{quantity_spec}@ Q>
  requires(Den != 0) && (!@\exposconceptnc{PointQuantitySpec}@<Q>) && (@\exposconceptnc{ScalarOperand}@<Q> || Num == Den)
consteval @\libconcept{quantity_spec}@ auto pow(Q q);
template<@\libconcept{quantity_spec}@ Q>
  requires(!@\exposconceptnc{PointQuantitySpec}@<Q>) && @\exposconceptnc{ScalarOperand}@<Q>
consteval @\libconcept{quantity_spec}@ auto sqrt(Q q);
template<@\libconcept{quantity_spec}@ Q>
  requires(!@\exposconceptnc{PointQuantitySpec}@<Q>) && @\exposconceptnc{ScalarOperand}@<Q>
consteval @\libconcept{quantity_spec}@ auto cbrt(Q q);

template<@\libconcept{quantity_spec}@ Q>
  requires(!@\exposconceptnc{PointQuantitySpec}@<Q>) && @\exposconceptnc{NonScalarOperand}@<Q>
consteval @\libconcept{quantity_spec}@ auto magnitude(Q q);

// \ref{qty.spec.hier.algos}, hierarchy algorithms

// \ref{qty.get.kind}, \tcode{get_kind}

template<@\libconcept{quantity_spec}@ Q>
consteval @\seebelownc@ get_kind(Q);

// \ref{get.common.qty.spec}, \tcode{get_common_quantity_spec}

consteval @\libconcept{quantity_spec}@ auto get_common_quantity_spec(@\libconcept{quantity_spec}@ auto... qs)
  requires @\seebelownc@;

// \ref{qty.unit}, unit

// \ref{qty.unit.mag}, magnitude

// \ref{qty.unit.mag.concepts}, concepts

template<typename T>
concept MagConstant = @\seebelownc@;

template<typename T>
concept unit_magnitude = @\seebelownc@;

// \ref{qty.unit.mag.types}, types

template<symbol_text Symbol, long double Value>
  requires(Value > 0)
struct mag_constant;

// \ref{qty.unit.mag.ops}, operations

template<@\exposconceptnc{MagArg}@ auto V>
constexpr @\libconcept{unit_magnitude}@ auto mag = @\seebelownc@;

template<std::intmax_t N, std::intmax_t D>
  requires(N > 0)
constexpr @\libconcept{unit_magnitude}@ auto mag_ratio = @\seebelownc@;

template<@\exposconceptnc{MagArg}@ auto Base, int Num, int Den = 1>
constexpr @\libconcept{unit_magnitude}@ auto mag_power = @\seebelownc@;

// constants

inline constexpr struct @\libglobal{pi}@ final :
    mag_constant<{u8"\u03C0" /* @\unicode{03c0}{GREEK SMALL LETTER PI}@ */, "pi"},
                 std::numbers::pi_v<long double>> {
} @\libglobal{pi}@;

inline constexpr auto \u03C0 /* @\unicode{03c0}{GREEK SMALL LETTER PI}@ */ = pi;

// \ref{qty.unit.traits}, traits

template<@\libconcept{unit}@ auto U>
constexpr bool space_before_unit_symbol = true;

template<>
inline constexpr bool @\libspec{space_before_unit_symbol}{one}@<one> = false;

// \ref{qty.unit.concepts}, concepts

template<typename T>
concept unit = @\seebelownc@;

template<typename T>
concept prefixable_unit = @\seebelownc@;

template<typename U, auto QS>
concept unit_of = @\seebelownc@;

// \ref{qty.unit.types}, types

// \ref{qty.scaled.unit}, scaled

template<@\libconcept{unit_magnitude}@ auto M, @\libconcept{unit}@ U>
  requires @\seebelownc@
struct scaled_unit;

// \ref{qty.named.unit}, named

template<symbol_text Symbol, auto...>
struct named_unit;  // \notdef

template<symbol_text Symbol, @\exposconceptnc{QuantityKindSpec}@ auto QS>
  requires @\seebelownc@
struct named_unit<Symbol, QS>;

template<symbol_text Symbol, @\exposconceptnc{QuantityKindSpec}@ auto QS, @\exposconceptnc{PointOrigin}@ auto PO>
  requires @\seebelownc@
struct named_unit<Symbol, QS, PO>;

template<symbol_text Symbol, @\libconcept{unit}@ auto U>
  requires @\seebelownc@
struct named_unit<Symbol, U>;

template<symbol_text Symbol, @\libconcept{unit}@ auto U, @\exposconceptnc{PointOrigin}@ auto PO>
  requires @\seebelownc@
struct named_unit<Symbol, U, PO>;

template<symbol_text Symbol, @\libconcept{unit}@ auto U, @\exposconceptnc{QuantityKindSpec}@ auto QS>
  requires @\seebelownc@
struct named_unit<Symbol, U, QS>;

template<symbol_text Symbol, @\libconcept{unit}@ auto U, @\exposconceptnc{QuantityKindSpec}@ auto QS,
         @\exposconceptnc{PointOrigin}@ auto PO>
  requires @\seebelownc@
struct named_unit<Symbol, U, QS, PO>;

// \ref{qty.prefixed.unit}, prefixed

template<symbol_text Symbol, @\libconcept{unit_magnitude}@ auto M, @\libconcept{prefixable_unit}@ auto U>
  requires @\seebelownc@
struct prefixed_unit;

// \ref{qty.common.unit}, common

template<@\libconcept{unit}@ U1, @\libconcept{unit}@ U2, @\libconcept{unit}@... Rest>
struct common_unit;

// \ref{qty.derived.unit}, derived

template<@\exposconceptnc{SymbolicConstant}@... Expr>
struct derived_unit;

// \ref{qty.unit.one}, one

struct one;
inline constexpr one @\libglobal{one}@{};

// named derived units of a quantity of dimension one

inline constexpr struct @\libglobal{percent}@ final : named_unit<"%", mag_ratio<1, 100> * one> {
} @\libglobal{percent}@;

inline constexpr struct @\libglobal{per_mille}@ final :
    named_unit<symbol_text{u8"\u2030" /* @\unicode{2030}{PER MILLE SIGN}@ */, "%o"},
               mag_ratio<1, 1000> * one> {
} @\libglobal{per_mille}@;

inline constexpr struct @\libglobal{parts_per_million}@ final :
    named_unit<"ppm", mag_ratio<1, 1'000'000> * one> {
} @\libglobal{parts_per_million}@;

inline constexpr auto @\libglobal{ppm}@ = parts_per_million;

// \ref{qty.unit.ops}, operations

consteval @\libconcept{unit}@ auto inverse(@\libconcept{unit}@ auto u);

template<std::intmax_t Num, std::intmax_t Den = 1, @\libconcept{unit}@ U>
  requires @\seebelownc@
consteval @\libconcept{unit}@ auto pow(U u);
consteval @\libconcept{unit}@ auto sqrt(@\libconcept{unit}@ auto u);
consteval @\libconcept{unit}@ auto cbrt(@\libconcept{unit}@ auto u);
consteval @\libconcept{unit}@ auto square(@\libconcept{unit}@ auto u);
consteval @\libconcept{unit}@ auto cubic(@\libconcept{unit}@ auto u);

// \ref{qty.unit.cmp}, comparison

template<@\libconcept{unit}@ From, @\libconcept{unit}@ To>
consteval bool convertible(From from, To to);

// \ref{qty.unit.obs}, observers

consteval @\libconcept{quantity_spec}@ auto get_quantity_spec(@\libconcept{unit}@ auto u);
consteval @\libconcept{unit}@ auto get_unit(@\libconcept{unit}@ auto u);

// \ref{qty.unit.sym.fmt}, symbol formatting

enum class @\libglobal{unit_symbol_solidus}@ {
  one_denominator,
  always,
  never
};

enum class @\libglobal{unit_symbol_separator}@ {
  space,
  half_high_dot
};

struct @\libglobal{unit_symbol_formatting}@ {
  character_set char_set = character_set::unicode;
  unit_symbol_solidus solidus = unit_symbol_solidus::one_denominator;
  unit_symbol_separator separator = unit_symbol_separator::space;
};

template<typename CharT = char, std::@\stdconcept{output_iterator}@<CharT> Out, @\libconcept{unit}@ U>
constexpr Out unit_symbol_to(Out out, U u, const unit_symbol_formatting& fmt = {});

template<unit_symbol_formatting fmt = {}, typename CharT = char, @\libconcept{unit}@ U>
consteval std::string_view unit_symbol(U);

// \ref{qty.ref.concepts}, concepts

template<typename T>
concept quantity_reference = @\seebelownc@;

template<typename T, auto QS>
concept quantity_reference_of = @\seebelownc@;

// \ref{qty.ref.syn}, class template \tcode{reference}

template<@\libconcept{quantity_spec}@ Q, @\libconcept{unit}@ U>
struct reference;

// \ref{qty.ref.ops}, operations

template<typename FwdRep, @\exposconceptnc{BareReference}@ R,
         @\libconcept{representation_of}@<get_quantity_spec(R{})> Rep = std::remove_cvref_t<FwdRep>>
  requires(!@\exposconceptnc{OffsetUnit}@<decltype(get_unit(R{}))>)
constexpr quantity<R{}, Rep> operator*(FwdRep&& lhs, R r);

template<typename FwdRep, @\exposconceptnc{BareReference}@ R,
         @\libconcept{representation_of}@<get_quantity_spec(R{})> Rep = std::remove_cvref_t<FwdRep>>
  requires(!@\exposconceptnc{OffsetUnit}@<decltype(get_unit(R{}))>)
constexpr @\libconcept{delta_quantity}@ auto operator/(FwdRep&& lhs, R);

template<typename FwdQ, @\exposconceptnc{BareReference}@ R, @\libconcept{delta_quantity}@ Q = std::remove_cvref_t<FwdQ>>
constexpr @\libconcept{delta_quantity}@ auto operator*(FwdQ&& q, R);

template<typename FwdQ, @\exposconceptnc{BareReference}@ R, @\libconcept{delta_quantity}@ Q = std::remove_cvref_t<FwdQ>>
constexpr @\libconcept{delta_quantity}@ auto operator/(FwdQ&& q, R);

template<@\exposconceptnc{BareReference}@ R, typename Rep>
  requires @\libconcept{representation_of}@<std::remove_cvref_t<Rep>, get_quantity_spec(R{})>
constexpr auto operator*(R, Rep&&) = delete;

template<@\exposconceptnc{BareReference}@ R, typename Rep>
  requires @\libconcept{representation_of}@<std::remove_cvref_t<Rep>, get_quantity_spec(R{})>
constexpr auto operator/(R, Rep&&) = delete;

template<@\exposconceptnc{BareReference}@ R, typename Q>
  requires @\libconcept{delta_quantity}@<std::remove_cvref_t<Q>>
constexpr auto operator*(R, Q&&) = delete;

template<@\exposconceptnc{BareReference}@ R, typename Q>
  requires @\libconcept{delta_quantity}@<std::remove_cvref_t<Q>>
constexpr auto operator/(R, Q&&) = delete;

// \ref{qty.ref.obs}, observers

template<typename Q, typename U>
consteval @\libconcept{quantity_spec}@ auto get_quantity_spec(reference<Q, U>);

template<typename Q, typename U>
consteval @\libconcept{unit}@ auto get_unit(reference<Q, U>);

consteval @\libconcept{unit}@ auto get_common_reference(@\libconcept{unit}@ auto u1, @\libconcept{unit}@ auto u2,
                                         @\libconcept{unit}@ auto... rest)
  requires @\seebelownc@;

template<@\exposconceptnc{BareReference}@ R1, @\exposconceptnc{BareReference}@ R2, @\exposconceptnc{BareReference}@... Rest>
consteval @\libconcept{quantity_reference}@ auto get_common_reference(R1 r1, R2 r2, Rest... rest)
  requires @\seebelownc@;

// \ref{qty.rep}, representation

enum class @\libglobal{quantity_character}@ { scalar, complex, vector, tensor };

// \ref{qty.rep.traits}, traits

// \ref{qty.fp.traits}, floating-point

template<typename Rep>
constexpr bool treat_as_floating_point = @\seebelownc@;

// \ref{qty.char.traits}, quantity character

template<typename T>
constexpr bool disable_scalar = false;
template<>
inline constexpr bool @\libspec{disable_scalar}{bool}@<bool> = true;
template<>
inline constexpr bool @\libspec{disable_scalar}{char}@<char> = true;
template<>
inline constexpr bool @\libspec{disable_scalar}{wchar_t}@<wchar_t> = true;
template<>
inline constexpr bool @\libspec{disable_scalar}{char8_t}@<char8_t> = true;
template<>
inline constexpr bool @\libspec{disable_scalar}{char16_t}@<char16_t> = true;
template<>
inline constexpr bool @\libspec{disable_scalar}{char32_t}@<char32_t> = true;
template<typename T>
constexpr bool @\libspec{disable_scalar}{std::complex}@<std::complex<T>> = true;

template<typename T>
constexpr bool disable_complex = false;

template<typename T>
constexpr bool disable_vector = false;

// \ref{qty.val.traits}, values

template<typename Rep>
struct representation_values;

// \ref{qty.rep.cpos}, customization point objects

namespace math {
  inline namespace @\unspec@ {
    inline constexpr @\unspec@ real = @\unspec@;
    inline constexpr @\unspec@ imag = @\unspec@;
    inline constexpr @\unspec@ modulus = @\unspec@;

    inline constexpr @\unspec@ magnitude = @\unspec@;
  }
}

// \ref{qty.rep.concepts}, concepts

template<typename T>
concept Representation = @\seebelownc@;

template<typename T, quantity_character Ch>
concept representation_of = @\seebelownc@;

// \ref{qty}, quantity

// \ref{qty.like}, interoperability

template<typename T>
struct quantity_like_traits;  // \notdef

// \ref{qty.syn}, class template \tcode{quantity}

template<typename T>
concept delta_quantity = @\seebelownc@;

template<typename T>
concept point_quantity = @\seebelownc@;

template<typename Q, auto V>
concept quantity_of = @\seebelownc@;

template<@\libconcept{quantity_reference}@ auto R, @\libconcept{representation_of}@<get_quantity_spec(R)> Rep = double>
class quantity;

// \ref{qty.delta}, variable template \tcode{delta}

template<@\exposconceptnc{BareReference}@ R>
struct delta_;

template<auto Arg>
  requires @\seebelownc@
constexpr auto @\libglobal{delta}@ = @\seebelownc@;

// \ref{qty.non.mem.conv}, non-member conversions

template<@\libconcept{unit}@ auto ToU, @\seebelownc@>
  requires @\seebelownc@
constexpr @\libconcept{delta_quantity}@ auto value_cast(@\seebelownc@ q);

template<@\libconcept{Representation}@ ToRep, @\seebelownc@>
  requires @\seebelownc@
constexpr quantity<@\seebelownc@, ToRep> value_cast(@\seebelownc@ q);

template<@\libconcept{unit}@ auto ToU, @\libconcept{Representation}@ ToRep, @\seebelownc@>
  requires @\seebelownc@
constexpr @\libconcept{delta_quantity}@ auto value_cast(@\seebelownc@ q);
template<@\libconcept{Representation}@ ToRep, @\libconcept{unit}@ auto ToU, @\seebelownc@>
  requires @\seebelownc@
constexpr @\libconcept{delta_quantity}@ auto value_cast(@\seebelownc@ q);

template<@\libconcept{delta_quantity}@ ToQ, @\seebelownc@>
  requires @\seebelownc@
constexpr @\libconcept{delta_quantity}@ auto value_cast(@\seebelownc@ q);

template<@\libconcept{quantity_spec}@ auto ToQS, @\seebelownc@>
  requires @\seebelownc@
constexpr @\libconcept{delta_quantity}@ auto quantity_cast(@\seebelownc@ q);

template<@\libconcept{unit}@ auto ToU, @\seebelownc@>
  requires @\seebelownc@
constexpr @\libconcept{point_quantity}@ auto value_cast(@\seebelownc@ qp);

template<@\libconcept{Representation}@ ToRep, @\seebelownc@>
  requires @\seebelownc@
constexpr quantity<@\seebelownc@, ToRep> value_cast(@\seebelownc@ qp);

template<@\libconcept{unit}@ auto ToU, @\libconcept{Representation}@ ToRep, @\seebelownc@>
  requires @\seebelownc@
constexpr @\libconcept{point_quantity}@ auto value_cast(@\seebelownc@ qp);
template<@\libconcept{Representation}@ ToRep, @\libconcept{unit}@ auto ToU, @\seebelownc@>
  requires @\seebelownc@
constexpr @\libconcept{point_quantity}@ auto value_cast(@\seebelownc@ qp);

template<@\libconcept{delta_quantity}@ ToQ, @\seebelownc@>
  requires @\seebelownc@
constexpr @\libconcept{point_quantity}@ auto value_cast(@\seebelownc@ qp);

template<@\libconcept{point_quantity}@ ToQP, @\seebelownc@>
  requires @\seebelownc@
constexpr @\libconcept{point_quantity}@ auto value_cast(@\seebelownc@ qp);

template<@\libconcept{quantity_spec}@ auto ToQS, @\seebelownc@>
  requires @\seebelownc@
constexpr @\libconcept{point_quantity}@ auto quantity_cast(@\seebelownc@ qp);

}

// \ref{qty.common.type}, \tcode{std::common_type} specializations

template<std::@\libconcept{delta_quantity}@ Q1, std::@\libconcept{delta_quantity}@ Q2>
  requires @\seebelownc@
struct std::common_type<Q1, Q2>;

template<std::@\libconcept{delta_quantity}@ Q, std::@\libconcept{Representation}@ Value>
  requires @\seebelownc@
struct std::common_type<Q, Value>;

template<std::@\libconcept{delta_quantity}@ Q, std::@\libconcept{Representation}@ Value>
  requires requires { typename std::common_type<Q, Value>; }
struct @\libspec{std::common_type}{quantity}@<Value, Q> : std::common_type<Q, Value> {};

export namespace std {

// \ref{qty.pt.orig}, point origin

// \ref{qty.pt.orig.concepts}, concepts

template<typename T, auto QS>
concept point_origin_for = @\seebelownc@;

// \ref{qty.pt.orig.types}, types

// \ref{qty.abs.pt.orig}, absolute

template<@\libconcept{quantity_spec}@ auto QS>
struct absolute_point_origin;

// \ref{qty.rel.pt.orig}, relative

template<@\libconcept{point_quantity}@ auto QP>
struct relative_point_origin;

// \ref{qty.natural.pt.orig}, natural

template<@\libconcept{quantity_spec}@ auto QS>
  requires @\seebelownc@
struct natural_point_origin_;

template<@\libconcept{quantity_spec}@ auto QS>
constexpr natural_point_origin_<@\exposidnc{remove-point}@(@\exposidnc{remove-delta}@(QS))> @\libglobal{natural_point_origin}@{};

// \ref{qty.def.pt.orig}, default

template<@\exposconceptnc{BareReference}@ R>
consteval @\libconcept{point_origin_for}@<get_quantity_spec(R{})> auto default_point_origin(R);

// \ref{qty.point}, point

template<@\exposconceptnc{BareReference}@ R, @\exposconceptnc{PointOrigin}@ PO>
struct point_<R, PO>;

template<auto Arg, auto... Rest>
  requires @\seebelownc@
constexpr auto @\libglobal{point}@ = @\seebelownc@;

}
\end{codeblock}

\rSec2[mp.units.systems.syn]{Module \tcode{mp_units.systems} synopsis}
\indexmodule{mp_units.systems}%
\begin{codeblock}
export module mp_units.systems;

export import mp_units.core;
import std;

export namespace std {

// \ref{qty.chrono}, \tcode{std::chrono} interoperability

template<typename Rep, typename Period>
struct quantity_like_traits<std::chrono::duration<Rep, Period>>;

template<typename Clock>
struct chrono_point_origin_;
template<typename Clock>
constexpr chrono_point_origin_<Clock> @\libglobal{chrono_point_origin}@{};

template<typename Clock, typename Rep, typename Period>
struct quantity_like_traits<
  std::chrono::time_point<Clock, std::chrono::duration<Rep, Period>>>;

}
\end{codeblock}

\rSec1[qty.utils]{Utilities}

\rSec2[qty.utils.non.types]{Non-types}

\indexlibraryglobalexpos{is-specialization-of}
\indexlibraryglobalexpos{is-derived-from-specialization-of}
\begin{itemdecl}
template<typename T, template<@\seebelow@> typename U>
consteval bool @\exposidnc{is-specialization-of}@();               // \expos
template<typename T, template<@\seebelow@> typename U>
consteval bool @\exposidnc{is-derived-from-specialization-of}@();  // \expos
\end{itemdecl}

\begin{itemdescr}
\pnum
\returns
\begin{itemize}
\item
For the first signature,
\tcode{true} of \tcode{T} is a specialization of \tcode{U}, and
\tcode{false} otherwise.
\item
For the second signature,
\tcode{true} if \tcode{T} has exactly one public base class
that is a specialization of \tcode{U}
and has no other base class that is a specialization of \tcode{U}, and
\tcode{false} otherwise.
\end{itemize}

\pnum
\remarks
An implementation provides enough overloads for all arguments to \tcode{U}.
\end{itemdescr}

\rSec2[qty.ratio]{Ratio}

\indexlibraryglobalexpos{ratio}
\indexlibrarymemberexpos{operator==}{ratio}
\indexlibrarymemberexpos{operator<=>}{ratio}
\indexlibrarymemberexpos{operator-}{ratio}
\indexlibrarymemberexpos{operator+}{ratio}
\indexlibrarymemberexpos{operator/}{ratio}
\indexlibraryexposmemberexpos{is-integral}{ratio}
\begin{codeblock}
namespace std {

struct @\exposidnc{ratio}@ {  // \expos
  std::intmax_t num;
  std::intmax_t den;

  consteval @\exposidnc{ratio}@(std::intmax_t n, std::intmax_t d = 1);

  friend consteval bool operator==(@\exposidnc{ratio}@, @\exposidnc{ratio}@) = default;
  friend consteval auto operator<=>(@\exposidnc{ratio}@ lhs, @\exposidnc{ratio}@ rhs) { return (lhs - rhs).num <=> 0; }

  friend consteval @\exposidnc{ratio}@ operator-(@\exposidnc{ratio}@ r) { return {-r.num, r.den}; }

  friend consteval @\exposidnc{ratio}@ operator+(@\exposidnc{ratio}@ lhs, @\exposidnc{ratio}@ rhs)
  {
    return {lhs.num * rhs.den + lhs.den * rhs.num, lhs.den * rhs.den};
  }

  friend consteval @\exposidnc{ratio}@ operator-(@\exposidnc{ratio}@ lhs, @\exposidnc{ratio}@ rhs) { return lhs + (-rhs); }

  friend consteval @\exposidnc{ratio}@ operator*(@\exposidnc{ratio}@ lhs, @\exposidnc{ratio}@ rhs);

  friend consteval @\exposidnc{ratio}@ operator/(@\exposidnc{ratio}@ lhs, @\exposidnc{ratio}@ rhs)
  {
    return lhs * @\exposidnc{ratio}@{rhs.den, rhs.num};
  }
};

consteval bool @\exposidnc{is-integral}@(@\exposidnc{ratio}@ r) { return r.num % r.den == 0; }

consteval @\exposidnc{ratio}@ @\exposidnc{common-ratio}@(@\exposidnc{ratio}@ r1, @\exposidnc{ratio}@ r2);

}
\end{codeblock}

\pnum
\exposid{ratio} represents the rational number $\tcode{num}/\tcode{den}$.

\pnum
Unless otherwise specified,
in the following descriptions,
let \tcode{R(r)} be \tcode{std::ratio<N, D>},
where \tcode{N} and \tcode{D} are the values of \tcode{r.num} and \tcode{r.den}.

\indexlibraryctorexpos{ratio}
\begin{itemdecl}
consteval @\exposidnc{ratio}@(std::intmax_t n, std::intmax_t d = 1);
\end{itemdecl}

\begin{itemdescr}
\pnum
Let \tcode{N} and \tcode{D} be the values of \tcode{n} and \tcode{d}.
Let \tcode{R} be \tcode{std::ratio<N, D>}.

\pnum
\effects
Equivalent to
\tcode{R}.

\pnum
\ensures
\tcode{num == R::num \&\& den == R::den} is \tcode{true}.
\end{itemdescr}

\indexlibrarymemberexpos{operator*}{ratio}
\begin{itemdecl}
friend consteval @\exposidnc{ratio}@ operator*(@\exposidnc{ratio}@ lhs, @\exposidnc{ratio}@ rhs);
\end{itemdecl}

\begin{itemdescr}
\pnum
Let \tcode{Res} be \tcode{std::ratio_multiply<R(lhs), R(rhs)>}.

\pnum
\effects
Equivalent to:
\tcode{return \{Res::num, Res::den\};}
\end{itemdescr}

\indexlibraryexposmemberexpos{common-ratio}{ratio}
\begin{itemdecl}
consteval @\exposidnc{ratio}@ @\exposidnc{common-ratio}@(@\exposidnc{ratio}@ r1, @\exposidnc{ratio}@ r2);
\end{itemdecl}

\begin{itemdescr}
\pnum
Let \tcode{Res} be equal to
\begin{codeblock}
std::common_type<std::chrono::duration<int, R(r1)>,
                 std::chrono::duration<int, R(r2)>>::type::period
\end{codeblock}

\pnum
\effects
Equivalent to:
\tcode{return \{Res::num, Res::den\};}
\end{itemdescr}

\rSec2[qty.sym.txt]{Symbol text}

\indexlibraryglobal{symbol_text}
\begin{codeblock}
namespace std {

template<std::size_t N, std::size_t M>
class symbol_text {
public:
  std::fixed_u8string<N> @\exposidnc{utf8}@;    // \expos
  std::fixed_string<M> @\exposidnc{portable}@;  // \expos

  // constructors
  constexpr symbol_text(char portable);
  consteval symbol_text(const char (&portable)[N + 1]);
  constexpr symbol_text(const std::fixed_string<N>& portable);
  consteval symbol_text(const char8_t (&utf8)[N + 1], const char (&portable)[M + 1]);
  constexpr symbol_text(const std::fixed_u8string<N>& utf8,
                        const std::fixed_string<M>& portable);

  // observers
  constexpr const auto& @\libmember{utf8}{symbol_text}@() const { return @\exposidnc{utf8}@; }
  constexpr const auto& @\libmember{portable}{symbol_text}@() const { return @\exposidnc{portable}@; }
  constexpr bool @\libmember{empty}{symbol_text}@() const { return utf8().empty(); }

  // string operations
  template<std::size_t N2, std::size_t M2>
  friend constexpr symbol_text<N + N2, M + M2> operator+(const symbol_text& lhs,
                                                         const symbol_text<N2, M2>& rhs);

  // comparison
  template<std::size_t N2, std::size_t M2>
  friend constexpr bool operator==(const symbol_text& lhs,
                                   const symbol_text<N2, M2>& rhs) noexcept;
  template<std::size_t N2, std::size_t M2>
  friend constexpr auto operator<=>(const symbol_text& lhs,
                                    const symbol_text<N2, M2>& rhs) noexcept;
};

symbol_text(char) -> symbol_text<1, 1>;

template<std::size_t N>
symbol_text(const char (&)[N]) -> symbol_text<N - 1, N - 1>;

template<std::size_t N>
symbol_text(const std::fixed_string<N>&) -> symbol_text<N, N>;

template<std::size_t N, std::size_t M>
symbol_text(const char8_t (&)[N], const char (&)[M]) -> symbol_text<N - 1, M - 1>;

template<std::size_t N, std::size_t M>
symbol_text(const std::fixed_u8string<N>&, const std::fixed_string<M>&) -> symbol_text<N, M>;

}
\end{codeblock}

\pnum
\tcode{symbol_text} represents a symbol text.
\exposid{utf8} stores its UTF-8 representation, and
\exposid{portable} stores its portable representation.
\tcode{symbol_text<N, M>} is a structural type\irefcppx{temp.param}{term.structural.type}.

\pnum
In the descriptions that follow,
it is a \Fundescx{Precondition} that
\begin{itemize}
\item
values of \tcode{char} are in the basic literal character set\irefcpp{lex.charset}, and
\item
for a parameter of the form \tcode{const \placeholdernc{CharT} (\&\placeholdernc{txt})[\placeholder{M}]},
\tcode{(\placeholdernc{txt}[\placeholdernc{M} - 1] == \placeholdernc{CharT}())} is \tcode{true}.
\end{itemize}

\indexlibraryctor{symbol_text}
\begin{itemdecl}
constexpr symbol_text(char portable);
consteval symbol_text(const char (&portable)[N + 1]);
constexpr symbol_text(const std::fixed_string<N>& portable);
consteval symbol_text(const char8_t (&utf8)[N + 1], const char (&portable)[M + 1]);
constexpr symbol_text(const std::fixed_u8string<N>& utf8, const std::fixed_string<M>& portable);
\end{itemdecl}

\begin{itemdescr}
\pnum
For the constructors without a parameter named \tcode{utf8},
let \tcode{utf8} be:
\begin{codeblock}
std::bit_cast<std::fixed_u8string<N>>(std::basic_fixed_string(portable))
\end{codeblock}

\pnum
\effects
Equivalent to the \fakegrammarterm{mem-initializer-list}:
\begin{codeblock}
@\exposid{utf8}@{utf8}, @\exposidnc{portable}@{portable}
\end{codeblock}
\end{itemdescr}

\indexlibrarymember{operator+}{symbol_text}
\begin{itemdecl}
template<std::size_t N2, std::size_t M2>
friend constexpr symbol_text<N + N2, M + M2> operator+(const symbol_text& lhs,
                                                       const symbol_text<N2, M2>& rhs);
\end{itemdecl}

\begin{itemdescr}
\pnum
\effects
Equivalent to:
\begin{codeblock}
return symbol_text<N + N2, M + M2>(lhs.utf8() + rhs.utf8(),
                                   lhs.portable() + rhs.portable());
\end{codeblock}
\end{itemdescr}

\indexlibrarymember{operator==}{symbol_text}
\indexlibrarymember{operator<=>}{symbol_text}
\begin{itemdecl}
template<std::size_t N2, std::size_t M2>
friend constexpr bool operator==(const symbol_text& lhs,
                                 const symbol_text<N2, M2>& rhs) noexcept;
template<std::size_t N2, std::size_t M2>
friend constexpr auto operator<=>(const symbol_text& lhs,
                                  const symbol_text<N2, M2>& rhs) noexcept;
\end{itemdecl}

\begin{itemdescr}
\pnum
Let \tcode{@} be the \fakegrammarterm{operator}.

\pnum
\effects
Equivalent to:
\begin{codeblock}
return std::make_tuple(std::cref(lhs.utf8()), std::cref(lhs.portable())) @\atsign@
       std::make_tuple(std::cref(rhs.utf8()), std::cref(rhs.portable()));
\end{codeblock}
\end{itemdescr}

\rSec2[qty.sym.expr]{Symbolic expressions}

\rSec3[qty.sym.expr.general]{General}

\pnum
Subclause \ref{qty.sym.expr} specifies the components
used to maintain ordered, simplified, and readable
argument lists in the names of specializations.
\begin{example}
\begin{codeblock}
using namespace si::unit_symbols;
int x = kg * km / square(h);  // error: cannot construct from
  // \tcode{derived_unit<si::kilo_<si::gram>, si::kilo_<si::metre>, \exposidnc{per}<\exposidnc{power}<non_si::hour, 2>>>}
\end{codeblock}
The library ensures \tcode{decltype(kg * km / square(h))} is styled-like as commented in diagnostics,
provided that, in the implementation-defined total order of types,
\tcode{decltype(kg)} is less than \tcode{decltype(km)}.
\end{example}

\rSec3[qty.sym.expr.concepts]{Concept \ecname{SymbolicConstant}}

\begin{itemdecl}
template<typename T>
concept @\defexposconceptnc{SymbolicConstant}@ =  // \expos
  std::is_empty_v<T> && std::is_final_v<T> && std::is_trivially_default_constructible_v<T> &&
  std::is_trivially_copy_constructible_v<T> && std::is_trivially_move_constructible_v<T> &&
  std::is_trivially_destructible_v<T>;
\end{itemdecl}

\pnum
The concept \exposconcept{SymbolicConstant}
is used to constrain the types
that are used in symbolic expressions.

\rSec3[qty.sym.expr.types]{Types}

\indexlibraryglobalexpos{per}
\begin{codeblock}
namespace std {

template<typename T, typename... Ts>
struct @\exposidnc{per}@ final {};  // \expos

}
\end{codeblock}

\pnum
\exposid{per} is used to store arguments with negative exponents.
A specialization of \exposid{per}
represents the product of the inverse of its template arguments.
A program that instantiates a specialization of \exposid{per}
that is not a possible result of the library specifications
is ill-formed, no diagnostic required.

\indexlibraryglobalexpos{power}
\begin{codeblock}
namespace std {

template<typename F, int Num, int... Den>
  requires @\seebelownc@
struct @\exposidnc{power}@ final {  // \expos
  using @\exposidnc{factor}@ = F;                              // \expos
  static constexpr @\exposidnc{ratio}@ @\exposidnc{exponent}@{Num, Den...};  // \expos
};

}
\end{codeblock}

\pnum
\exposid{power} represents a power\irefiev{102-02-08}
of the form $\tcode{F}^{\tcode{Num}/\tcode{Den}}$.
\begin{note}
\tcode{Den} is optional to shorten the type name when \tcode{Den} is \tcode{1}.
\end{note}
A program that instantiates a specialization of \exposid{power}
that is not a possible result of the library specifications
is ill-formed, no diagnostic required.

\pnum
Let \tcode{r} be \tcode{\exposidnc{ratio}\{Num, Den...\}}.
Let \tcode{\placeholder{is-valid-ratio}} be
\tcode{true} if \tcode{r} is a valid constant expression, and
\tcode{false} otherwise.
The expression in the \fakegrammarterm{requires-clause} is equivalent to:
\begin{codeblock}
@\placeholdernc{is-valid-ratio}@ && (r > @\exposidnc{ratio}@{0}) && (r != @\exposidnc{ratio}@{1})
\end{codeblock}

\rSec3[qty.sym.expr.algos]{Algorithms}

\indexlibraryglobalexpos{expr-type}
\begin{codeblock}
template<typename T>
using @\exposidnc{expr-type}@ = @\seebelownc@;  // \expos
\end{codeblock}

\pnum
\tcode{\exposidnc{expr-type}<T>} denotes
\tcode{U} if \tcode{T} is of the form \tcode{\exposidnc{power}<U, Ints...>}, and
\tcode{T} otherwise.

\indexlibraryglobalexpos{type-less-impl}
\begin{itemdecl}
template<typename T, typename U>
consteval bool @\exposidnc{type-less-impl}@();  // \expos
\end{itemdecl}

\begin{itemdescr}
\pnum
\returns
\tcode{true} if \tcode{T} is less than \tcode{U}
in an implementation-defined total order for types, and
\tcode{false} otherwise.
\end{itemdescr}

\indexlibraryglobalexpos{type-less}
\begin{codeblock}
template<typename Lhs, typename Rhs>
struct @\exposidnc{type-less}@ :  // \expos
    std::bool_constant<@\exposidnc{is-specialization-of}@<Rhs, @\exposidnc{power}@>() ||
                       @\exposidnc{type-less-impl}@<@\exposidnc{expr-type}@<Lhs>, @\exposidnc{expr-type}@<Rhs>>()> {};
\end{codeblock}

\pnum
\exposid{type-less} meets the requirements of
the \tcode{Pred} parameter of the symbolic expression algorithms below.

\indexlibraryglobalexpos{type-list}
\indexlibraryglobalexpos{expr-fractions}
\begin{codeblock}
template<typename... Ts>
struct @\exposidnc{type-list}@ {};      // \expos

template<typename OneType, typename... Ts>
struct @\exposidnc{expr-fractions}@ {   // \expos
  using @\exposidnc{num}@ = @\seebelownc@;  // \expos
  using @\exposidnc{den}@ = @\seebelownc@;  // \expos
}
\end{codeblock}

\pnum
\exposid{expr-fractions} divides a symbolic expression to numerator and denominator parts.
Let \tcode{EF} be a specialization of \exposid{expr-fractions}.
\begin{itemize}
\item
If \tcode{EF} is of the form \tcode{\exposidnc{expr-fractions}<OneType, Ts..., \exposidnc{per}<Us...>>},
then
\begin{itemize}
\item
\tcode{EF::\exposidnc{num}} denotes \tcode{\exposidnc{type-list}<Ts...>}, and
\item
\tcode{EF::\exposidnc{den}} denotes \tcode{\exposidnc{type-list}<Us...>}.
\end{itemize}
\item
Otherwise, \tcode{EF} is of the form \tcode{\exposidnc{expr-fractions}<OneType, Ts...>}, and
\begin{itemize}
\item
\tcode{EF::\exposidnc{num}} denotes \tcode{\exposidnc{type-list}<Ts...>}, and
\item
\tcode{EF::\exposidnc{den}} denotes \tcode{\exposidnc{type-list}<>}.
\end{itemize}
\end{itemize}

\pnum
The symbolic expression algorithms perform operations on symbolic constants.
A symbolic constant is a type that is a model of \exposconcept{SymbolicConstant}.
\begin{example}
The dimension \tcode{dim_length}, the quantity \tcode{time}, and the unit \tcode{one} are symbolic constants.
\end{example}
The algorithms also support
powers with a symbolic constant base and a rational exponent,
products thereof, and
fractions thereof.

\indexlibraryglobalexpos{expr-multiply}
\indexlibraryglobalexpos{expr-divide}
\indexlibraryglobalexpos{expr-invert}
\indexlibraryglobalexpos{expr-pow}
\begin{itemdecl}
template<template<typename...> typename To, typename OneType,
         template<typename, typename> typename Pred = @\exposidnc{type-less}@, typename Lhs, typename Rhs>
consteval auto @\exposidnc{expr-multiply}@(Lhs, Rhs);        // \expos

template<template<typename...> typename To, typename OneType,
         template<typename, typename> typename Pred = @\exposidnc{type-less}@, typename Lhs, typename Rhs>
consteval auto @\exposidnc{expr-divide}@(Lhs lhs, Rhs rhs);  // \expos

template<template<typename...> typename To, typename OneType, typename T>
consteval auto @\exposidnc{expr-invert}@(T);                 // \expos

template<std::intmax_t Num, std::intmax_t Den, template<typename...> typename To,
         typename OneType, template<typename, typename> typename Pred = @\exposidnc{type-less}@, typename T>
  requires(Den != 0)
consteval auto @\exposidnc{expr-pow}@(T);                    // \expos
\end{itemdecl}

\begin{itemdescr}
\pnum
\mandates
\begin{itemize}
\item
\tcode{OneType} is the neutral element\irefiev{102-01-19} of the operation, and
\item
\tcode{Pred} is a \oldconcept{BinaryTypeTrait}\irefcpp{meta.rqmts}
with a base characteristic of \tcode{std::bool_constant<B>}.
\tcode{Pred<T, U>} implements a total order for types;
\tcode{B} is \tcode{true} if \tcode{T} is ordered before \tcode{U}, and \tcode{false} otherwise.
\end{itemize}

\pnum
\effects
\pnum
First, inputs to the operations are obtained from the types of the function parameters.
If the type of a function parameter is:
\begin{itemize}
\item
A specialization of \tcode{To},
then its input is the product of its template arguments, and
the following also apply.
\item
A specialization of \exposid{per},
then its input is the product of the inverse of its template arguments, and
the following also apply.
\item
A specialization of the form \tcode{\exposidnc{power}<F, Num>},
then its input is $\tcode{F}^\tcode{Num}$, or
a specialization of the form \tcode{\exposidnc{power}<F, Num, Den>},
then its input is $\tcode{F}^{\tcode{Num}/\tcode{Den}}$, and
the following also applies.
\item
Otherwise, the input is the symbolic constant itself.
\end{itemize}
\begin{example}
Item by item, this algorithm step goes from the \Cpp{} parameter type
\tcode{decltype(km / square(h))},
styled in diagnostics like
\tcode{derived_unit<si::kilo_<si::metre>, \exposidnc{per}<\exposidnc{power}<non_si::hour, 2>>},
\begin{itemize}
\item
to $\tcode{decltype(km)} \times \tcode{\exposidnc{per}<\exposidnc{power}<decltype(h), 2>}$ (product of \tcode{To}'s arguments),
\item
to $\tcode{decltype(km)} \times 1/\tcode{\exposidnc{power}<decltype(h), 2>}$ (product of inverse of \exposid{per}'s arguments),
\item
to $\tcode{decltype(km)} \times 1/\tcode{decltype(h)}^2$ (\exposid{power}s as powers),
\item
to $\txtrm{a} \times 1/\txtrm{b}^2$ where $\txtrm{a} = \tcode{decltype(km)}$ and $\txtrm{b} = \tcode{decltype(h)}$ (symbolic substitution)
in the mathematical domain.
\end{itemize}
\end{example}

\pnum
Then, the operation takes place:
\begin{itemize}
\item
\exposid{expr-multiply} multiplies its inputs,
\item
\exposid{expr-divide} divides the input of its first parameter by the input of its second parameter,
\item
\exposid{expr-invert} divides $1$ by its input, and
\item
\exposid{expr-pow} raises its input to the $\tcode{Num}/\tcode{Den}$.
\end{itemize}

\pnum
Finally, let $r$ be the result of the operation simplified as follows:
\begin{itemize}
\item
All terms are part of the same fraction (if any).
\item
There is at most a single term with a given symbolic constant.
\item
There are no negative exponents.
\item
$1$ is only present as $r$ and as a numerator with a denominator not equal to $1$.
\end{itemize}
\begin{example}
Item by item: \\
$x \times 1/y \times 1/x^2$ \\
$= x/(y x^2)$ (single fraction) \\
$= x^{-1}/y$ (unique symbolic constants) \\
$= 1/(x^1 y)$ (positive exponents) \\
$= 1/(xy)$ (non-redundant $1$s) \\
\end{example}

\pnum
\returns
$r$ is mapped to the return type:
\begin{itemize}
\item
If $r = 1$, returns \tcode{OneType\{\}}.
\item
Otherwise, if $r$ is a symbolic constant, returns $r$.
\item
Otherwise, first applies the following mappings to the terms of $r$:
\begin{itemize}
\item
$x^{n/d}$ is mapped to \tcode{\exposidnc{power}<$x$, $n$, $d$>}, and
$x^{n}$ is mapped to \tcode{\exposidnc{power}<$x$, $n$>}, and
\item
$1$ is mapped to \tcode{OneType\{\}}.
\end{itemize}
\item
Then, a denominator $x$ of $r$ (if any) is mapped to \tcode{\exposidnc{per}<$x$>}.
\item
Then, sorts $r$ without \exposid{per} (if any) and
the template arguments of \exposid{per} (if any)
according to \tcode{Pred}.
\item
Finally, returns \tcode{To<$r$>\{\}}, where \exposid{per} (if any) is the last argument.
\end{itemize}

\pnum
\remarks
A valid template argument list for \tcode{To} and \exposid{per}
is formed by interspersing commas between each mapped term.
If a mapping to \tcode{std::intmax_t} is not representable,
the program is ill-formed.
\end{itemdescr}

\pnum
\exposid{expr-map} maps the contents of one symbolic expression to another resulting in a different type list.

\indexlibraryglobalexpos{expr-map}
\begin{itemdecl}
template<template<typename> typename Proj, template<typename...> typename To, typename OneType,
         template<typename, typename> typename Pred = @\exposidnc{type-less}@, typename T>
consteval auto @\exposidnc{expr-map}@(T);  // \expos
\end{itemdecl}

\begin{itemdescr}
\pnum
Let
\begin{itemize}
\item
\tcode{\exposidnc{expr-type-map}<U>} be
\tcode{\exposidnc{power}<Proj<F>, Ints...>} if \tcode{U} is of the form \tcode{\exposidnc{power}<F, Ints...>}, and
\tcode{Proj<U>} otherwise,
\item
\tcode{\exposidnc{map-power}(u)} be
\tcode{pow<Ints...>(F\{\})} if \tcode{decltype(u)} is of the form \tcode{\exposidnc{power}<F, Ints...>}, and
\tcode{u} otherwise, and
\item
\tcode{Nums} and \tcode{Dens}
be packs denoting the template arguments of
\tcode{T::\exposidnc{nums}} and \tcode{T::\exposidnc{dens}}, respectively.
\end{itemize}

\pnum
\returns
\begin{codeblock}
(OneType{} * ... * @\exposidnc{map-power}@(@\exposidnc{expr-type-map}@<Nums>{})) /
(OneType{} * ... * @\exposidnc{map-power}@(@\exposidnc{expr-type-map}@<Dens>{}))
\end{codeblock}
\end{itemdescr}

\rSec1[qty.ref]{Reference}

\rSec2[qty.ref.general]{General}

\pnum
Subclause \ref{qty.ref} specifies the components
for describing the reference of a quantity\irefiev{112-01-01}.

\rSec2[qty.dim]{Dimension}

\rSec3[qty.dim.general]{General}

\pnum
Subclause \ref{qty.dim} specifies the components
for defining the dimension of a quantity\irefiev{112-01-11}.

\rSec3[qty.dim.concepts]{Concepts}

\begin{itemdecl}
template<typename T>
concept @\deflibconcept{dimension}@ = @\exposconceptnc{SymbolicConstant}@<T> && std::@\stdconcept{derived_from}@<T, @\exposidnc{dimension-interface}@>;

template<typename T>
concept @\defexposconceptnc{BaseDimension}@ =  // \expos
  @\libconcept{dimension}@<T> && (@\exposidnc{is-derived-from-specialization-of}@<T, base_dimension>());
\end{itemdecl}

\rSec3[qty.dim.types]{Types}

\begin{codeblock}
namespace std {

template<symbol_text Symbol>
struct @\libglobal{base_dimension}@ : @\exposidnc{dimension-interface}@ {
  static constexpr auto @\exposidnc{symbol}@ = Symbol;  // \expos
};

}
\end{codeblock}

\pnum
\tcode{base_dimension} is used
to define the dimension of a base quantity\irefiev{112-01-08}.
\tcode{Symbol} is its symbolic representation.
\begin{example}
\begin{codeblock}
inline constexpr struct dim_length final : base_dimension<"L"> {} dim_length;
\end{codeblock}
\end{example}

\indexlibraryglobal{derived_dimension}
\begin{codeblock}
namespace std {

template<typename... Expr>
struct @\exposidnc{derived-dimension-impl}@  // \expos
    : @\exposidnc{expr-fractions}@<struct dimension_one, Expr...> {};

template<@\exposconceptnc{SymbolicConstant}@... Expr>
struct derived_dimension final : @\exposidnc{dimension-interface}@, @\exposidnc{derived-dimension-impl}@<Expr...> {};

}
\end{codeblock}

\pnum
\tcode{derived_dimension} is used by the library
to represent the dimension of a derived quantity\irefiev{112-01-10}.
\begin{example}
\begin{codeblock}
constexpr auto dim_acceleration = isq::speed.dimension / isq::dim_time;
int x = dim_acceleration;  // error: cannot construct from
  // \tcode{derived_dimension<isq::dim_length, \exposidnc{per}<\exposidnc{power}<isq::dim_time, 2>>>}
\end{codeblock}
\end{example}
A program that instantiates a specialization of \tcode{derived_dimension}
that is not a possible result of the library specifications
is ill-formed, no diagnostic required.

\indexlibraryglobal{dimension_one}
\begin{codeblock}
namespace std {

struct dimension_one final : @\exposidnc{dimension-interface}@, @\exposidnc{derived-dimension-impl}@<> {};

}
\end{codeblock}

\pnum
\tcode{dimension_one} represents the dimension of a quantity of dimension one\irefiev{112-01-13}.

\rSec3[qty.dim.ops]{Operations}

\indexlibraryglobalexpos{dimension-interface}
\begin{codeblock}
namespace std {

struct @\exposidnc{dimension-interface}@ {  // \expos
  template<@\libconcept{dimension}@ Lhs, @\libconcept{dimension}@ Rhs>
  friend consteval @\libconcept{dimension}@ auto operator*(Lhs, Rhs);

  template<@\libconcept{dimension}@ Lhs, @\libconcept{dimension}@ Rhs>
  friend consteval @\libconcept{dimension}@ auto operator/(Lhs, Rhs);

  template<@\libconcept{dimension}@ Lhs, @\libconcept{dimension}@ Rhs>
  friend consteval bool operator==(Lhs, Rhs);
};

}
\end{codeblock}

\indexlibrarymember{operator*}{dimension}
\begin{itemdecl}
template<@\libconcept{dimension}@ Lhs, @\libconcept{dimension}@ Rhs>
friend consteval @\libconcept{dimension}@ auto operator*(Lhs, Rhs);
\end{itemdecl}

\begin{itemdescr}
\pnum
\returns
\tcode{\exposidnc{expr-multiply}<derived_dimension, struct dimension_one>(Lhs\{\}, Rhs\{\})}.
\end{itemdescr}

\indexlibrarymember{operator/}{dimension}
\begin{itemdecl}
template<@\libconcept{dimension}@ Lhs, @\libconcept{dimension}@ Rhs>
friend consteval @\libconcept{dimension}@ auto operator/(Lhs, Rhs);
\end{itemdecl}

\begin{itemdescr}
\pnum
\returns
\tcode{\exposidnc{expr-divide}<derived_dimension, struct dimension_one>(Lhs\{\}, Rhs\{\})}.
\end{itemdescr}

\indexlibrarymember{operator==}{dimension}
\begin{itemdecl}
template<@\libconcept{dimension}@ Lhs, @\libconcept{dimension}@ Rhs>
friend consteval bool operator==(Lhs, Rhs);
\end{itemdecl}

\begin{itemdescr}
\pnum
\returns
\tcode{std::is_same_v<Lhs, Rhs>}.
\end{itemdescr}

\indexlibrarymember{inverse}{dimension}
\begin{itemdecl}
consteval @\libconcept{dimension}@ auto inverse(@\libconcept{dimension}@ auto d);
\end{itemdecl}

\begin{itemdescr}
\pnum
\returns
\tcode{dimension_one / d}.
\end{itemdescr}

\indexlibrarymember{pow}{dimension}
\begin{itemdecl}
template<std::intmax_t Num, std::intmax_t Den = 1, @\libconcept{dimension}@ D>
  requires(Den != 0)
consteval @\libconcept{dimension}@ auto pow(D d);
\end{itemdecl}

\begin{itemdescr}
\pnum
\returns
\tcode{\exposidnc{expr-pow}<Num, Den, derived_dimension, struct dimension_one>(d)}.
\end{itemdescr}

\indexlibrarymember{sqrt}{dimension}
\begin{itemdecl}
consteval @\libconcept{dimension}@ auto sqrt(@\libconcept{dimension}@ auto d);
\end{itemdecl}

\begin{itemdescr}
\pnum
\returns
\tcode{pow<1, 2>(d)}.
\end{itemdescr}

\indexlibrarymember{cbrt}{dimension}
\begin{itemdecl}
consteval @\libconcept{dimension}@ auto cbrt(@\libconcept{dimension}@ auto d);
\end{itemdecl}

\begin{itemdescr}
\pnum
\returns
\tcode{pow<1, 3>(d)}.
\end{itemdescr}

\rSec3[qty.dim.sym.fmt]{Symbol formatting}

\indexlibrarymember{dimension_symbol_to}{dimension}
\begin{itemdecl}
template<typename CharT = char, std::@\stdconcept{output_iterator}@<CharT> Out, @\libconcept{dimension}@ D>
constexpr Out dimension_symbol_to(Out out, D d, const dimension_symbol_formatting& fmt = {});
\end{itemdecl}

\begin{itemdescr}
\pnum
\effects
TBD.

\pnum
\returns
TBD.
\end{itemdescr}

\indexlibrarymember{dimension_symbol}{dimension}
\begin{itemdecl}
template<dimension_symbol_formatting fmt = {}, typename CharT = char, @\libconcept{dimension}@ D>
consteval std::string_view dimension_symbol(D);
\end{itemdecl}

\begin{itemdescr}
\pnum
\effects
Equivalent to:
\begin{codeblock}
TBD.
\end{codeblock}
\end{itemdescr}

\rSec2[qty.spec]{Quantity specification}

\rSec3[qty.spec.general]{General}

\pnum
Subclause \ref{qty.spec} specifies the components
for defining a quantity\irefiev{112-01-01}.

\rSec3[qty.spec.concepts]{Concepts}

\begin{itemdecl}
template<typename T>
concept @\deflibconcept{quantity_spec}@ = @\exposconceptnc{SymbolicConstant}@<T> && std::@\stdconcept{derived_from}@<T, @\exposidnc{quantity-spec-interface}@>;

template<typename T>
concept @\defexposconceptnc{QuantityKindSpec}@ =                                   // \expos
  @\libconcept{quantity_spec}@<T> && @\exposidnc{is-specialization-of}@<T, kind_of_>();

template<typename T>
concept @\defexposconceptnc{PointQuantitySpec}@ =                                  // \expos
  @\libconcept{quantity_spec}@<T> && @\exposidnc{is-specialization-of}@<T, point_>();

template<typename T>
concept @\defexposconceptnc{DeltaQuantitySpec}@ =                                  // \expos
  @\libconcept{quantity_spec}@<T> && @\exposidnc{is-specialization-of}@<T, @\exposidnc{delta-spec-t}@>();

template<typename T>
concept @\defexposconceptnc{NamedQuantitySpec}@ =                                  // \expos
  @\libconcept{quantity_spec}@<T> && @\exposidnc{is-derived-from-specialization-of}@<T, named_quantity_spec>() &&
  (!@\exposconceptnc{QuantityKindSpec}@<T>);

template<typename T>
concept @\defexposconceptnc{DerivedQuantitySpec}@ =                                // \expos
  @\libconcept{quantity_spec}@<T> &&
  (@\exposidnc{is-specialization-of}@<T, derived_quantity_spec>() ||
   (@\exposconceptnc{QuantityKindSpec}@<T> &&
    @\exposidnc{is-specialization-of}@<decltype(auto(T::@\exposidnc{quantity-spec}@)), derived_quantity_spec>()));

template<auto Child, auto Parent>
concept @\defexposconceptnc{ChildQuantitySpecOf}@ = (@\exposidnc{is-child-of}@(Child, Parent));  // \expos

template<auto To, auto From>
concept @\defexposconceptnc{NestedQuantityKindSpecOf}@ =                           // \expos
  @\libconcept{quantity_spec}@<decltype(From)> && @\libconcept{quantity_spec}@<decltype(To)> &&
  (get_kind(From) != get_kind(To)) && @\exposconceptnc{ChildQuantitySpecOf}@<To, get_kind(From).@\exposidnc{quantity-spec}@>;

template<auto From, auto To>
concept @\defexposconceptnc{QuantitySpecConvertibleTo}@ =                          // \expos
  @\libconcept{quantity_spec}@<decltype(From)> && @\libconcept{quantity_spec}@<decltype(To)> && @\exposid{implicitly-convertible}@(From, To);

template<auto From, auto To>
concept @\defexposconceptnc{QuantitySpecExplicitlyConvertibleTo}@ =                // \expos
  @\libconcept{quantity_spec}@<decltype(From)> && @\libconcept{quantity_spec}@<decltype(To)> && @\exposid{explicitly-convertible}@(From, To);

template<auto From, auto To>
concept @\defexposconceptnc{QuantitySpecCastableTo}@ =                             // \expos
  @\libconcept{quantity_spec}@<decltype(From)> && @\libconcept{quantity_spec}@<decltype(To)> && @\exposid{castable}@(From, To);

template<typename T, auto QS>
concept @\deflibconcept{quantity_spec_of}@ =
  @\libconcept{quantity_spec}@<T> && @\libconcept{quantity_spec}@<decltype(QS)> && @\exposconceptnc{QuantitySpecConvertibleTo}@<T{}, QS> &&
  !@\exposconceptnc{NestedQuantityKindSpecOf}@<T{}, QS> &&
  (@\exposconceptnc{QuantityKindSpec}@<T> || !@\exposconceptnc{NestedQuantityKindSpecOf}@<QS, T{}>);

template<typename T>
concept @\defexposconceptnc{QSProperty}@ = (!@\libconcept{quantity_spec}@<T>);                     // \expos

consteval bool @\exposidnc{admits-character}@(@\libconcept{quantity_spec}@ auto q, quantity_character ch);      // \expos
consteval bool @\exposidnc{characters-overlap}@(@\libconcept{quantity_spec}@ auto q1, @\libconcept{quantity_spec}@ auto q2);  // \expos

template<typename T>
concept @\defexposconceptnc{ScalarOperand}@=                                      // \expos
  @\libconcept{quantity_spec}@<T> && (@\exposid{admits-character}@(T{}, quantity_character::scalar) ||
                        @\exposid{admits-character}@(T{}, quantity_character::complex));

template<typename T>
concept @\defexposconceptnc{NonScalarOperand}@ =                                   // \expos
  @\libconcept{quantity_spec}@<T> && (@\exposid{admits-character}@(T{}, quantity_character::vector) ||
                        @\exposid{admits-character}@(T{}, quantity_character::tensor));
\end{itemdecl}

\pnum
The character of a quantity is a property of its quantity specification.
For a quantity specification \tcode{q},
let \tcode{Q} be \tcode{\exposid{remove-point}(\exposid{remove-delta}(q))}\iref{qty.spec.utils}.
\tcode{q} \defnx{admits}{character!admitted} every value of \tcode{quantity_character}
if \tcode{decltype(Q)} models \exposconceptnc{QuantityKindSpec}
or is a specialization of \tcode{derived_quantity_spec} with a kind factor\iref{derived.qty},
and the value \tcode{Q.character} only otherwise.
\tcode{\exposid{admits-character}(q, ch)} is \tcode{true} if \tcode{q} admits \tcode{ch},
and \tcode{false} otherwise.
\tcode{\exposid{characters-overlap}(q1, q2)} is \tcode{true} if
some value of \tcode{quantity_character} is admitted by both \tcode{q1} and \tcode{q2},
and \tcode{false} otherwise.
The values \tcode{scalar} and \tcode{complex} are of the tensor order of a scalar,
\tcode{vector} is of the tensor order of a vector, and \tcode{tensor} of the tensor order of a tensor.
\tcode{\exposconceptnc{ScalarOperand}<T>} is satisfied if \tcode{T\{\}} admits a value of the tensor order of a scalar,
and \tcode{\exposconceptnc{NonScalarOperand}<T>} if it admits a value of the tensor order of a vector or a tensor.
\begin{note}
A kind admits every character,
because only the representation type of a quantity of a kind says which character it has\iref{qty.rep.concepts}.
A kind satisfies both \exposconceptnc{ScalarOperand} and \exposconceptnc{NonScalarOperand}.
\end{note}

\rSec3[qty.spec.types]{Types}

\rSec4[named.qty]{Named}

\indexlibraryglobal{subkind}
\indexlibraryglobal{named_quantity_spec}
\begin{codeblock}
namespace std {

struct subkind {};

template<@\exposconceptnc{BaseDimension}@ auto Dim, @\exposconceptnc{QSProperty}@ auto... Args>
struct named_quantity_spec<Dim, Args...> : @\exposidnc{quantity-spec-interface}@ {
  using @\exposidnc{base-type}@ = named_quantity_spec;
  static constexpr @\exposconceptnc{BaseDimension}@ auto dimension = Dim;
  static constexpr quantity_character character = @\seebelownc@;
};

template<@\exposconceptnc{DerivedQuantitySpec}@ auto Eq, @\exposconceptnc{QSProperty}@ auto... Args>
struct named_quantity_spec<Eq, Args...> : @\exposidnc{quantity-spec-interface}@ {
  using @\exposidnc{base-type}@ = named_quantity_spec;
  static constexpr auto @\exposidnc{equation}@ = Eq;
  static constexpr @\libconcept{dimension}@ auto dimension = Eq.dimension;
  static constexpr quantity_character character = @\seebelownc@;
};

template<@\exposconceptnc{NamedQuantitySpec}@ auto QS, @\exposconceptnc{QSProperty}@ auto... Args>
struct named_quantity_spec<QS, Args...> : @\exposidnc{quantity-spec-interface}@ {
  using @\exposidnc{base-type}@ = named_quantity_spec;
  static constexpr auto @\exposidnc{parent}@ = QS;
  static constexpr auto @\exposidnc{equation}@ = @\exposidnc{parent}@.@\exposidnc{equation}@;  // \expos, present only
    // if the \fakegrammarterm{qualified-id} \tcode{\exposidnc{parent}.\exposidnc{equation}} is valid and denotes an object
  static constexpr @\libconcept{dimension}@ auto dimension = @\exposidnc{parent}@.dimension;
  static constexpr quantity_character character = @\seebelownc@;
};

template<@\exposconceptnc{NamedQuantitySpec}@ auto QS, @\exposconceptnc{DerivedQuantitySpec}@ auto Eq, @\exposconceptnc{QSProperty}@ auto... Args>
  requires @\exposconceptnc{QuantitySpecExplicitlyConvertibleTo}@<Eq, QS>
struct named_quantity_spec<QS, Eq, Args...> : @\exposidnc{quantity-spec-interface}@ {
  using @\exposidnc{base-type}@ = named_quantity_spec;
  static constexpr auto @\exposidnc{parent}@ = QS;
  static constexpr auto @\exposidnc{equation}@ = Eq;
  static constexpr @\libconcept{dimension}@ auto dimension = @\exposidnc{parent}@.dimension;
  static constexpr quantity_character character = @\seebelownc@;
};

}
\end{codeblock}

\pnum
A \defnadj{named}{quantity} is a type that models \exposconceptnc{NamedQuantitySpec}.
A specialization of \tcode{named_quantity_spec} is used as a base type when defining a named quantity.

\pnum
In the following descriptions, let \tcode{Q} be a named quantity defined with an alluded signature.
The identifier of \tcode{Q} represents its quantity name\irefiev{112-01-02}.

\pnum
Let \tcode{Ch} be an enumerator value of \tcode{quantity_character}.
The possible arguments to \tcode{named_quantity_spec} are
\begin{itemize}
\item
$(\text{a base quantity dimension}, \opt{\tcode{Ch}})$,
\item
$(\text{a quantity calculus}, \opt{\tcode{Ch}})$,
\item
$(\text{a named quantity}, \opt{\tcode{Ch}}, \opt{\tcode{subkind}})$, and
\item
$(\text{a named quantity}, \text{a quantity calculus}, \opt{\tcode{Ch}}, \opt{\tcode{subkind}})$.
\end{itemize}

\pnum
If the first argument is a base quantity dimension,
then \tcode{Q} is that base quantity\irefiev{112-01-08}.
If an argument is a quantity calculus\irefiev{112-01-30} \placeholder{C},
then \tcode{Q} is implicitly convertible from \placeholder{C}.
If the first argument is a named quantity,
then \tcode{Q} is of its kind\irefiev{112-01-04}.

\pnum
The member \tcode{character} represents
the set of the numerical value of \tcode{Q}\iref{qty.char.traits}
and is equal to
\begin{itemize}
\item
\tcode{Ch} if specified,
\item
otherwise, \tcode{quantity_character::real_scalar} for the first signature, and
\item
otherwise, \tcode{(BC).character},
where \tcode{BC} is the argument preceding \tcode{Ch} in the signatures above.
\end{itemize}

\pnum
\tcode{subkind} specifies \tcode{Q} to start a new hierarchy tree of a kind.

\pnum
Optional arguments may appear in any order.

\pnum
\begin{example}
\begin{codeblock}
// The first signature defines a base quantity.
inline constexpr struct length final : named_quantity_spec<dim_length> {
} length;  // Length is a base quantity.

// The second signature defines a derived quantity.
inline constexpr struct area final : named_quantity_spec<pow<2>(length)> {
} area;  // An area equals length by length.

// The third and fourth signatures add a leaf to a hierarchy of kinds.
inline constexpr struct width final : named_quantity_spec<length> {
} width;  // Width is a kind of length.

// The fourth signature also refines the calculus required for implicit conversions.
inline constexpr struct angular_measure final :
    named_quantity_spec<dimensionless, arc_length / radius, subkind> {
} angular_measure;  // Requires an arc length per radius, not just any quantity of dimension one.
\end{codeblock}
\end{example}

\rSec4[derived.qty]{Derived}

\indexlibraryglobal{derived_quantity_spec}
\begin{codeblock}
namespace std {

template<@\libconcept{quantity_spec}@ Q>
using @\exposidnc{to-dimension}@ = decltype(auto(Q::dimension));  // \expos

template<typename... Expr>
struct @\exposidnc{derived-quantity-spec-impl}@ :                 // \expos
    @\exposidnc{quantity-spec-interface}@,
    @\exposidnc{expr-fractions}@<struct dimensionless, Expr...> {
  using @\exposidnc{base-type}@ = @\exposidnc{derived-quantity-spec-impl}@;
  using @\exposidnc{base}@ = @\exposidnc{expr-fractions}@<struct dimensionless, Expr...>;

  static constexpr @\libconcept{dimension}@ auto dimension =
    @\exposidnc{expr-map}@<@\exposidnc{to-dimension}@, derived_dimension, struct dimension_one>(@\exposidnc{base}@{});
  static constexpr quantity_character character = @\seebelownc@;
};

template<@\exposconceptnc{SymbolicConstant}@... Expr>
struct derived_quantity_spec final : @\exposidnc{derived-quantity-spec-impl}@<Expr...> {};

}
\end{codeblock}

\pnum
\tcode{derived_quantity_spec} is used by the library
to represent the result of a quantity calculus not equal to a named quantity.
\begin{example}
\begin{codeblock}
constexpr auto area = pow<2>(isq::length);
int x = area;  // error: cannot construct from \tcode{derived_quantity_spec<\exposidnc{power}<isq::length, 2>>}
\end{codeblock}
\end{example}
A program that instantiates a specialization of \tcode{derived_quantity_spec}
that is not a possible result of the library specifications
is ill-formed, no diagnostic required.

\pnum
An element of \tcode{Expr} that is a specialization of \tcode{kind_of_},
or a specialization of \exposid{power} whose first template argument is one,
is a \defnx{kind factor}{kind factor}.
\begin{note}
A kind factor results from an operation on a kind and a quantity specification that is not a kind
\iref{qty.spec.ops}, and denotes a factor that can be any quantity of the tree below its kind root:
\tcode{isq::height / kind_of<isq::duration>} is
\tcode{derived_quantity_spec<isq::height, \exposidnc{per}<kind_of_<isq::duration>>>}.
A kind factor and a named quantity are different elements even when the kind factor wraps that quantity,
so \tcode{kind_of<isq::length> / isq::length} is not \tcode{dimensionless},
as \tcode{isq::width / isq::length} is not.
\end{note}

\pnum
Let
\begin{itemize}
\item
\tcode{Nums} and \tcode{Dens}
be packs denoting the template arguments of
\tcode{\exposidnc{base}::\exposidnc{nums}} and \tcode{\exposidnc{base}::\exposidnc{dens}}, respectively,
\item
\tcode{\placeholdernc{QUANTITY-CHARACTER-OF}(Pack)} be
\begin{codeblock}
std::max({quantity_character::real_scalar, @\exposidnc{expr-type}@<Pack>::character...})
\end{codeblock}
and
\item
\tcode{num_char} be \tcode{\placeholdernc{QUANTITY-CHARACTER-OF}(Nums)} and
\tcode{den_char} be \tcode{\placeholdernc{QUANTITY-CHARACTER-OF}(Dens)}.
\end{itemize}
The member \tcode{character} is equal to
\tcode{quantity_character::real_scalar} if \tcode{num_char == den_char} is \tcode{true}, and
\tcode{std::max(num_char, den_char)} otherwise.

\rSec4[dimless.qty]{Base quantity of dimension one}

\indexlibraryglobal{dimensionless}
\begin{codeblock}
namespace std {

struct dimensionless final : named_quantity_spec<derived_quantity_spec<>{}> {};

}
\end{codeblock}

\pnum
\tcode{dimensionless} represents the base quantity of dimension one\irefiev{112-01-13}.

\rSec4[kind.of.qty]{Kind of}

\indexlibraryglobal{kind_of_}
\begin{codeblock}
namespace std {

template<@\libconcept{quantity_spec}@ Q>
  requires(!@\exposconceptnc{QuantityKindSpec}@<Q>) && (!@\exposconceptnc{PointQuantitySpec}@<Q>) &&
           (!@\exposconceptnc{DeltaQuantitySpec}@<Q>) && (@\exposidnc{get-kind-tree-root}@(Q{}) == Q{})
struct kind_of_ final : Q::@\exposidnc{base-type}@ {
  using @\exposidnc{base-type}@ = kind_of_;                 // \expos
  static constexpr auto @\exposidnc{quantity-spec}@ = Q{};  // \expos
};

}
\end{codeblock}

\pnum
\tcode{kind_of<Q>} represents a kind of quantity\irefiev{112-01-04} \tcode{Q}.
\begin{note}
The quantity specification of a point is never a kind,
so \tcode{kind_of<point<isq::length>>} is ill-formed,
and so is \tcode{kind_of<delta<isq::length>>}\iref{delta.qty}.
\tcode{get_quantity_spec(point<si::metre>)} is \tcode{point<kind_of<isq::length>>}\iref{point.qty}.
\end{note}

\rSec4[delta.qty]{Delta}

\indexlibraryglobalexpos{delta-spec-t}
\begin{codeblock}
namespace std {

template<@\libconcept{quantity_spec}@ Q>
  requires(!@\exposconceptnc{PointQuantitySpec}@<Q>) && (!@\exposconceptnc{DeltaQuantitySpec}@<Q>)
struct @\exposidnc{delta-spec-t}@ : @\exposidnc{quantity-spec-interface}@ {  // \expos
  using @\exposidnc{base-type}@ = @\exposidnc{delta-spec-t}@;                    // \expos
  static constexpr auto @\exposidnc{quantity-spec}@ = Q{};           // \expos
  static constexpr @\libconcept{dimension}@ auto dimension = Q::dimension;
  static constexpr quantity_character character = Q::character;
};

}
\end{codeblock}

\pnum
\tcode{\exposid{delta-spec-t}<Q>} represents the quantity specification of a delta\iref{qty.general}
of the quantity \tcode{Q}, a \defnadj{delta}{quantity specification}.
For a quantity specification \tcode{QS} that models neither \exposconceptnc{PointQuantitySpec}
nor \exposconceptnc{DeltaQuantitySpec},
\tcode{delta<QS>}\iref{qty.delta} is \tcode{\exposid{delta-spec-t}<decltype(QS)>\{\}}.
\begin{note}
\tcode{delta<delta<isq::height>>}, \tcode{delta<point<isq::height>>}, and
\tcode{point<delta<isq::height>>} are ill-formed.
\tcode{QS} can be a kind: \tcode{delta<kind_of<isq::length>>} is well-formed.
\end{note}

\pnum
\tcode{delta<QS>} differs from \tcode{QS} only in the units its subscript operator accepts:
for a unit \tcode{U} that models \exposconceptnc{OffsetUnit}\iref{qty.unit.concepts},
\tcode{delta<QS>[U]} is the reference of a delta, while \tcode{QS[U]} is ill-formed\iref{qty.spec.ops}.
Except where stated otherwise,
an entity specified in \ref{quantities} that is given a delta quantity specification \tcode{q}
as a function argument or as a template argument
behaves as if it were given \tcode{\exposid{remove-delta}(q)}\iref{qty.spec.utils}.
The exceptions are
the subscript operator and the call operators of a quantity specification\iref{qty.spec.ops},
\tcode{operator==} on quantity specifications\iref{qty.spec.ops},
\tcode{kind_of_}\iref{kind.of.qty}, \tcode{\exposid{delta-spec-t}}, \tcode{point_}\iref{point.qty},
and the variable templates \tcode{delta}\iref{qty.delta} and \tcode{point}\iref{qty.point}.
\begin{note}
The quantity specification of the reference of a delta is never a delta quantity specification:
\tcode{get_quantity_spec(delta<isq::width>[m])} is \tcode{isq::width}.
\end{note}
\begin{example}
\begin{codeblock}
static_assert(delta<isq::height> != isq::height);
static_assert(delta<isq::height> * isq::duration == isq::height * isq::duration);
static_assert(get_kind(delta<isq::height>) == kind_of<isq::length>);
static_assert(@\exposid{implicitly-convertible}@(delta<isq::height>, isq::length));
static_assert(delta<isq::height>[m] == isq::height[m]);
static_assert(get_quantity_spec(delta<isq::thermodynamic_temperature>[deg_C]) ==
              isq::thermodynamic_temperature);
\end{codeblock}
\end{example}

\rSec4[point.qty]{Point}

\indexlibraryglobal{point_}
\begin{codeblock}
namespace std {

template<@\libconcept{quantity_spec}@ Q>
  requires(!@\exposconceptnc{PointQuantitySpec}@<Q>) && (!@\exposconceptnc{DeltaQuantitySpec}@<Q>)
struct point_<Q> : @\exposidnc{quantity-spec-interface}@ {
  using @\exposidnc{base-type}@ = point_;                   // \expos
  static constexpr auto @\exposidnc{quantity-spec}@ = Q{};  // \expos
  static constexpr @\libconcept{dimension}@ auto dimension = Q::dimension;
  static constexpr quantity_character character = Q::character;
};

template<@\libconcept{quantity_spec}@ Q, @\exposconceptnc{PointOrigin}@ PO>
  requires(!@\exposconceptnc{PointQuantitySpec}@<Q>) && (!@\exposconceptnc{DeltaQuantitySpec}@<Q>) &&
          @\libconcept{point_origin_for}@<PO, Q{}>
struct point_<Q, PO> : @\exposidnc{quantity-spec-interface}@ {
  using @\exposidnc{base-type}@ = point_;                   // \expos
  static constexpr auto @\exposidnc{quantity-spec}@ = Q{};  // \expos
  static constexpr PO @\exposidnc{point-origin}@{};         // \expos
  static constexpr @\libconcept{dimension}@ auto dimension = Q::dimension;
  static constexpr quantity_character character = Q::character;
};

}
\end{codeblock}

\pnum
\tcode{point_<Q>} represents the quantity specification of a point,
a quantity measured from a point origin\iref{qty.pt.orig} on the scale of the quantity \tcode{Q},
measured from \tcode{natural_point_origin<Q\{\}>}\iref{qty.natural.pt.orig}.
\tcode{point_<Q, PO>} represents the quantity specification of a point on the scale of \tcode{Q}
measured from the point origin \tcode{PO\{\}}.
\tcode{Q} is the \defnx{scale}{scale!of a point quantity specification} of either,
and the point origin it is measured from is its \defnx{origin}{origin!of a point quantity specification}.
For a quantity specification \tcode{QS} that models neither \exposconceptnc{PointQuantitySpec}
nor \exposconceptnc{DeltaQuantitySpec} and a point origin \tcode{po},
\tcode{point<QS>}\iref{qty.point} is \tcode{point_<decltype(QS)>\{\}} and
\tcode{point<QS, po>} is \tcode{\exposid{normalized-point-spec}(QS, po)}.
\begin{note}
\tcode{point<isq::width, mean_sea_level>}
for an absolute point origin \tcode{mean_sea_level} of \tcode{isq::altitude} is ill-formed.
\end{note}

\indexlibrarymemberexpos{quantity_spec}{point-spec-origin}
\begin{itemdecl}
template<@\exposconceptnc{PointQuantitySpec}@ Q>
consteval @\exposconceptnc{PointOrigin}@ auto @\exposidnc{point-spec-origin}@(Q qs);  // \expos
\end{itemdecl}

\begin{itemdescr}
\pnum
\returns
\tcode{Q::\exposid{point-origin}} if \tcode{Q} is a specialization \tcode{point_<S, PO>},
and \tcode{natural_point_origin<Q::\exposid{quantity-spec}>} otherwise.
\end{itemdescr}

\indexlibrarymemberexpos{quantity_spec}{normalized-point-spec}
\begin{itemdecl}
template<@\libconcept{quantity_spec}@ QS, @\exposconceptnc{PointOrigin}@ PO>
  requires(!@\exposconceptnc{PointQuantitySpec}@<QS>) && @\libconcept{point_origin_for}@<PO, QS{}>
consteval @\exposconceptnc{PointQuantitySpec}@ auto @\exposidnc{normalized-point-spec}@(QS qs, PO po);  // \expos
\end{itemdecl}

\begin{itemdescr}
\pnum
Let \tcode{S} be \tcode{\exposid{remove-delta}(qs)}.

\pnum
\returns
\tcode{point_<decltype(S)>\{\}} if \tcode{po == natural_point_origin<S>} is \tcode{true}\iref{qty.pt.orig.ops},
and \tcode{point_<decltype(S), PO>\{\}} otherwise.
\begin{note}
A point quantity specification has one spelling:
\tcode{point<isq::length, natural_point_origin<isq::length>>} is \tcode{point<isq::length>}.
\end{note}
\end{itemdescr}

\pnum
\begin{note}
A point is a quantity whose reference models \exposconceptnc{PointReference}\iref{qty.ref.concepts}.
The quantity specification of a point names its origin\iref{qty.point},
so a point and a delta never share a quantity specification.
A point never converts to a delta, nor a delta to a point.
Two points convert only when their origins have the same absolute point origin,
and then as the quantity specifications of their scales do,
along one branch of the hierarchy only\iref{qty.spec.conv}.
The quantity specification of a point cannot be multiplied, divided, or raised to a power,
it has no magnitude\iref{qty.spec.ops},
and it is never a kind\iref{kind.of.qty}.
\end{note}
\begin{example}
\begin{codeblock}
inline constexpr struct mean_sea_level final : absolute_point_origin<isq::altitude> {} mean_sea_level;

static_assert(get_quantity_spec(point<isq::altitude>[m]) == point<isq::altitude>);
static_assert(point<isq::altitude, mean_sea_level> != point<isq::altitude>);
static_assert(point<isq::length, natural_point_origin<isq::length>> == point<isq::length>);
static_assert(get_kind(point<isq::altitude, mean_sea_level>) ==
              point<kind_of<isq::length>, mean_sea_level>);
static_assert(@\exposid{implicitly-convertible}@(point<isq::altitude>, point<isq::length>));
static_assert(!@\exposid{implicitly-convertible}@(point<isq::altitude>, isq::length));
static_assert(!@\exposid{castable}@(point<isq::height>, point<isq::width>));
static_assert(!@\exposid{castable}@(point<isq::altitude, mean_sea_level>, point<isq::altitude>));
\end{codeblock}
\end{example}

\rSec4[magnitude.qty]{Magnitude}

\indexlibraryglobalexpos{magnitude-t}
\begin{codeblock}
namespace std {

template<@\libconcept{quantity_spec}@ V>
  requires(!@\exposconceptnc{QuantityKindSpec}@<V>) && (!@\exposconceptnc{PointQuantitySpec}@<V>) &&
          (!@\exposconceptnc{DeltaQuantitySpec}@<V>) && (!@\exposidnc{is-specialization-of}@<V, @\exposidnc{magnitude-t}@>()) &&
          @\exposconceptnc{NonScalarOperand}@<V>
struct @\exposidnc{magnitude-t}@ : @\exposidnc{quantity-spec-interface}@ {  // \expos
  using @\exposidnc{base-type}@ = @\exposidnc{magnitude-t}@;                    // \expos
  static constexpr auto @\exposidnc{quantity-spec}@ = V{};          // \expos
  static constexpr auto @\exposidnc{equation}@ = @\seebelownc@;            // \expos, present only if \tcode{V} has a magnitude equation
  static constexpr @\libconcept{dimension}@ auto dimension = V::dimension;
  static constexpr quantity_character character = quantity_character::scalar;
};

}
\end{codeblock}

\pnum
\tcode{\exposid{magnitude-t}<V>} is the type of \tcode{magnitude(v)}\iref{qty.spec.ops}
for an object \tcode{v} of type \tcode{V}.
It represents the quantity specification of the magnitude of a vector or tensor quantity of \tcode{V},
a scalar, which is a \defnx{magnitude}{magnitude!quantity specification}.
A magnitude can be a factor of a derived quantity specification\iref{derived.qty}
and so of the equation of a named quantity\iref{named.qty}.

\pnum
In \ref{qty.spec.conv} and \ref{get.common.qty.spec},
a magnitude \tcode{magnitude(v)} is a named quantity without a parent, which is not a declared kind,
whose \exposid{equation} is the magnitude equation of \tcode{V} (see below) if \tcode{V} has one,
and whose kind root is the kind root of \tcode{V}.
The paragraph of \ref{qty.spec.conv} that applies \tcode{\exposid{tree-rank}}
to two quantity specifications that both model \exposconceptnc{NamedQuantitySpec}
does not apply to a magnitude.

\pnum
The \defnx{magnitude equation}{magnitude equation} of \tcode{V} is determined as follows,
where a quantity specification \defnx{admits only a real scalar}{admits only a real scalar} if
it does not model \exposconceptnc{QuantityKindSpec} and its member \tcode{character} is \tcode{quantity_character::scalar}:
\begin{itemize}
\item
If \tcode{V} is a named quantity with a parent \tcode{P} that does not satisfy \exposconceptnc{ScalarOperand},
\tcode{magnitude(P)}.
\item
Otherwise, if \tcode{V} is a named quantity with a parent \tcode{P} that admits only a real scalar,
\tcode{P}.
\item
Otherwise, if \tcode{V} models \exposconceptnc{DerivedQuantitySpec} and
exactly one of its factors\iref{qty.spec.conv}, \tcode{X}, satisfies \exposconceptnc{NonScalarOperand},
and \tcode{X} is a numerator factor with exponent 1,
the derived quantity specification obtained from \tcode{V} by replacing \tcode{X} with \tcode{magnitude(X)}.
\item
Otherwise, if \tcode{V} is a named quantity whose \exposid{equation} \tcode{E}
has the same member \tcode{character} as \tcode{V} and is of the form of the previous item,
the magnitude equation of \tcode{E}.
\item
Otherwise, if \tcode{V} is a named quantity whose \exposid{equation} \tcode{E} admits only a real scalar,
\tcode{E}; such a magnitude equation is \defnx{relaxed}{magnitude equation!relaxed}.
\item
Otherwise, \tcode{V} has no magnitude equation.
\end{itemize}

\pnum
In \ref{qty.spec.conv}, expanding a factor whose named quantity is a magnitude \tcode{M}
ranks \tcode{explicit_only} if the magnitude equation of \tcode{M} is relaxed.
Otherwise, expanding such a factor of $T$ ranks \tcode{implicit}
if \tcode{from-kind} is \tcode{true}, or
if the magnitude equation of \tcode{M} is determined by the third or the fourth item above and
\tcode{V} models \exposconceptnc{DerivedQuantitySpec} or is its own defining ancestor,
and \tcode{explicit_only} otherwise.
\begin{note}
The first two items make the magnitude of a vector a special case of the magnitude of its parent,
and the magnitude of a vector defined below a scalar quantity a special case of that scalar.
So \tcode{magnitude(v)} converts implicitly to the nearest scalar quantity above \tcode{v},
and that quantity converts to \tcode{magnitude(v)} only explicitly,
as not every speed is the magnitude of a velocity.
A quantity of \tcode{v} itself never converts to \tcode{magnitude(v)},
as their characters differ\iref{qty.spec.conv}.
\end{note}
\begin{example}
\begin{codeblock}
inline constexpr struct pressure final : named_quantity_spec<magnitude(isq::force) / isq::area> {} pressure;

static_assert(magnitude(isq::velocity).dimension == isq::velocity.dimension);
static_assert(@\exposid{implicitly-convertible}@(magnitude(isq::velocity), isq::speed));
static_assert(!@\exposid{implicitly-convertible}@(isq::speed, magnitude(isq::velocity)));
static_assert(@\exposid{explicitly-convertible}@(isq::speed, magnitude(isq::velocity)));
static_assert(!@\exposid{castable}@(isq::velocity, isq::speed));
static_assert(!@\exposid{castable}@(isq::velocity, magnitude(isq::velocity)));
static_assert(@\exposid{implicitly-convertible}@(magnitude(isq::force) / isq::area, pressure));
static_assert(!@\exposid{castable}@(isq::force / isq::area, pressure));
\end{codeblock}
where \tcode{isq::velocity} is a vector quantity below the scalar \tcode{isq::speed}
and \tcode{isq::force} is a vector quantity.
\end{example}

\rSec3[qty.spec.utils]{Utilities}

\indexlibrarymemberexpos{quantity_spec}{open-kinds}
\begin{itemdecl}
template<@\libconcept{quantity_spec}@ Q>
consteval @\libconcept{quantity_spec}@ auto @\exposidnc{open-kinds}@(Q q);    // \expos
\end{itemdecl}

\begin{itemdescr}
\pnum
Let \exposid{to-kind-factor} be defined as follows:
\begin{codeblock}
template<@\exposconceptnc{NamedQuantitySpec}@ Q>
using @\exposidnc{to-kind-factor}@ = kind_of_<Q>;  // \expos
\end{codeblock}

\pnum
\effects
Equivalent to:
\begin{codeblock}
if constexpr (!@\exposconceptnc{QuantityKindSpec}@<Q>)
  return q;
else if constexpr (Q::@\exposidnc{quantity-spec}@ == dimensionless)
  return dimensionless;
else if constexpr (@\exposconceptnc{NamedQuantitySpec}@<decltype(Q::@\exposidnc{quantity-spec}@)>)
  return q;
else
  return @\exposidnc{expr-map}@<@\exposidnc{to-kind-factor}@, derived_quantity_spec, struct dimensionless>(Q::@\exposidnc{quantity-spec}@);
\end{codeblock}
\begin{note}
A kind of a product of named quantities becomes the product of the kinds of those quantities,
so that an operation on quantity specifications keeps a kind per factor.
\end{note}
\end{itemdescr}

\indexlibrarymemberexpos{quantity_spec}{close-kinds}
\begin{itemdecl}
template<@\libconcept{quantity_spec}@ auto... Operands, @\libconcept{quantity_spec}@ Q>
consteval @\libconcept{quantity_spec}@ auto @\exposidnc{close-kinds}@(Q q);    // \expos
\end{itemdecl}

\begin{itemdescr}
\pnum
Let \exposid{from-kind-factor} be defined as follows:
\begin{codeblock}
template<typename Q>
struct @\exposidnc{from-kind-factor-impl}@ : std::type_identity<Q> {};  // \expos

template<typename N>
struct @\exposidnc{from-kind-factor-impl}@<kind_of_<N>> : std::type_identity<N> {};

template<typename Q>
using @\exposidnc{from-kind-factor}@ = @\exposidnc{from-kind-factor-impl}@<Q>::type;  // \expos
\end{codeblock}

\pnum
\effects
Equivalent to:
\begin{codeblock}
if constexpr (@\exposconceptnc{QuantityKindSpec}@<Q>)
  return q;
else if constexpr (Q{} == dimensionless) {
  if constexpr ((... && @\exposconceptnc{QuantityKindSpec}@<decltype(Operands)>))
    return kind_of<Q{}>;
  else
    return q;
} else if constexpr (@\placeholder{ALL-KIND-FACTORS}@)
  return kind_of<@\exposidnc{expr-map}@<@\exposidnc{from-kind-factor}@, derived_quantity_spec, struct dimensionless>(q)>;
else
  return q;
\end{codeblock}
where \tcode{\placeholder{ALL-KIND-FACTORS}} is \tcode{true} if
\tcode{Q} models \exposconceptnc{DerivedQuantitySpec} and
every element of \tcode{Q::\exposidnc{num}} and of \tcode{Q::\exposidnc{den}}
is a specialization of \tcode{kind_of_} or
a specialization of \exposid{power} whose first template argument is one,
and \tcode{false} otherwise.
\begin{note}
A result whose factors are all kinds is the kind of the product of what they wrap,
while a result with at least one factor that is not a kind keeps its kind factors as they are.
\end{note}
\end{itemdescr}

\indexlibrarymemberexpos{quantity_spec}{remove-kind}
\begin{itemdecl}
template<@\libconcept{quantity_spec}@ Q>
consteval auto @\exposidnc{remove-kind}@(Q q);                   // \expos
\end{itemdecl}

\begin{itemdescr}
\pnum
\effects
Equivalent to:
\begin{codeblock}
if constexpr (@\exposconceptnc{QuantityKindSpec}@<Q>)
  return Q::@\exposidnc{quantity-spec}@;
else
  return q;
\end{codeblock}
\end{itemdescr}

\indexlibrarymemberexpos{quantity_spec}{remove-point}
\begin{itemdecl}
template<@\libconcept{quantity_spec}@ Q>
consteval @\libconcept{quantity_spec}@ auto @\exposidnc{remove-point}@(Q q);                  // \expos
\end{itemdecl}

\begin{itemdescr}
\pnum
\effects
Equivalent to:
\begin{codeblock}
if constexpr (@\exposconceptnc{PointQuantitySpec}@<Q>)
  return Q::@\exposidnc{quantity-spec}@;
else
  return q;
\end{codeblock}
\end{itemdescr}

\indexlibrarymemberexpos{quantity_spec}{remove-delta}
\begin{itemdecl}
template<@\libconcept{quantity_spec}@ Q>
consteval @\libconcept{quantity_spec}@ auto @\exposidnc{remove-delta}@(Q q);                  // \expos
\end{itemdecl}

\begin{itemdescr}
\pnum
\effects
Equivalent to:
\begin{codeblock}
if constexpr (@\exposconceptnc{DeltaQuantitySpec}@<Q>)
  return Q::@\exposidnc{quantity-spec}@;
else
  return q;
\end{codeblock}
\end{itemdescr}

\indexlibrarymemberexpos{quantity_spec}{as-point}
\begin{itemdecl}
template<@\libconcept{quantity_spec}@ Q>
consteval @\libconcept{quantity_spec}@ auto @\exposidnc{as-point}@(Q q);                      // \expos
\end{itemdecl}

\begin{itemdescr}
\pnum
\effects
Equivalent to:
\begin{codeblock}
if constexpr (@\exposconceptnc{PointQuantitySpec}@<Q>)
  return q;
else
  return point_<Q>{};
\end{codeblock}
\end{itemdescr}

\indexlibrarymemberexpos{quantity_spec}{make-reference}
\indexlibrarymemberexpos{unit}{make-reference}
\begin{itemdecl}
template<@\libconcept{quantity_spec}@ QS, @\libconcept{unit}@ U>
  requires @\libconcept{unit_of}@<U, QS{}>
consteval @\libconcept{quantity_reference}@ auto @\exposidnc{make-reference}@(QS, U u);  // \expos
\end{itemdecl}

\begin{itemdescr}
\pnum
\effects
Equivalent to:
\begin{codeblock}
if constexpr (@\exposconceptnc{PointQuantitySpec}@<QS>)
  return @\exposidnc{make-point-spec-reference}@(QS{}, u);
else if constexpr (@\exposconceptnc{DeltaQuantitySpec}@<QS>)
  return @\exposidnc{make-reference}@(@\exposidnc{remove-delta}@(QS{}), u);
else if constexpr (requires { requires get_quantity_spec(U{}) == QS{}; })
  return u;
else
  return reference<QS, U>{};
\end{codeblock}
\begin{note}
\exposid{make-point-spec-reference} is specified in \ref{qty.point}.
\end{note}
\end{itemdescr}

\rSec3[qty.spec.ops]{Operations}

\indexlibraryglobalexpos{quantity-spec-interface}
\begin{codeblock}
namespace std {

struct @\exposidnc{quantity-spec-interface}@ {  // \expos
  template<@\libconcept{quantity_spec}@ Lhs, @\libconcept{quantity_spec}@ Rhs>
    requires(!@\exposconceptnc{PointQuantitySpec}@<Lhs>) && (!@\exposconceptnc{PointQuantitySpec}@<Rhs>) &&
            (@\exposconceptnc{ScalarOperand}@<Lhs> || @\exposconceptnc{ScalarOperand}@<Rhs>)
  friend consteval @\libconcept{quantity_spec}@ auto operator*(Lhs lhs, Rhs rhs);

  template<@\libconcept{quantity_spec}@ Lhs, @\libconcept{quantity_spec}@ Rhs>
    requires(!@\exposconceptnc{PointQuantitySpec}@<Lhs>) && (!@\exposconceptnc{PointQuantitySpec}@<Rhs>) &&
            @\exposconceptnc{ScalarOperand}@<Rhs>
  friend consteval @\libconcept{quantity_spec}@ auto operator/(Lhs lhs, Rhs rhs);

  template<typename Self, @\libconcept{unit_of}@<Self{}> U>
  consteval @\libconcept{quantity_reference}@ auto operator[](this Self self, U u);

  template<typename Self, typename FwdQ, @\libconcept{delta_quantity}@ Q = std::remove_cvref_t<FwdQ>>
    requires @\exposconceptnc{QuantitySpecExplicitlyConvertibleTo}@<Q::quantity_spec, Self{}>
  constexpr @\libconcept{delta_quantity}@ auto operator()(this Self self, FwdQ&& q);

  template<typename Self, typename FwdQ, @\libconcept{point_quantity}@ Q = std::remove_cvref_t<FwdQ>>
    requires @\seebelownc@
  constexpr @\libconcept{point_quantity}@ auto operator()(this Self self, FwdQ&& q);

  template<@\libconcept{quantity_spec}@ Lhs, @\libconcept{quantity_spec}@ Rhs>
  friend consteval bool operator==(Lhs, Rhs);
};

}
\end{codeblock}

\pnum
\tcode{*} and \tcode{/} are the operations of scalars.
A product is formed only when at least one operand admits a character of the tensor order of a scalar,
that is, when at most one operand is declared a vector or a tensor;
a quotient only when the divisor admits a character of the tensor order of a scalar;
and \tcode{inverse}, \tcode{sqrt}, \tcode{cbrt}, and \tcode{pow<Num, Den>} with \tcode{Num != Den}
only for an operand that admits a character of the tensor order of a scalar\iref{qty.spec.concepts}.
\begin{note}
A vector or a tensor has no reciprocal,
and the product of two of them is one of several products that ISO 80000-2 defines, none of which is \tcode{*}.
A kind admits every character and is accepted: its representation type decides.
\end{note}

\indexlibrarymember{operator*}{quantity_spec}
\begin{itemdecl}
template<@\libconcept{quantity_spec}@ Lhs, @\libconcept{quantity_spec}@ Rhs>
  requires(!@\exposconceptnc{PointQuantitySpec}@<Lhs>) && (!@\exposconceptnc{PointQuantitySpec}@<Rhs>) &&
          (@\exposconceptnc{ScalarOperand}@<Lhs> || @\exposconceptnc{ScalarOperand}@<Rhs>)
friend consteval @\libconcept{quantity_spec}@ auto operator*(Lhs lhs, Rhs rhs);
\end{itemdecl}

\begin{itemdescr}
\pnum
\returns
\begin{codeblock}
@\exposidnc{close-kinds}@<Lhs{}, Rhs{}>(@\exposidnc{expr-multiply}@<derived_quantity_spec, struct dimensionless>(
  @\exposidnc{open-kinds}@(lhs), @\exposidnc{open-kinds}@(rhs)))
\end{codeblock}
\end{itemdescr}

\indexlibrarymember{operator/}{quantity_spec}
\begin{itemdecl}
template<@\libconcept{quantity_spec}@ Lhs, @\libconcept{quantity_spec}@ Rhs>
  requires(!@\exposconceptnc{PointQuantitySpec}@<Lhs>) && (!@\exposconceptnc{PointQuantitySpec}@<Rhs>) &&
          @\exposconceptnc{ScalarOperand}@<Rhs>
friend consteval @\libconcept{quantity_spec}@ auto operator/(Lhs lhs, Rhs rhs);
\end{itemdecl}

\begin{itemdescr}
\pnum
\returns
\begin{codeblock}
@\exposidnc{close-kinds}@<Lhs{}, Rhs{}>(@\exposidnc{expr-divide}@<derived_quantity_spec, struct dimensionless>(
  @\exposidnc{open-kinds}@(lhs), @\exposidnc{open-kinds}@(rhs)))
\end{codeblock}
\end{itemdescr}

\indexlibrarymember{operator[]}{quantity_spec}
\begin{itemdecl}
template<typename Self, @\libconcept{unit_of}@<Self{}> U>
consteval @\libconcept{quantity_reference}@ auto operator[](this Self self, U u);
\end{itemdecl}

\begin{itemdescr}
\pnum
\constraints
\begin{itemize}
\item
If \tcode{Self} is a specialization \tcode{point_<S>} and
\tcode{U} models \exposconceptnc{OffsetUnit}\iref{qty.unit.concepts},
\tcode{U::\exposid{point-origin} == natural_point_origin<S\{\}>} is \tcode{true}\iref{qty.pt.orig.ops}.
\item
If \tcode{Self} models neither \exposconceptnc{PointQuantitySpec} nor \exposconceptnc{DeltaQuantitySpec},
\tcode{U} does not model \exposconceptnc{OffsetUnit}.
\end{itemize}

\pnum
\returns
\tcode{\exposidnc{make-reference}(self, u)}:
for a point quantity specification, the reference of a point of \tcode{self} in \tcode{u};
otherwise, the reference of a delta of \tcode{\exposid{remove-delta}(self)} in \tcode{u}.

\pnum
\begin{note}
A numerical value in a unit that models \exposconceptnc{OffsetUnit}
does not say whether a delta or a point on the scale of the unit is meant,
so \tcode{QS[U]} is ill-formed for such a unit,
and \tcode{delta<QS>[U]} and \tcode{point<QS, PO>[U]} say which.
For any other unit, \tcode{delta<QS>[U]} and \tcode{QS[U]} are the same reference.
\tcode{point<QS>[U]} is measured from the natural point origin of \tcode{QS},
so a unit whose own origin is another zero is not bound to it:
\tcode{point<QS, PO>[U]} names the origin, and \tcode{U} gives only the step of the value.
\end{note}
\begin{example}
\begin{codeblock}
static_assert(delta<isq::length>[m] == isq::length[m]);
static_assert(@\libconcept{quantity_reference}@<decltype(delta<isq::thermodynamic_temperature>[deg_C])>);
static_assert(@\libconcept{quantity_reference}@<decltype(point<isq::thermodynamic_temperature, si::ice_point>[deg_C])>);
auto r1 = isq::thermodynamic_temperature[deg_C];          // error: \tcode{deg_C} has a point origin
auto r2 = point<isq::thermodynamic_temperature>[deg_C];   // error: \tcode{si::ice_point} is not the natural point origin
\end{codeblock}
\end{example}
\end{itemdescr}

\indexlibrarymember{operator()}{quantity_spec}
\begin{itemdecl}
template<typename Self, typename FwdQ, @\libconcept{delta_quantity}@ Q = std::remove_cvref_t<FwdQ>>
  requires @\exposconceptnc{QuantitySpecExplicitlyConvertibleTo}@<Q::quantity_spec, Self{}>
constexpr @\libconcept{delta_quantity}@ auto operator()(this Self self, FwdQ&& q);
\end{itemdecl}

\begin{itemdescr}
\pnum
\returns
\begin{codeblock}
quantity{std::forward<FwdQ>(q).@\exposidnc{numerical-value}@, @\exposidnc{make-reference}@(self, Q::unit)}
\end{codeblock}
\begin{note}
A point quantity specification converts no delta\iref{qty.spec.conv}.
\tcode{isq::thermodynamic_temperature(delta<deg_C>(5))} is a
\tcode{quantity<delta<isq::thermodynamic_temperature>[deg_C], int>},
and \tcode{delta<isq::length>(isq::height(1 * m))} a \tcode{quantity<isq::length[m], int>}.
\end{note}
\end{itemdescr}

\indexlibrarymemberexpos{quantity_spec}{point-target}
\indexlibrarymemberexpos{quantity_spec}{conversion-origin}
\begin{itemdecl}
template<@\libconcept{quantity_spec}@ QS, @\exposconceptnc{PointOrigin}@ PO>
consteval @\exposconceptnc{PointQuantitySpec}@ auto @\exposidnc{point-target}@(QS qs, PO po);      // \expos

template<@\libconcept{quantity_spec}@ QS, @\exposconceptnc{PointOrigin}@ PO>
consteval @\exposconceptnc{PointOrigin}@ auto @\exposidnc{conversion-origin}@(QS qs, PO po);      // \expos
\end{itemdecl}

\begin{itemdescr}
\pnum
\returns
For \exposid{point-target}, the point quantity specification
that a point measured from \tcode{po} is converted to by \tcode{qs}:
\tcode{qs} if \tcode{QS} is a specialization \tcode{point_<S, PO2>};
otherwise, \tcode{\exposid{normalized-point-spec}(\exposid{remove-point}(qs), po)}\iref{point.qty}
if \tcode{\libconcept{point_origin_for}<PO, \exposid{remove-point}(qs)>} is satisfied;
otherwise, \tcode{\exposid{as-point}(qs)}.
For \exposid{conversion-origin}, the point origin from which the result of that conversion is measured:
\tcode{po} if \tcode{\libconcept{point_origin_for}<PO, \exposid{remove-point}(qs)>} is satisfied and
\tcode{QS} is not a specialization \tcode{point_<S, PO2>};
otherwise, \tcode{\exposid{point-spec-origin}(\exposid{as-point}(qs))}\iref{point.qty}.
\end{itemdescr}

\indexlibrarymember{operator()}{quantity_spec}
\begin{itemdecl}
template<typename Self, typename FwdQ, @\libconcept{point_quantity}@ Q = std::remove_cvref_t<FwdQ>>
  requires @\seebelownc@
constexpr @\libconcept{point_quantity}@ auto operator()(this Self self, FwdQ&& q);
\end{itemdecl}

\begin{itemdescr}
\pnum
Let \tcode{po} be \tcode{\exposid{conversion-origin}(Self\{\}, Q::point_origin)}.

\pnum
\returns
\begin{codeblock}
quantity{@\exposidnc{remove-point}@(self)(std::forward<FwdQ>(q).quantity_from(po)), po}
\end{codeblock}

\pnum
\remarks
The expression in the \fakegrammarterm{requires-clause} is equivalent to:
\begin{codeblock}
(!@\exposconceptnc{DeltaQuantitySpec}@<Self>) &&
  @\exposconceptnc{QuantitySpecExplicitlyConvertibleTo}@<Q::quantity_spec, @\exposidnc{point-target}@(Self{}, Q::point_origin)> &&
  requires(const Q& qp) { qp.quantity_from(@\exposidnc{conversion-origin}@(Self{}, Q::point_origin)); }
\end{codeblock}

\pnum
\begin{note}
A point converts to a point of \tcode{Self},
whether \tcode{Self} is a point quantity specification or the quantity specification of its scale,
up or down one branch of the hierarchy, and only from a compatible origin\iref{qty.spec.conv}.
The result keeps the origin of the point when that origin is a point origin for the scale of \tcode{Self}
and \tcode{Self} names no other origin;
otherwise, it is measured from the origin of the target.
\tcode{delta<QS>} is the quantity specification of a delta and converts no point.
\end{note}
\begin{example}
\begin{codeblock}
inline constexpr struct mean_sea_level final : absolute_point_origin<isq::altitude> {} mean_sea_level;

auto p1 = isq::altitude(point<m>(42));                      // measured from \tcode{natural_point_origin<kind_of<isq::length>>}
auto p2 = isq::height(point<isq::altitude, mean_sea_level>[m](42));
                                                           // \tcode{quantity<point<isq::height, mean_sea_level>[m], int>}
auto p3 = isq::length(point<isq::height>[m](1.));           // measured from \tcode{natural_point_origin<isq::length>}
auto p4 = isq::length(point<isq::altitude, mean_sea_level>[m](5.));
                                                           // error: \tcode{mean_sea_level} is not compatible with
                                                           // \tcode{natural_point_origin<isq::length>}
auto p5 = point<isq::length>(42 * m);                       // error: a delta never converts to a point
auto p6 = isq::width(point<isq::height>[m](42));            // error: a point converts along one branch only
auto p7 = delta<isq::altitude>(point<m>(42));               // error: \tcode{delta<isq::altitude>} converts no point
\end{codeblock}
\end{example}
\end{itemdescr}

\indexlibrarymember{operator==}{quantity_spec}
\begin{itemdecl}
template<@\libconcept{quantity_spec}@ Lhs, @\libconcept{quantity_spec}@ Rhs>
friend consteval bool operator==(Lhs, Rhs);
\end{itemdecl}

\begin{itemdescr}
\pnum
\returns
\tcode{std::is_same_v<Lhs, Rhs>}.
\end{itemdescr}

\indexlibrarymember{inverse}{quantity_spec}
\begin{itemdecl}
template<@\libconcept{quantity_spec}@ Q>
  requires(!@\exposconceptnc{PointQuantitySpec}@<Q>) && @\exposconceptnc{ScalarOperand}@<Q>
consteval @\libconcept{quantity_spec}@ auto inverse(Q q);
\end{itemdecl}

\begin{itemdescr}
\pnum
\returns
\tcode{dimensionless / q}.
\end{itemdescr}

\indexlibrarymember{pow}{quantity_spec}
\begin{itemdecl}
template<std::intmax_t Num, std::intmax_t Den = 1, @\libconcept{quantity_spec}@ Q>
  requires(Den != 0) && (!@\exposconceptnc{PointQuantitySpec}@<Q>) && (@\exposconceptnc{ScalarOperand}@<Q> || Num == Den)
consteval @\libconcept{quantity_spec}@ auto pow(Q q);
\end{itemdecl}

\begin{itemdescr}
\pnum
\returns
\begin{codeblock}
@\exposidnc{close-kinds}@<Q{}>(
  @\exposidnc{expr-pow}@<Num, Den, derived_quantity_spec, struct dimensionless>(@\exposidnc{open-kinds}@(q)));
\end{codeblock}
\end{itemdescr}

\indexlibrarymember{sqrt}{quantity_spec}
\begin{itemdecl}
template<@\libconcept{quantity_spec}@ Q>
  requires(!@\exposconceptnc{PointQuantitySpec}@<Q>) && @\exposconceptnc{ScalarOperand}@<Q>
consteval @\libconcept{quantity_spec}@ auto sqrt(Q q);
\end{itemdecl}

\begin{itemdescr}
\pnum
\returns
\tcode{pow<1, 2>(q)}.
\end{itemdescr}

\indexlibrarymember{cbrt}{quantity_spec}
\begin{itemdecl}
template<@\libconcept{quantity_spec}@ Q>
  requires(!@\exposconceptnc{PointQuantitySpec}@<Q>) && @\exposconceptnc{ScalarOperand}@<Q>
consteval @\libconcept{quantity_spec}@ auto cbrt(Q q);
\end{itemdecl}

\begin{itemdescr}
\pnum
\returns
\tcode{pow<1, 3>(q)}.
\end{itemdescr}

\begin{example}
\begin{codeblock}
static_assert(@\libconcept{quantity_spec}@<decltype(isq::velocity * isq::duration)>);
static_assert(pow<1>(isq::velocity) == isq::velocity);
auto q1 = isq::velocity * isq::displacement;  // error: two vector operands
auto q2 = isq::duration / isq::velocity;      // error: a vector divisor
auto q3 = pow<2>(isq::velocity);              // error: not a scalar
\end{codeblock}
where \tcode{isq::velocity} and \tcode{isq::displacement} are vector quantities.
\end{example}

\indexlibrarymember{magnitude}{quantity_spec}
\begin{itemdecl}
template<@\libconcept{quantity_spec}@ Q>
  requires(!@\exposconceptnc{PointQuantitySpec}@<Q>) && @\exposconceptnc{NonScalarOperand}@<Q>
consteval @\libconcept{quantity_spec}@ auto magnitude(Q q);
\end{itemdecl}

\begin{itemdescr}
\pnum
\returns
\tcode{magnitude(\exposid{remove-delta}(q))} if \tcode{Q} models \exposconceptnc{DeltaQuantitySpec};
otherwise, \tcode{q} if \tcode{Q} models \exposconceptnc{QuantityKindSpec};
otherwise, \tcode{\exposid{magnitude-t}<Q>\{\}}\iref{magnitude.qty}.

\pnum
\begin{note}
A kind has no character of its own,
so the magnitude of a quantity of a kind is a quantity of that kind:
\tcode{magnitude(kind_of<isq::speed>)} is \tcode{kind_of<isq::speed>}.
\tcode{magnitude(isq::speed)} and \tcode{magnitude(point<isq::velocity>)} are ill-formed.
\end{note}
\end{itemdescr}

\rSec3[qty.spec.hier.algos]{Hierarchy algorithms}

\rSec4[qty.spec.conv]{Conversion}

\pnum
In \ref{qty.spec.conv}, let \tcode{from} and \tcode{to} denote
the arguments of the conversion being determined,
objects of types modeling \libconcept{quantity_spec}, and
let \tcode{N} and \tcode{M} denote
objects of types modeling \exposconceptnc{NamedQuantitySpec}.

\pnum
The \defnx{parent chain}{parent chain} of \tcode{N} is the sequence
\tcode{N}, \tcode{N.\exposid{parent}}, \tcode{N.\exposidnc{parent}.\exposidnc{parent}}, \ldots,
ending with the first element \tcode{E} for which \tcode{E.\exposidnc{parent}} is not a valid expression.
That last element is the \defnx{root}{root} of \tcode{N}.
\tcode{N} is an \defnx{ancestor}{ancestor} of \tcode{M} if
\tcode{N} appears in the parent chain of \tcode{M};
\tcode{N} is a \defnx{proper ancestor}{proper ancestor} of \tcode{M} if
it is an ancestor of \tcode{M} and \tcode{N != M}.
\tcode{N} and \tcode{M} are \defnx{in the same tree}{tree, same} if they have the same root.

\pnum
The \defnx{kind root}{kind root} of \tcode{N} is
\tcode{\exposid{get-kind-tree-root}(N)}.

\pnum
The \defnx{carried kinds}{carried kinds} of a named quantity \tcode{N} are a mapping from named quantities
to rational numbers, called multiplicities, zero for all but finitely many:
every element \tcode{E} of the parent chain of \tcode{N} such that
the specialization of \tcode{named_quantity_spec} \tcode{E} is derived from
has a template argument equal to \tcode{subkind}
(each such element being a \defnx{declared kind}{declared kind}), each with multiplicity 1,
plus the carried kinds of \tcode{N.\exposidnc{equation}}
if that is a valid expression.
The carried kinds of a quantity specification \tcode{Q} with factors (defined below) $(\tcode{N}_i, e_i)$ are
the sum over $i$ of the carried kinds of $\tcode{N}_i$ with every multiplicity multiplied by $e_i$,
where multiplicities of the same named quantity add.
\begin{example}
\tcode{isq::phase_angle} carries \tcode{isq::phase_angle} and \tcode{isq::angular_measure};
\tcode{isq::angular_velocity} carries \tcode{isq::angular_measure} through its equation;
\tcode{isq::speed} carries nothing.
\end{example}

\pnum
The \defnx{defining ancestor}{defining ancestor} of \tcode{N} is
the first element \tcode{A} of the parent chain of \tcode{N}
for which \tcode{A.\exposid{equation}} is a valid expression and
either \tcode{A.\exposidnc{parent}} is not a valid expression or
\tcode{A.\exposidnc{parent}.\exposidnc{equation}} is not a valid expression or
\tcode{A.\exposidnc{parent}.\exposidnc{equation} != A.\exposidnc{equation}}.
If there is no such element, \tcode{N} has no defining ancestor.
\begin{note}
A quantity whose equation is inherited from its parent has the same defining ancestor as
its parent.
\end{note}

\pnum
The \defnx{depth}{depth} of a quantity specification \tcode{Q} is
\begin{itemize}
\item
the depth of \tcode{Q.\exposid{quantity-spec}} if \tcode{Q} models \exposconceptnc{QuantityKindSpec},
\item
otherwise, \tcode{0} if \tcode{Q} is a named quantity without a defining ancestor,
\item
otherwise, one more than the depth of \tcode{A.\exposidnc{equation}} if
\tcode{Q} is a named quantity with defining ancestor \tcode{A}, and
\item
otherwise, the greatest depth among the factors of \tcode{Q} (defined below) if
\tcode{Q} models \exposconceptnc{DerivedQuantitySpec}, or
\tcode{0} if it has no factors.
\end{itemize}

\pnum
The \defnx{conversion rank}{conversion rank} of a conversion is
a value of the exposition-only enumeration
\begin{codeblock}
enum class @\exposidnc{rank}@ {  // \expos
  none, cast, explicit_only, implicit
};
\end{codeblock}
whose enumerators are ordered as declared, \tcode{none} being the lowest.
Where a conversion is composed of several conversion operations, its rank is the lowest rank of any of them.
Where this subclause leaves a choice of pairings, of self-cancellings, or of the factor to expand,
the rank of the conversion is the highest rank obtainable over all such choices.

\pnum
For a named quantity \tcode{K} and a quantity specification \tcode{Q},
let $m$ be the multiplicity of \tcode{K} in the carried kinds of \tcode{Q};
the positive part of $m$ is $\max(m, 0)$ and the negative part is $\max(-m, 0)$.
The carried kinds of \tcode{Q1} \defnx{exceed}{carried kinds, exceed} those of \tcode{Q2} if,
for some named quantity \tcode{K},
the positive part or the negative part of its multiplicity in \tcode{Q1} is
greater than the corresponding part in \tcode{Q2}.
\tcode{\exposid{kind-boundary-rank}(Q1, Q2, k)}, where \tcode{k} is a \tcode{bool} value, denotes
\tcode{none} if the carried kinds of \tcode{Q1} exceed those of \tcode{Q2} and
the carried kinds of \tcode{Q2} exceed those of \tcode{Q1},
\tcode{explicit_only} if only the carried kinds of \tcode{Q1} exceed those of \tcode{Q2},
or if only the carried kinds of \tcode{Q2} exceed those of \tcode{Q1} and \tcode{k} is \tcode{false},
and \tcode{implicit} otherwise.
When \tcode{k} is \tcode{true}, a negative contribution that a factor of the numerator of \tcode{Q2}
makes through its \exposidnc{equation} is left out of the sum and counted in the negative part on its own.
\begin{note}
The multiplicities that the numerator and the denominator of a quantity specification contribute
to one declared kind add up before the comparison:
\tcode{isq::angular_displacement / isq::duration * isq::radius / isq::angular_measure / isq::speed}
carries no \tcode{isq::angular_measure}.
From a kind, a gained kind is not a loss even where the target nets it against one it already has:
a source kind that carries a declared kind once in its numerator converts implicitly to
a target whose first numerator factor carries that declared kind
and whose second numerator factor, defined over the first, takes it away again.
\end{note}
\ednote{
The example of this rule was the kind of kg m\textsuperscript{2} rad/s converting implicitly to
\tcode{isq::angular_velocity * isq::moment_of_inertia}.
P3045R11 makes that product ill-formed, as both factors are vectors or tensors\iref{qty.spec.ops}.
The paper marks a new example as pending; the note above states the rule without one.
}

\pnum
\tcode{\exposid{tree-rank}(N, M, k)}, where \tcode{k} is a \tcode{bool} value, denotes
the rank determined by the first of the following that applies:
\begin{itemize}
\item
If \tcode{N == M}, the rank is \tcode{implicit}.
\item
Otherwise, if \tcode{N} and \tcode{M} are not in the same tree, the rank is \tcode{none}.
\item
Otherwise, if \tcode{M} is a proper ancestor of \tcode{N},
the rank is \tcode{implicit} if \tcode{N} and \tcode{M} have the same kind root,
and \tcode{explicit_only} otherwise.
\item
Otherwise, if \tcode{N} is a proper ancestor of \tcode{M},
the rank is \tcode{implicit} if \tcode{k} is \tcode{true}, and \tcode{explicit_only} otherwise.
\item
Otherwise, the rank is \tcode{cast} if \tcode{N} and \tcode{M} have the same kind root,
and \tcode{none} otherwise.
\end{itemize}

\pnum
\tcode{\exposid{conversion-rank}(from, to, from-kind)} denotes
the rank determined below for converting
a quantity of specification \tcode{from} into a quantity of specification \tcode{to},
where \tcode{from-kind} is \tcode{true} when the conversion is reached from a kind
as specified below.
\tcode{\exposidnc{conversion-rank}(from, to)} denotes
\tcode{\exposidnc{conversion-rank}(from, to, false)}.

\pnum
The rank of converting \tcode{from} to \tcode{to} is determined by
the first of the following paragraphs that applies.

\pnum
If \tcode{from == to}, the rank is \tcode{implicit}.
Otherwise, if \tcode{decltype(from)} or \tcode{decltype(to)} models \exposconceptnc{DeltaQuantitySpec}\iref{delta.qty},
the rank is
\tcode{\exposidnc{conversion-rank}(\exposidnc{remove-delta}(from), \exposidnc{remove-delta}(to), from-kind)}.
Otherwise, if exactly one of \tcode{decltype(from)} and \tcode{decltype(to)}
models \exposconceptnc{PointQuantitySpec}, the rank is \tcode{none}.
Otherwise, if both model \exposconceptnc{PointQuantitySpec}:
if \tcode{\exposidnc{same-absolute-point-origins}(\exposidnc{point-spec-origin}(from), \exposidnc{point-spec-origin}(to))}
is \tcode{false}\iref{qty.same.abs.pt.origs,point.qty}, the rank is \tcode{none};
otherwise, let \tcode{r} be
\tcode{\exposidnc{conversion-rank}(from.\exposidnc{quantity-spec}, to.\exposidnc{quantity-spec})};
the rank is \tcode{none} if \tcode{r} is \tcode{cast}, and \tcode{r} otherwise.
Otherwise, if \tcode{from.dimension != to.dimension}, the rank is \tcode{none}.
Otherwise, if \tcode{\exposid{characters-overlap}(from, to)} is \tcode{false}\iref{qty.spec.concepts},
the rank is \tcode{none}.
\begin{note}
A point never converts to a delta nor the other way round, even with \tcode{quantity_cast}.
A point converts to a point of another origin only when both share their absolute point origin;
the natural point origins of a quantity, of its ancestors and descendants, and of its kind compare equal\iref{qty.natural.pt.orig},
so they count as one origin.
A point converts up or down one branch of the hierarchy only:
the natural origins of two siblings are the zeros of different scales,
along axes that can point in different directions,
so \tcode{point<isq::height>} and \tcode{point<isq::width>} do not convert to each other,
even with \tcode{quantity_cast}, while \tcode{isq::height} and \tcode{isq::width} do, with a cast.
A quantity does not convert to a quantity of another character, even with \tcode{quantity_cast}:
a vector or a tensor and a scalar, a vector and a tensor, and a complex and a real quantity.
A kind admits every character and passes this check;
the scalar that belongs to a vector is its magnitude\iref{magnitude.qty}.
\end{note}
\begin{example}
\begin{codeblock}
static_assert(@\exposid{implicitly-convertible}@(delta<isq::height>, isq::length));
static_assert(@\exposid{implicitly-convertible}@(point<isq::height>, point<isq::length>));
static_assert(@\exposid{explicitly-convertible}@(point<isq::length>, point<isq::height>));
static_assert(!@\exposid{castable}@(point<isq::height>, point<isq::width>));
static_assert(!@\exposid{castable}@(point<isq::height>, isq::height));
static_assert(@\exposid{implicitly-convertible}@(point<isq::altitude, ground>, point<isq::altitude, mean_sea_level>));
static_assert(!@\exposid{castable}@(point<isq::altitude, other_zero>, point<isq::altitude, mean_sea_level>));
static_assert(!@\exposid{castable}@(isq::velocity, isq::speed));
static_assert(!@\exposid{castable}@(isq::speed, isq::velocity));
static_assert(@\exposid{implicitly-convertible}@(kind_of<isq::length / isq::time>, isq::velocity));
\end{codeblock}
where \tcode{mean_sea_level} and \tcode{other_zero} are two absolute point origins of \tcode{isq::altitude},
\tcode{ground} is a relative point origin defined from \tcode{mean_sea_level}\iref{qty.pt.orig},
and \tcode{isq::velocity} is a vector quantity.
\end{example}

\pnum
If \tcode{decltype(from)} models \exposconceptnc{QuantityKindSpec},
let \tcode{K} be \tcode{from.\exposidnc{quantity-spec}}.
The rank is \tcode{\exposidnc{conversion-rank}(K, to.\exposidnc{quantity-spec}, true)}
if \tcode{decltype(to)} models \exposconceptnc{QuantityKindSpec}, and
\tcode{\exposidnc{conversion-rank}(K, to, true)} otherwise.
\begin{note}
A kind denotes a quantity that can be any quantity of the tree below its kind root,
so a downcast from it is implicit, while leaving a declared kind is not.
\end{note}

\pnum
If \tcode{decltype(to)} models \exposconceptnc{QuantityKindSpec},
the rank is \tcode{\exposidnc{conversion-rank}(from, to.\exposidnc{quantity-spec})}.

\pnum
If \tcode{decltype(from)} and \tcode{decltype(to)} both model \exposconceptnc{NamedQuantitySpec},
let \tcode{r} be \tcode{\exposidnc{tree-rank}(from, to, from-kind)}.
Let \tcode{b} be \tcode{\exposidnc{kind-boundary-rank}(from, to, from-kind)}.
The rank is \tcode{explicit_only} if \tcode{b} is \tcode{none}
and \tcode{r} is \tcode{implicit} or \tcode{explicit_only},
and the lower of \tcode{r} and \tcode{b} otherwise.
\begin{note}
Equations participate in the conversion between two named quantities
only through the carried kinds.
\tcode{isq::phase_coefficient}, defined as \tcode{phase_angle / path_length},
is a child of \tcode{isq::propagation_coefficient} and converts to it with \tcode{explicit_only},
because it carries \tcode{isq::phase_angle} and its parent does not.
Two quantities of one tree, neither an ancestor of the other,
that each carry a declared kind the other lacks do not convert at all;
on a hierarchy edge the two quantities are ancestor and descendant rather than unrelated kinds,
so the rank is \tcode{explicit_only} rather than \tcode{none}.
\end{note}

\pnum
Otherwise, the rank is determined by the following procedure.
Let $F$ and $T$ be lists of \defnx{factors}{factor},
each factor being a named quantity, a non-zero rational exponent, and
whether the factor is \defnx{open}{factor, open}.
Initially, $F$ holds the factors of \tcode{from} and $T$ the factors of \tcode{to},
where the factors of a quantity specification \tcode{Q} are
\begin{itemize}
\item
the single factor $(\tcode{Q}, 1)$ if \tcode{Q} is a named quantity other than \tcode{dimensionless},
\item
no factors if \tcode{Q} is \tcode{dimensionless}, and
\item
otherwise one factor per element of \tcode{decltype(Q)::\exposidnc{num}} with its exponent and
one factor per element of \tcode{decltype(Q)::\exposidnc{den}} with its exponent negated,
where the exponent of an element that is not a specialization of \exposid{power} is 1.
An element equal to \tcode{dimensionless} yields no factor.
A factor is open if its element is a kind factor \iref{derived.qty},
and its named quantity is then the template argument of the specialization of \tcode{kind_of_};
otherwise the factor is not open.
\end{itemize}
\begin{note}
Only a kind factor of a derived quantity specification yields an open factor.
A \tcode{kind_of<K>} on either side of the conversion is handled above,
where \tcode{from-kind} records that the source is a kind.
\end{note}
Factors with a positive exponent are \defnx{numerator factors}{numerator factor};
the others are \defnx{denominator factors}{denominator factor}.
The depth of a factor is the depth of its named quantity.
Whenever the exponent of a factor becomes zero, the factor is removed from its list.
Replacing a factor $(\tcode{N}, e)$ by an expression $E$ denotes
removing it and inserting one factor $(\tcode{M}, e \cdot x)$ for each factor $(\tcode{M}, x)$ of
$E$, each inserted factor being open if the replaced factor is open,
after which factors of the same list with the same named quantity
that are either both open or both not open are merged by adding their exponents.

\pnum
Before any conversion operation, factors common to both lists are canceled:
for every pair of a factor $(\tcode{N}, e)$ in $F$ and a factor $(\tcode{N}, f)$ in $T$
such that $e$ and $f$ have the same sign and the two factors are either both open or both not open,
the one of $e$ and $f$ with the smaller magnitude is subtracted from both exponents.
Let $F_0$ and $T_0$ denote the lists at this point.
If $F_0$ consists of a single factor $(\tcode{N}, 1)$ and $T_0$ of a single factor $(\tcode{M}, 1)$,
the rank is \tcode{\exposidnc{conversion-rank}(N, M, k)},
where \tcode{k} is \tcode{true} if \tcode{from-kind} is \tcode{true} or the factor of $F_0$ is open,
and the conversion operations below are not applied.
\begin{note}
A factor common to both sides of a conversion, with exponents of the same sign, does not change its rank:
\tcode{isq::phase_coefficient * isq::mass} converts to \tcode{isq::propagation_coefficient * isq::mass}
with the rank of \tcode{isq::phase_coefficient} to \tcode{isq::propagation_coefficient}.
\end{note}

\pnum
Each \defnx{conversion operation}{operation!conversion} transforms the lists and has a rank:
\begin{itemize}
\item
\defnx{Pairing}{pairing}: given a factor $(\tcode{N}, e)$ in $F$ and
a factor $(\tcode{M}, f)$ in $T$ such that $e$ and $f$ have the same sign and
\tcode{N} and \tcode{M} are in the same tree,
let $c$ be the one of $e$ and $f$ with the smaller magnitude.
$c$ is subtracted from both exponents.
The rank is \tcode{\exposidnc{tree-rank}(N, M, k)},
where \tcode{k} is \tcode{true} if \tcode{from-kind} is \tcode{true} or the factor of $F$ is open.
\begin{note}
An open factor of $F$ pairs as a kind does, and an open factor of $T$ as its named quantity does.
\tcode{kind_of<isq::length> / isq::duration} converts to a child of \tcode{isq::speed}
defined over \tcode{isq::height} with \tcode{implicit}, while
\tcode{isq::length / isq::duration} converts to it with \tcode{explicit_only}.
\end{note}
\item
\defnx{Self-cancelling}{self-cancelling}: given a numerator factor $(\tcode{N}, e)$ and
a denominator factor $(\tcode{M}, f)$ of the same list such that
\tcode{N.dimension == M.dimension},
let $c$ be the smaller of $e$ and $-f$.
$e$ is reduced by $c$ and $f$ is increased by $c$.
The rank is \tcode{implicit} for factors of $F$, and for factors of $T$ it is \tcode{implicit}
if \tcode{from-kind} is \tcode{true} and \tcode{explicit_only} otherwise.
\item
\defnx{Expanding}{expanding}: given a factor $(\tcode{N}, e)$ in $F$ or in $T$
such that \tcode{N} has a defining ancestor \tcode{A},
the factor is replaced by \tcode{A.\exposidnc{equation}}.
The rank is \tcode{implicit} for a factor of $F$ and
\tcode{\exposidnc{tree-rank}(A, N, from-kind)} for a factor of $T$.
\end{itemize}

\pnum
The following is repeated until $F$ and $T$ hold the same factors,
an iteration applies no conversion operation, or a conversion operation ranks \tcode{none}:
\begin{itemize}
\item
If $F$ consists of a single factor $(\tcode{N}, e)$,
\tcode{N} has a defining ancestor \tcode{A}, and
$T$ holds, in any order, exactly the factors of \tcode{A.\exposidnc{equation}}
with every exponent multiplied by $e$, whether a factor is open being disregarded,
expanding is applied to that factor.
Otherwise, if $T$ consists of a single factor $(\tcode{M}, f)$,
\tcode{M} has a defining ancestor \tcode{A}, and
$F$ holds, in any order, exactly the factors of \tcode{A.\exposidnc{equation}}
with every exponent multiplied by $f$, whether a factor is open being disregarded,
expanding is applied to that factor.
\begin{note}
A quantity is thereby convertible from and to its own equation
even when that equation contains factors of the quantity's own tree.
\end{note}
\item
Otherwise, if $F$ consists of a single factor $(\tcode{N}, e)$ such that
\tcode{N.\exposidnc{equation}} is a valid expression with at least one factor and
the named quantity of every factor of it is in the same tree as \tcode{N},
expanding is applied to that factor;
otherwise, if $T$ consists of a single such factor $(\tcode{M}, f)$, expanding is applied to it.
\begin{note}
In the ISQ this occurs in the tree of \tcode{dimensionless}, where every quantity of dimension one
that is not a declared kind is a sibling of every other, and for \tcode{isq::intrinsic_carrier_density},
defined as \tcode{sqrt(electron_density * hole_density)} within the tree of \tcode{number_density}.
A product of such quantities is treated as that product rather than as a sibling.
\end{note}
\item
Otherwise, if a pairing applies that is not deferred, such a pairing is applied.
A pairing of \tcode{N} and \tcode{M} is deferred while its rank would be \tcode{cast} and
either list holds a factor whose named quantity \tcode{P} is not in the tree of \tcode{N} and
for which \tcode{P.\exposidnc{equation}} is a valid expression with
a factor whose named quantity is in the tree of \tcode{N},
and, when \tcode{from-kind} is \tcode{true}, while its rank would be \tcode{none} and either list holds such a factor,
or one whose named quantity \tcode{P} is not in the tree of \tcode{N} and
for which \tcode{P.\exposidnc{equation}} is a valid expression with a factor
whose named quantity, recursively, has such an equation.
\begin{note}
The sideways relation between two siblings is decided once every factor of their tree is present.
The named quantity of a factor that defers a pairing has an equation and therefore a defining ancestor,
so expanding applies to that factor;
a deferred pairing is applied once it is no longer deferred.
\end{note}
\item
Otherwise, if $F$ holds a factor $(\tcode{N}, e)$ and $T$ a factor $(\tcode{M}, f)$ such that
$e$ and $f$ have the same sign, \tcode{N} and \tcode{M} are not in the same tree,
\tcode{N.dimension == M.dimension}, and neither \tcode{N.\exposidnc{equation}} nor
\tcode{M.\exposidnc{equation}} is a valid expression with a factor whose named quantity is
in the tree of the other, no conversion operation is applied.
\begin{note}
Two roots of one dimension that the tree rule keeps apart stay apart inside products:
\tcode{isq::frequency * isq::width} converts to \tcode{isq::activity * isq::length} with rank \tcode{none}.
A root whose equation mentions the other's tree is left to the equations:
\tcode{isq::radiant_intensity * isq::solid_angular_measure} converts to \tcode{isq::radiant_flux} with rank \tcode{implicit}.
\end{note}
\item
Otherwise, if self-cancelling applies to a list whose counterpart is empty,
it is applied to that list.
\item
Otherwise, if expanding applies,
it is applied to a factor of greatest depth among those it applies to,
taken from the first of the following groups that holds such a factor:
numerator factors of $F$, denominator factors of $F$,
numerator factors of $T$, denominator factors of $T$.
Any factor of greatest depth in that group may be chosen.
When a pairing was deferred in this step and a factor of $T$ that expanding applies to
has a depth less than the greatest depth among the factors of $F$ and $T$,
expanding is applied instead to a factor of $T$ of least depth among those,
numerator factors before denominator factors.
\begin{note}
The rank is the highest over the factors of greatest depth in that group,
so the names of the quantities do not affect it.
A deferred pairing waits for the factors whose equations mention its tree,
and the factor of least depth is the one whose own equation brings that tree in:
it meets the factors that wait before a deeper factor defined over it cancels them.
\end{note}
\end{itemize}

\pnum
When the repetition stops, the rank is \tcode{none}
if a conversion operation ranked \tcode{none} or $F$ and $T$ do not hold the same factors,
and otherwise the lower of the lowest rank of the conversion operations applied
(\tcode{implicit} if none was applied) and
\tcode{\exposidnc{kind-boundary-rank}($P_F$, $P_T$, from-kind)},
where $P_F$ and $P_T$ are the products of the factors of $F_0$ and of $T_0$, respectively.
For that comparison, the carried kinds of $P_T$ do not exceed those of $P_F$
by a named quantity \tcode{K} of which the named quantity of an open factor of $F_0$ is an ancestor.
\begin{note}
A declared kind gained below an open factor is not created by the conversion;
the open factor commits to it, as a kind does.
\tcode{kind_of<isq::length> * isq::mass} converts to \tcode{fluid_head * isq::mass} with \tcode{implicit}
for a declared kind \tcode{fluid_head} below \tcode{isq::length}, while
\tcode{isq::length * isq::mass} converts to it with \tcode{explicit_only}.
\end{note}
\begin{note}
A declared kind is a boundary in both directions, and the carried kinds are compared
between the two sides once their common factors have canceled, so the result does not depend on the conversion operations.
When the conversion is from a kind, gaining a declared kind does not lower the rank,
which is what lets an uncommitted kind enter a declared kind.
\tcode{isq::arc_length / isq::radius} converts to \tcode{isq::angular_measure} with
\tcode{explicit_only}, \tcode{isq::angular_measure * isq::radius} converts to
\tcode{isq::arc_length} with \tcode{explicit_only}, and
\tcode{isq::angular_displacement / isq::duration / isq::frequency} converts to \tcode{dimensionless} with
\tcode{explicit_only} but to \tcode{isq::angular_measure} with \tcode{implicit}.
\tcode{dimensionless / isq::angular_measure} converts to \tcode{dimensionless} with
\tcode{explicit_only}, since a denominator angle is lost rather than created.
\end{note}
\begin{note}
The procedure terminates because every conversion operation either removes a factor or
replaces a factor by factors of strictly smaller depth,
so the multiset of depths of the factors decreases in a well-founded order.
Pairing takes precedence so that two quantities of one tree relate as the tree says even
when they occur as factors, and the deepest factor is reduced first so that
both sides meet at the most specific common equation.
\end{note}
\begin{example}
\begin{codeblock}
static_assert(@\exposid{implicitly-convertible}@(isq::width * isq::height, isq::area));
  // \tcode{area} expands into \tcode{pow<2>(length)} (\tcode{implicit}),
  // then \tcode{width} and \tcode{height} pair with \tcode{length} (\tcode{implicit})
static_assert(!@\exposid{implicitly-convertible}@(inverse(isq::time), isq::frequency));
static_assert(@\exposid{explicitly-convertible}@(inverse(isq::time), isq::frequency));
  // \tcode{frequency} expands into \tcode{inverse(period_duration)} (\tcode{implicit}),
  // then \tcode{time} pairs with \tcode{period_duration} (\tcode{explicit_only})
static_assert(!@\exposid{implicitly-convertible}@(isq::angular_displacement / isq::duration * isq::radius /
                                      isq::speed, dimensionless));
static_assert(@\exposid{explicitly-convertible}@(isq::angular_displacement / isq::duration * isq::radius /
                                     isq::speed, dimensionless));
  // every conversion operation ranks \tcode{implicit}: \tcode{angular_displacement}, that is
  // \tcode{rotational_displacement}, expands
  // into \tcode{path_length / radius}, \tcode{speed} into \tcode{length / duration},
  // and with $T$ empty \tcode{path_length} self-cancels against \tcode{length};
  // the carried kinds of \tcode{from} include \tcode{angular_measure} and those of
  // \tcode{to} do not, so the rank is \tcode{explicit_only}
\end{codeblock}
\end{example}

\indexlibraryglobalexpos{implicitly-convertible}
\begin{itemdecl}
consteval bool @\exposid{implicitly-convertible}@(quantity_spec auto from, quantity_spec auto to);
\end{itemdecl}

\begin{itemdescr}
\pnum
\returns
\tcode{\exposidnc{conversion-rank}(from, to) == \exposidnc{rank}::implicit}.
\end{itemdescr}

\indexlibraryglobalexpos{explicitly-convertible}
\begin{itemdecl}
consteval bool @\exposid{explicitly-convertible}@(quantity_spec auto from, quantity_spec auto to);
\end{itemdecl}

\begin{itemdescr}
\pnum
\returns
\tcode{\exposidnc{conversion-rank}(from, to) >= \exposidnc{rank}::explicit_only}.
\end{itemdescr}

\indexlibraryglobalexpos{castable}
\begin{itemdecl}
consteval bool @\exposid{castable}@(quantity_spec auto from, quantity_spec auto to);
\end{itemdecl}

\begin{itemdescr}
\pnum
\returns
\tcode{\exposidnc{conversion-rank}(from, to) >= \exposidnc{rank}::cast}.
\end{itemdescr}

\indexlibraryglobalexpos{interconvertible}
\begin{itemdecl}
consteval bool @\exposid{interconvertible}@(quantity_spec auto qs1, quantity_spec auto qs2);
\end{itemdecl}

\begin{itemdescr}
\pnum
\returns
\tcode{\exposid{implicitly-convertible}(qs1, qs2) \&\& \exposid{implicitly-convertible}(qs2, qs1)}.
\end{itemdescr}

\rSec4[qty.get.kind]{Get kind}

\indexlibrarymemberexpos{quantity_spec}{get-kind-tree-root}
\begin{itemdecl}
template<@\libconcept{quantity_spec}@ Q>
consteval @\libconcept{quantity_spec}@ auto @\exposidnc{get-kind-tree-root}@(Q q);  // \expos
\end{itemdecl}

\begin{itemdescr}
\pnum
\returns
\begin{itemize}
\item
If \tcode{\exposconceptnc{QuantityKindSpec}<Q>} is \tcode{true},
returns \tcode{\exposidnc{remove-kind}(q)}.
\item
Otherwise, if \tcode{Q} is a specialization of \exposid{magnitude-t}\iref{magnitude.qty},
returns \tcode{\exposidnc{get-kind-tree-root}(Q::\exposidnc{quantity-spec})}.
\item
Otherwise, if
\tcode{\exposidnc{is-derived-from-specialization-of}<Q, named_quantity_spec>()}
is \tcode{true}, and
the specialization of \tcode{Q::named_quantity_spec} has a template argument equal to \tcode{subkind},
returns \tcode{q}.
\item
Otherwise, if \tcode{Q::\exposidnc{parent}} is a valid expression,
returns \tcode{\exposidnc{get-kind-tree-root}(Q::\exposidnc{parent})}.
\item
Otherwise, if \tcode{\exposconceptnc{DerivedQuantitySpec}<Q>} is \tcode{true},
returns
\begin{codeblock}
@\exposidnc{expr-map}@<@\exposidnc{to-kind}@, derived_quantity_spec, struct dimensionless>(q)
\end{codeblock}
where \exposid{to-kind} is defined as follows:
\begin{codeblock}
template<@\libconcept{quantity_spec}@ Q>
using @\exposidnc{to-kind}@ = decltype(@\exposidnc{get-kind-tree-root}@(Q{}));  // \expos
\end{codeblock}
\item
Otherwise, returns \tcode{q}.
\end{itemize}
\end{itemdescr}

\indexlibrarymember{get_kind}{quantity_spec}
\begin{itemdecl}
template<@\libconcept{quantity_spec}@ Q>
consteval @\libconcept{quantity_spec}@ auto get_kind(Q);
\end{itemdecl}

\begin{itemdescr}
\pnum
\returns
\tcode{\exposidnc{normalized-point-spec}(get_kind(Q::\exposidnc{quantity-spec}), \exposidnc{point-spec-origin}(Q\{\}))}\iref{point.qty}
if \tcode{Q} models \exposconceptnc{PointQuantitySpec},
and \tcode{kind_of<\exposidnc{get-kind-tree-root}(Q\{\})>} otherwise.
\end{itemdescr}

\rSec4[get.common.qty.spec]{Get common quantity specification}

\pnum
In \ref{get.common.qty.spec}, the term named quantity has the meaning given in \ref{named.qty}.
The terms factor, proper ancestor, kind root, declared kind, defining ancestor, depth, and conversion rank
have the meaning given in \ref{qty.spec.conv}, and
\tcode{q1} and \tcode{q2} denote objects of types modeling \libconcept{quantity_spec}.

\pnum
A quantity specification \tcode{Q} \defn{reaches} a quantity specification \tcode{C}
if \tcode{\exposid{conversion-rank}(Q, C)} is \tcode{implicit}.
\tcode{Q} and \tcode{C} are \defn{interconvertible} if each reaches the other.
\tcode{Q} is \defn{more specific} than \tcode{C}
if \tcode{Q} reaches \tcode{C} and \tcode{C} does not reach \tcode{Q}.

\pnum
A quantity specification is \defn{committed} if its type does not model
\exposconceptnc{QuantityKindSpec}.
A committed quantity specification is a \defn{power form}
if its factors consist of a single factor $(\tcode{N}, e)$ with $e \neq 1$.
\begin{note}
\tcode{pow<2>(isq::speed)} and \tcode{inverse(isq::period_duration)} are power forms;
\tcode{isq::speed} is a named quantity, and \tcode{isq::length / isq::time} is neither.
\end{note}

\pnum
The \defn{shared factors} of two quantity specifications \tcode{A} and \tcode{B} are
the factors $(\tcode{N}, c)$ such that
$(\tcode{N}, e)$ is a factor of \tcode{A}, $(\tcode{N}, f)$ is a factor of \tcode{B},
$e$ and $f$ have the same sign,
$c$ is the one of $e$ and $f$ with the smaller magnitude, or $e$ if they are equal, and $c$ is an integer.
The \defn{join} of \tcode{A} and \tcode{B}, neither of which models
\exposconceptnc{QuantityKindSpec} and at most one of which is a named quantity,
is determined as follows.
Let $S$ be their shared factors.
If both are named quantities, or $S$ is empty, there is no join.
Otherwise, let \tcode{R1} and \tcode{R2} be the quantity specifications whose factors are
the factors of \tcode{A} and of \tcode{B}, respectively,
with the exponent of every factor in $S$ subtracted from the exponent of the factor with the same named quantity.
If \tcode{get_common_quantity_spec(R1, R2)} is not a valid expression, there is no join.
Otherwise, the join is \tcode{get_common_quantity_spec(R1, R2) * P},
where \tcode{P} is the quantity specification whose factors are $S$.
\begin{note}
The join calls \tcode{get_common_quantity_spec} on quantity specifications with fewer factors in total than \tcode{A} and \tcode{B} together, so the recursion ends.
\end{note}
\begin{example}
The join of \tcode{special_speed * isq::width} and \tcode{special_speed * isq::height},
with \tcode{special_speed} a child of \tcode{isq::speed},
is \tcode{special_speed * isq::length}.
\end{example}

\pnum
The \defn{expansion} of a committed quantity specification is
the quantity specification obtained by replacing, repeatedly until none is left,
every factor $(\tcode{N}, e)$ whose named quantity has a defining ancestor \tcode{A}
and is not a declared kind by \tcode{A.\exposidnc{equation}}
(replacing a factor is specified in \ref{qty.spec.conv}).

\pnum
The \defn{fold} of two committed quantity specifications \tcode{A} and \tcode{B}
that are interconvertible and neither of which is a named quantity
is determined by the following procedure.
Let $F_A$ and $F_B$ be the lists of factors of \tcode{A} and of \tcode{B}.
The following is repeated until an iteration changes neither list:
for the first factor $(\tcode{M}, f)$ of $F_B$, in the order of $F_B$, such that
\tcode{M} is its own defining ancestor and not a declared kind,
the exponent $g$ of \tcode{M} in $F_A$ (zero if absent) is zero or has the sign of $f$ and a smaller magnitude,
and for every factor $(\tcode{N}, x)$ of \tcode{M.\exposidnc{equation}} the list $F_A$ holds
a factor $(\tcode{N}, y)$ such that $y$ has the sign of $x \cdot (f - g)$
and a magnitude of at least $|x \cdot (f - g)|$,
each such $y$ is reduced by $x \cdot (f - g)$ and
the factor $(\tcode{M}, f - g)$ is merged into $F_A$;
then the same is done with the roles of $F_A$ and $F_B$ exchanged.
Afterwards:
\begin{itemize}
\item
If $F_A$ and $F_B$ hold the same factors,
the fold is the quantity specification whose factors are $F_A$.
\item
Otherwise, let the depth sum of a quantity specification be
the sum of the depths of the named quantities of its factors.
If the depth sums of \tcode{A} and \tcode{B}, as given, differ, the fold is the one with the greater depth sum.
\item
Otherwise, if \tcode{A} and \tcode{B} have different numbers of factors,
the fold is the one with fewer factors.
\item
Otherwise, the fold is the expansion of \tcode{A}.
\end{itemize}
\begin{note}
Every fold replaces factors by one factor of greater depth, so the repetition ends.
\end{note}
\begin{example}
The fold of \tcode{isq::area * pow<3>(isq::length)} and \tcode{pow<2>(isq::length) * isq::volume}
is \tcode{isq::area * isq::volume}.
The fold of \tcode{isq::area * isq::speed} and \tcode{isq::volume / isq::time}
is \tcode{isq::area * isq::speed}.
\end{example}

\pnum
The \defn{representative} of a non-empty set $E$ of quantity specifications,
each reached by both \tcode{q1} and \tcode{q2}
and each interconvertible with a designated element \tcode{P} of $E$,
is determined as follows:
\begin{itemize}
\item
If $E$ contains committed quantity specifications, those that are not committed are discarded.
Otherwise, every element \tcode{kind_of<K>} of $E$ is replaced by \tcode{K},
and the quantity specification \tcode{C} determined by the following items,
applied to the resulting set in place of $E$, yields the representative \tcode{kind_of<C>}.
\item
If an element of $E$ is a named quantity, it is the representative.
\item
Otherwise, if exactly one element of $E$ is a power form, it is the representative.
\item
Otherwise, the representative is the fold of the elements of $E$, folded from the first to the last
in the order in which they were listed as candidates below.
\end{itemize}
A representative that is not an element of $E$ is the representative only if \tcode{q1} and \tcode{q2} reach it;
otherwise the first element of $E$ in the order of candidates is the representative.

\indexlibrarymember{get_common_quantity_spec}{quantity_spec}
\begin{itemdecl}
consteval @\libconcept{quantity_spec}@ auto get_common_quantity_spec(@\libconcept{quantity_spec}@ auto... qs)
  requires @\seebelownc@;
\end{itemdecl}

\begin{itemdescr}
\pnum
Let
\begin{itemize}
\item
\tcode{q1} be \tcode{qs...[0]},
\item
\tcode{q2} be \tcode{qs...[1]},
\item
\tcode{QQ1} be \tcode{\exposid{remove-kind}(q1)},
\item
\tcode{QQ2} be \tcode{\exposidnc{remove-kind}(q2)}, and
\item
\tcode{rest} be a pack denoting the elements of \tcode{qs} without \tcode{q1} and \tcode{q2}.
\end{itemize}

\pnum
The \defnx{candidates}{candidate} for \tcode{q1} and \tcode{q2} are, in this order and each once:
\begin{itemize}
\item
\tcode{q1} and \tcode{q2},
\item
the join of \tcode{QQ1} and \tcode{QQ2}, if any,
\item
\tcode{\exposid{get-kind-tree-root}(QQ1)},
\item
every proper ancestor of \tcode{QQ1}, if \tcode{QQ1} is a named quantity,
\item
\tcode{\exposidnc{get-kind-tree-root}(QQ2)}, and
\item
every proper ancestor of \tcode{QQ2}, if \tcode{QQ2} is a named quantity.
\end{itemize}
If both \tcode{decltype(q1)} and \tcode{decltype(q2)} model \exposconceptnc{QuantityKindSpec},
the candidates are instead \tcode{kind_of<C>} for every candidate \tcode{C} determined as above
for \tcode{QQ1} and \tcode{QQ2} such that \tcode{\exposidnc{get-kind-tree-root}(C) == C}.

\pnum
The \defn{common quantity specification} of \tcode{q1} and \tcode{q2}
is determined as follows.
If \tcode{decltype(q1)} and \tcode{decltype(q2)} are the same type, it is \tcode{q1}.
Otherwise, if \tcode{decltype(q1)} or \tcode{decltype(q2)} models \exposconceptnc{DeltaQuantitySpec},
it is the common quantity specification of
\tcode{\exposidnc{remove-delta}(q1)} and \tcode{\exposidnc{remove-delta}(q2)}.
Otherwise, if exactly one of \tcode{decltype(q1)} and \tcode{decltype(q2)}
models \exposconceptnc{PointQuantitySpec}, there is no common quantity specification.
Otherwise, if \tcode{\exposid{characters-overlap}(q1, q2)} is \tcode{false}\iref{qty.spec.concepts},
there is no common quantity specification.
Otherwise, if both model \exposconceptnc{PointQuantitySpec},
let \tcode{C} be the common quantity specification of
\tcode{q1.\exposidnc{quantity-spec}} and \tcode{q2.\exposidnc{quantity-spec}},
let \tcode{O1} and \tcode{O2} be \tcode{\exposid{point-spec-origin}(q1)} and \tcode{\exposid{point-spec-origin}(q2)}\iref{point.qty},
and let \tcode{O} be \tcode{O1} if \tcode{O1} and \tcode{O2} have the same type,
and \tcode{\exposid{get-absolute-point-origin}(O1)}
if \tcode{\exposid{same-absolute-point-origins}(O1, O2)} is \tcode{true}\iref{qty.same.abs.pt.origs} otherwise.
There is none if \tcode{C} or \tcode{O} does not exist.
Otherwise, let \tcode{P} be \tcode{\exposid{normalized-point-spec}(C, O)}
if \tcode{\libconcept{point_origin_for}<decltype(O), C>} is satisfied,
and \tcode{point<C>} if \tcode{O == natural_point_origin<C>} is \tcode{true} otherwise.
The common quantity specification is \tcode{P} if both \tcode{q1} and \tcode{q2} reach \tcode{P};
otherwise, and if \tcode{P} does not exist, there is none.
Otherwise, let $R$ be the set of candidates that both \tcode{q1} and \tcode{q2} reach,
with every element removed for which another element of $R$ is more specific.
If $R$ is empty, there is no common quantity specification.
Otherwise, let \tcode{P} be the first element of $R$, in the order of candidates,
of the highest of the following categories, listed from the highest:
a committed named quantity, a committed power form, any other committed quantity specification,
and, for a specification \tcode{kind_of<C>}, the category of \tcode{C} ranked below all three.
The common quantity specification is the representative of the set of
\tcode{P} and every element of $R$ interconvertible with \tcode{P}.
\begin{note}
Both \tcode{q1} and \tcode{q2} reach the common quantity specification.
A candidate not reached by one of the operands is never the result,
even when it is the kind root of that operand.
\end{note}
\begin{example}
\begin{codeblock}
static_assert(get_common_quantity_spec(isq::width, isq::height) == isq::length);
static_assert(get_common_quantity_spec(kind_of<isq::length>, isq::width) == isq::width);
static_assert(get_common_quantity_spec(isq::speed, isq::length / isq::time) == isq::speed);
static_assert(get_common_quantity_spec(isq::frequency, inverse(isq::time)) == inverse(isq::time));
static_assert(get_common_quantity_spec(isq::width / isq::period_duration, isq::height / isq::period_duration) ==
              isq::length / isq::period_duration);
static_assert(get_common_quantity_spec(isq::area * pow<3>(isq::length), pow<2>(isq::length) * isq::volume) ==
              isq::area * isq::volume);
static_assert(get_common_quantity_spec(point<isq::height>, point<isq::length>) == point<isq::length>);
static_assert(get_common_quantity_spec(delta<isq::height>, isq::width) == isq::length);
\end{codeblock}
\tcode{isq::angular_measure} and \tcode{dimensionless} have no common quantity specification,
because a declared kind reaches nothing above itself and \tcode{dimensionless} reaches nothing below itself.
A point and a delta have no common quantity specification either,
nor do \tcode{point<isq::height>} and \tcode{point<isq::width>},
whose natural point origins are not the same origin,
nor \tcode{isq::velocity} and \tcode{isq::speed}, whose characters do not overlap.
\end{example}

\pnum
\returns
\tcode{q1} if \tcode{sizeof...(qs) == 1};
otherwise, the common quantity specification of \tcode{q1} and \tcode{q2} if \tcode{sizeof...(qs) == 2};
otherwise, \tcode{get_common_quantity_spec(get_common_quantity_spec(q1, q2), rest...)},
which every element of \tcode{qs} reaches.

\pnum
\tcode{\exposid{has-common-quantity-spec}(q1, q2)} is \tcode{true}
if \tcode{q1} and \tcode{q2} have a common quantity specification, and \tcode{false} otherwise.
\tcode{\exposid{reached-by-all}(c, qs...)} is \tcode{true}
if every element of \tcode{qs} reaches \tcode{c}, and \tcode{false} otherwise.

\pnum
\remarks
The expression in the \fakegrammarterm{requires-clause} is equivalent to:
\begin{codeblock}
(sizeof...(qs) != 0 &&
 (sizeof...(qs) == 1 ||
  (sizeof...(qs) == 2 && @\exposid{has-common-quantity-spec}@(q1, q2)) ||
  (requires { get_common_quantity_spec(get_common_quantity_spec(q1, q2), rest...); } &&
   @\exposid{reached-by-all}@(get_common_quantity_spec(get_common_quantity_spec(q1, q2), rest...), qs...))))
\end{codeblock}
\end{itemdescr}

\rSec4[qty.is.child.of]{Is child of}

\indexlibrarymemberexpos{quantity_spec}{is-child-of}
\begin{itemdecl}
template<@\libconcept{quantity_spec}@ Child, @\libconcept{quantity_spec}@ Parent>
consteval bool @\exposidnc{is-child-of}@(Child ch, Parent p);  // \expos
\end{itemdecl}

\begin{itemdescr}
\pnum
\returns
If $h(\tcode{p})$ has more elements than $h(\tcode{ch})$, returns \tcode{false}.
Otherwise, let $C$ be a tuple of the last $s$ elements of $h(\tcode{ch})$,
where $s$ is the number of elements in $h(\tcode{p})$.
Returns \tcode{$C_0$ == p}.
\end{itemdescr}

\rSec2[qty.unit]{Unit}

\rSec3[qty.unit.general]{General}

\pnum
Subclause \ref{qty.unit} specifies the components
for defining a unit of measurement\irefiev{112-01-14}.

\rSec3[qty.unit.mag]{Magnitude}

\rSec4[qty.unit.mag.general]{General}

\pnum
Subclause \ref{qty.unit.mag} specifies the components
used to represent the numerical value\irefiev{112-01-29} of a unit
with support for powers\irefiev{102-02-08} of real numbers\irefiev{102-02-05}.

\rSec4[qty.unit.mag.concepts]{Concepts}

\begin{itemdecl}
template<typename T>
concept @\deflibconcept{MagConstant}@ = @\exposconceptnc{SymbolicConstant}@<T> && @\exposidnc{is-derived-from-specialization-of}@<T, mag_constant>();

template<typename T>
concept @\deflibconcept{unit_magnitude}@ = (@\exposidnc{is-specialization-of}@<T, @\exposidnc{unit-magnitude}@>());

template<typename T>
concept @\defexposconceptnc{MagArg}@ = std::@\stdconcept{integral}@<T> || @\libconcept{MagConstant}@<T>;  // \expos
\end{itemdecl}

\rSec4[qty.unit.mag.types]{Types}

\indexlibraryglobal{mag_constant}
\begin{codeblock}
namespace std {

template<symbol_text Symbol, long double Value>
  requires(Value > 0)
struct mag_constant {
  static constexpr auto @\exposidnc{symbol}@ = Symbol;       // \expos
  static constexpr long double @\exposidnc{value}@ = Value;  // \expos
};

}
\end{codeblock}

A specialization of \tcode{mag_constant} represents a real number\irefiev{102-02-05}.
\tcode{Symbol} is its symbol, and
\tcode{Value} is (an approximation of) its value.

\indexlibraryglobalexpos{unit-magnitude}
\begin{codeblock}
namespace std {

template<auto... Ms>
struct @\exposidnc{unit-magnitude}@ {  // \expos
  // \ref{qty.unit.mag.ops}, operations

  template<@\libconcept{unit_magnitude}@ M>
  friend consteval @\libconcept{unit_magnitude}@ auto operator*(@\exposidnc{unit-magnitude}@ lhs, M rhs);

  friend consteval auto operator/(@\exposidnc{unit-magnitude}@ lhs, @\libconcept{unit_magnitude}@ auto rhs);

  template<@\libconcept{unit_magnitude}@ Rhs>
  friend consteval bool operator==(@\exposidnc{unit-magnitude}@, Rhs);

  template<int Num, int Den = 1>
  friend consteval auto @\exposidnc{pow}@(@\exposidnc{unit-magnitude}@);                         // \expos

  // \ref{qty.unit.mag.utils}, utilities

  friend consteval bool @\exposidnc{is-positive-integral-power}@(@\exposidnc{unit-magnitude}@);  // \expos

  template<auto... Ms2>
  friend consteval auto @\exposidnc{common-magnitude}@(@\exposidnc{unit-magnitude}@,             // \expos
                                         @\exposidnc{unit-magnitude}@<Ms2...>);
};

}
\end{codeblock}

\pnum
A specialization of \exposid{unit-magnitude}
represents the product of its template arguments.

\pnum
For the purposes of specifying the implementation-defined limits,
let the representation of the terms of \exposid{unit-magnitude} be the structure
\begin{codeblock}
struct {
  @\exposidnc{ratio}@ exp;
  @\exposidnc{base-type}@ base;
};
\end{codeblock}
representing the number $\tcode{base}^\tcode{exp}$,
where \exposid{base-type} is a model of \exposconcept{MagArg}.
\begin{itemize}
\item
There is a single term for each \exposid{base-type}.
\item
\tcode{exp.num} is not expanded into base.
\begin{note}
$2^3 = 8$ is not permitted.
\end{note}
\item
\tcode{exp.den} can reduce the base.
\begin{note}
$4^{1/2} = 2$ is permitted.
\end{note}
\item
If the result of an operation on \tcode{std::intmax_t} values is undefined,
the behavior is
\impldef{behavior of \exposid{unit-magnitude} operations that do not fit in a \tcode{std::intmax_t}}.
\end{itemize}

\rSec4[qty.unit.mag.ops]{Operations}

\indexlibrarymember{operator*}{unit_magnitude}
\begin{itemdecl}
template<@\libconcept{unit_magnitude}@ M>
friend consteval @\libconcept{unit_magnitude}@ auto operator*(@\exposidnc{unit-magnitude}@ lhs, M rhs);
\end{itemdecl}

\begin{itemdescr}
\pnum
\returns
\begin{itemize}
\item
If \tcode{sizeof...(Ms) == 0} is \tcode{true}, returns \tcode{rhs}.
\item
Otherwise, if \tcode{std::is_same_v<M, \exposidnc{unit-magnitude}<>>}, returns \tcode{lhs}.
\item
Otherwise, returns an unspecified value equal to $\tcode{lhs} \times \tcode{rhs}$.
\end{itemize}
\end{itemdescr}

\indexlibrarymember{operator/}{unit_magnitude}
\begin{itemdecl}
friend consteval auto operator/(@\exposidnc{unit-magnitude}@ lhs, @\libconcept{unit_magnitude}@ auto rhs);
\end{itemdecl}

\begin{itemdescr}
\pnum
\returns
\tcode{lhs * \exposid{pow}<-1>(rhs)}.
\end{itemdescr}

\indexlibrarymember{operator==}{unit_magnitude}
\begin{itemdecl}
template<@\libconcept{unit_magnitude}@ Rhs>
friend consteval bool operator==(@\exposidnc{unit-magnitude}@, Rhs);
\end{itemdecl}

\begin{itemdescr}
\pnum
\returns
\tcode{std::is_same_v<\exposidnc{unit-magnitude}, Rhs>}.
\end{itemdescr}

\indexlibrarymemberexpos{unit_magnitude}{pow}
\begin{itemdecl}
template<int Num, int Den = 1>
friend consteval auto @\exposidnc{pow}@(@\exposidnc{unit-magnitude}@ base);  // \expos
\end{itemdecl}

\begin{itemdescr}
\pnum
\returns
\begin{itemize}
\item
If \tcode{Num == 0} is \tcode{true}, returns \tcode{\exposidnc{unit-magnitude}<>\{\}}.
\item
Otherwise, returns an unspecified value equal to $\tcode{base}^{\tcode{Num}/\tcode{Den}}$.
\end{itemize}
\end{itemdescr}

\indexlibraryglobal{mag}
\begin{itemdecl}
template<@\exposconceptnc{MagArg}@ auto V>
constexpr @\libconcept{unit_magnitude}@ auto mag = @\seebelownc@;
\end{itemdecl}

\begin{itemdescr}
\pnum
\constraints
\tcode{V} is greater than $0$.

\pnum
\effects
If \tcode{\libconcept{MagConstant}<decltype(V)>} is satisfied,
initializes \tcode{mag} with \tcode{\exposidnc{unit-magnitude}<V>\{\}}.
Otherwise, initializes \tcode{mag} with
an unspecified value equal to \tcode{V}.
\end{itemdescr}

\indexlibraryglobal{mag_ratio}
\begin{itemdecl}
template<std::intmax_t N, std::intmax_t D>
  requires(N > 0)
constexpr @\libconcept{unit_magnitude}@ auto mag_ratio = @\seebelownc@;
\end{itemdecl}

\begin{itemdescr}
\pnum
\effects
Initializes \tcode{mag_ratio} with
an unspecified value equal to $\tcode{N}/\tcode{D}$.
\end{itemdescr}

\indexlibraryglobal{mag_power}
\begin{itemdecl}
template<@\exposconceptnc{MagArg}@ auto Base, int Num, int Den = 1>
constexpr @\libconcept{unit_magnitude}@ auto mag_power = @\exposidnc{pow}@<Num, Den>(mag<Base>);
\end{itemdecl}

\begin{itemdescr}
\pnum
\constraints
\tcode{Base} is greater than $0$.
\end{itemdescr}

\rSec4[qty.unit.mag.utils]{Utilities}

\indexlibrarymemberexpos{unit_magnitude}{is-positive-integral-power}
\begin{itemdecl}
friend consteval bool @\exposidnc{is-positive-integral-power}@(@\exposidnc{unit-magnitude}@ x);              // \expos
\end{itemdecl}

\begin{itemdescr}
\pnum
\returns
\tcode{false} if \tcode{x} has a negative or rational exponent, and
\tcode{true} otherwise.
\end{itemdescr}

\indexlibrarymemberexpos{unit_magnitude}{common-magnitude}
\begin{itemdecl}
template<auto... Ms2>
friend consteval auto @\exposidnc{common-magnitude}@(@\exposidnc{unit-magnitude}@, @\exposidnc{unit-magnitude}@<Ms2...>);  // \expos
\end{itemdecl}

\begin{itemdescr}
\pnum
\returns
The largest magnitude \tcode{C}
such that each input magnitude is expressible
by only positive powers relative to \tcode{C}.
\end{itemdescr}

\rSec3[qty.unit.traits]{Traits}

\indexlibraryglobal{space_before_unit_symbol}
\begin{itemdecl}
template<@\libconcept{unit}@ auto U>
constexpr bool space_before_unit_symbol = true;
\end{itemdecl}

\begin{itemdescr}
\pnum
The formatting functions\iref{qty.unit.sym.fmt} use \tcode{space_before_unit_symbol}
to determine whether there is a space
between the numerical value and the unit symbol.

\pnum
\remarks
Pursuant to \refcpp{namespace.std}\iref{spec.ext},
users may specialize \tcode{space_before_unit_symbol}
for cv-unqualified program-defined types.
Such specializations shall be usable in constant expressions\irefcpp{expr.const}
and have type \tcode{const bool}.
\end{itemdescr}

\rSec3[qty.unit.concepts]{Concepts}

\begin{itemdecl}
template<typename T>
concept @\deflibconcept{unit}@ = @\exposconceptnc{SymbolicConstant}@<T> && std::@\stdconcept{derived_from}@<T, @\exposidnc{unit-interface}@>;

template<typename T>
concept @\deflibconcept{prefixable_unit}@ = @\libconcept{unit}@<T> && @\exposidnc{is-derived-from-specialization-of}@<T, named_unit>();

template<typename U, auto QS>
concept @\deflibconcept{unit_of}@ =
  @\libconcept{unit}@<U> && @\libconcept{quantity_spec}@<decltype(QS)> &&
  @\exposconceptnc{QuantitySpecConvertibleTo}@<get_quantity_spec(U{}), @\exposidnc{remove-point}@(QS)> &&
  (get_kind(@\exposidnc{remove-point}@(QS)) == get_kind(get_quantity_spec(U{})) ||
   !@\exposconceptnc{NestedQuantityKindSpecOf}@<get_quantity_spec(U{}), @\exposidnc{remove-point}@(QS)>);

template<auto From, auto To>
concept @\defexposconceptnc{UnitConvertibleTo}@ =                                     // \expos
  @\libconcept{unit}@<decltype(From)> && @\libconcept{unit}@<decltype(To)> && (convertible(From, To));

template<typename U, auto FromU, auto QS>
concept @\defexposconceptnc{UnitCompatibleWith}@ =                                    // \expos
  @\libconcept{unit}@<U> && @\libconcept{unit}@<decltype(FromU)> && @\libconcept{quantity_spec}@<decltype(QS)> &&
  @\libconcept{unit_of}@<U, QS> && @\exposconceptnc{UnitConvertibleTo}@<FromU, U{}>;

template<typename T>
concept @\defexposconceptnc{OffsetUnit}@ = @\libconcept{unit}@<T> && requires { T::@\exposidnc{point-origin}@; };  // \expos

template<typename From, typename To>
concept @\defexposconceptnc{PotentiallyConvertibleTo}@ =                              // \expos
  @\libconcept{unit}@<From> && @\libconcept{unit}@<To> &&
  (@\exposid{implicitly-convertible}@(get_quantity_spec(From{}), get_quantity_spec(To{})));
\end{itemdecl}

\rSec3[qty.unit.types]{Types}

\rSec4[qty.canon.unit]{Canonical}

\indexlibraryglobalexpos{canonical-unit}
\begin{codeblock}
namespace std {

template<@\libconcept{unit_magnitude}@ M, @\libconcept{unit}@ U>
struct @\exposidnc{canonical-unit}@ {  // \expos
  M mag;
  U reference_unit;
};

}
\end{codeblock}

\pnum
\exposid{canonical-unit} represents a unit expressed in terms of base units\irefiev{112-01-18}.
\begin{note}
Other types representing units are equal only if they have the same type.
\exposid{canonical-unit} is used to implement binary relations other than equality.
\end{note}
\tcode{reference_unit} is simplified\iref{qty.sym.expr.algos}.

\indexlibrarymemberexpos{unit}{get-canonical-unit}
\begin{itemdecl}
consteval auto @\exposidnc{get-canonical-unit}@(@\libconcept{unit}@ auto u);  // \expos
\end{itemdecl}

\begin{itemdescr}
\pnum
\returns
The instantiation of \exposid{canonical-unit} for \tcode{u}.
\end{itemdescr}

\rSec4[qty.scaled.unit]{Scaled}

\indexlibraryglobal{scaled_unit}
\begin{codeblock}
namespace std {

template<@\libconcept{unit_magnitude}@ auto M, @\libconcept{unit}@ U>
  requires(M != @\exposidnc{unit-magnitude}@<>{} && M != mag<1>)
struct scaled_unit final : @\exposidnc{unit-interface}@ {
  using @\exposidnc{base-type}@ = scaled_unit;                         // \expos
  static constexpr @\libconcept{unit_magnitude}@ auto @\exposidnc{mag}@ = M;           // \expos
  static constexpr U @\exposidnc{reference-unit}@{};                   // \expos
  static constexpr auto @\exposidnc{point-origin}@ = U::point_origin;  // \expos, present only
    // if the \fakegrammarterm{qualified-id} \tcode{U::point_origin} is valid and denotes an object
};

}
\end{codeblock}

\pnum
\tcode{scaled_unit<M, U>} is used by the library
to represent the unit $\tcode{M} \times \tcode{U}$.

\rSec4[qty.named.unit]{Named}

\indexlibraryglobal{named_unit}
\begin{codeblock}
namespace std {

template<symbol_text Symbol, @\exposconceptnc{QuantityKindSpec}@ auto QS>
  requires(!Symbol.empty()) &&
          (@\exposconceptnc{BaseDimension}@<decltype(get_dimension(QS))> || @\exposidnc{remove-kind}@(QS) == dimensionless)
struct named_unit<Symbol, QS> : @\exposidnc{unit-interface}@ {
  using @\exposidnc{base-type}@ = named_unit;              // \expos
  static constexpr auto @\exposidnc{symbol}@ = Symbol;     // \expos
  static constexpr auto @\exposidnc{quantity-spec}@ = QS;  // \expos
};

template<symbol_text Symbol, @\exposconceptnc{QuantityKindSpec}@ auto QS, @\exposconceptnc{PointOrigin}@ auto PO>
  requires(!Symbol.empty()) &&
          (@\exposconceptnc{BaseDimension}@<decltype(get_dimension(QS))> || @\exposidnc{remove-kind}@(QS) == dimensionless)
struct named_unit<Symbol, QS, PO> : @\exposidnc{unit-interface}@ {
  using @\exposidnc{base-type}@ = named_unit;              // \expos
  static constexpr auto @\exposidnc{symbol}@ = Symbol;     // \expos
  static constexpr auto @\exposidnc{quantity-spec}@ = QS;  // \expos
  static constexpr auto @\exposidnc{point-origin}@ = PO;   // \expos
};

template<symbol_text Symbol, @\libconcept{unit}@ auto U>
  requires(!Symbol.empty()) && @\exposidnc{is-positive-canonical-unit-mag}@<U>
struct named_unit<Symbol, U> : decltype(U)::@\exposidnc{base-type}@ {
  using @\exposidnc{base-type}@ = named_unit;              // \expos
  static constexpr auto @\exposidnc{symbol}@ = Symbol;     // \expos
};

template<symbol_text Symbol, @\libconcept{unit}@ auto U, @\exposconceptnc{PointOrigin}@ auto PO>
  requires(!Symbol.empty()) && @\exposidnc{is-positive-canonical-unit-mag}@<U>
struct named_unit<Symbol, U, PO> : decltype(U)::@\exposidnc{base-type}@ {
  using @\exposidnc{base-type}@ = named_unit;              // \expos
  static constexpr auto @\exposidnc{symbol}@ = Symbol;     // \expos
  static constexpr auto @\exposidnc{point-origin}@ = PO;   // \expos
};

template<symbol_text Symbol, @\libconcept{unit}@ auto U, @\exposconceptnc{QuantityKindSpec}@ auto QS>
  requires(!Symbol.empty()) &&
          (get_dimension(QS) == get_dimension(@\exposidnc{get-associated-quantity}@(U))) &&
          @\exposidnc{is-positive-canonical-unit-mag}@<U>
struct named_unit<Symbol, U, QS> : decltype(U)::@\exposidnc{base-type}@ {
  using @\exposidnc{base-type}@ = named_unit;              // \expos
  static constexpr auto @\exposidnc{symbol}@ = Symbol;     // \expos
  static constexpr auto @\exposidnc{quantity-spec}@ = QS;  // \expos
};

template<symbol_text Symbol, @\libconcept{unit}@ auto U, @\exposconceptnc{QuantityKindSpec}@ auto QS,
         @\exposconceptnc{PointOrigin}@ auto PO>
  requires(!Symbol.empty()) &&
          (get_dimension(QS) == get_dimension(@\exposidnc{get-associated-quantity}@(U))) &&
          @\exposidnc{is-positive-canonical-unit-mag}@<U>
struct named_unit<Symbol, U, QS, PO> : decltype(U)::@\exposidnc{base-type}@ {
  using @\exposidnc{base-type}@ = named_unit;              // \expos
  static constexpr auto @\exposidnc{symbol}@ = Symbol;     // \expos
  static constexpr auto @\exposidnc{quantity-spec}@ = QS;  // \expos
  static constexpr auto @\exposidnc{point-origin}@ = PO;   // \expos
};

}
\end{codeblock}

\pnum
A \defnadj{named}{unit} is a type that models \libconcept{prefixable_unit}.
A specialization of \tcode{named_unit} is used as a base type when defining a named unit.

\pnum
In the following descriptions, let \tcode{U} be a named unit defined with an alluded signature.
The identifier of \tcode{U} represents
its unit name\irefiev{112-01-15}
or special unit name\irefiev{112-01-16}.
\tcode{Symbol} is its unit symbol\irefiev{112-01-17}.

\pnum
The possible arguments to \tcode{named_unit} are
\begin{itemize}
\item
$(\text{the unit symbol}, \text{a kind of base quantity}, \opt{\text{a point origin}})$,
\item
$(\text{the unit symbol}, \text{a unit expression}, \opt{\text{a point origin}})$, and
\item
$(\text{the unit symbol}, \text{a unit expression}, \text{a kind of quantity}, \opt{\text{a point origin}})$.
\end{itemize}

\pnum
The first signature defines the unit of a base quantity
without a unit prefix\irefiev{112-01-26}.
The second and third signatures with a unit expression argument \placeholder{E}
define \tcode{U} as implicitly convertible from \placeholder{E}.
The first and third signatures with a kind of quantity\irefiev{112-01-04} \placeholder{Q}
also restrict \tcode{U} to \placeholder{Q}.

\pnum
A point origin argument specifies the default point origin of \tcode{U}\iref{qty.def.pt.orig}.

\pnum
\begin{example}
\begin{codeblock}
// The first signature defines a base unit restricted to a kind of base quantity.
inline constexpr struct second final : named_unit<"s", kind_of<time>> {
} second;

// The second and third signatures give a name to the unit argument.
inline constexpr struct minute final : named_unit<"min", mag<60> * second> {
} minute;  // $\txtrm{min} = 60 \txtrm{ s}$.

// The third signature also further restricts the kind of quantity.
inline constexpr struct hertz final : named_unit<"Hz", inverse(second), kind_of<frequency>> {
} hertz;  // $\txtrm{Hz}$ can't measure becquerel, activity,
          // or any other quantity with dimension $\txtrm{T}^{-1}$
          // that isn't a kind of frequency.
\end{codeblock}
\end{example}

\rSec4[qty.prefixed.unit]{Prefixed}

\indexlibraryglobal{prefixed_unit}
\begin{codeblock}
namespace std {

template<symbol_text Symbol, @\libconcept{unit_magnitude}@ auto M, @\libconcept{prefixable_unit}@ auto U>
  requires(!Symbol.empty())
struct prefixed_unit : decltype(M * U)::@\exposidnc{base-type}@ {
  using @\exposidnc{base-type}@ = prefixed_unit;                   // \expos
  static constexpr auto @\exposidnc{symbol}@ = Symbol + U.@\exposidnc{symbol}@;  // \expos
};

}
\end{codeblock}

\pnum
\tcode{prefixed_unit<Symbol, M, U>} represents the unit \tcode{U} with a unit prefix\irefiev{112-01-26}.
\tcode{Symbol} is the symbol of the unit prefix.
\tcode{M} is the factor of the unit prefix.
A specialization of \tcode{prefixed_unit} is used as a base type when defining a unit prefix.
\begin{example}
\begin{codeblock}
template<@\libconcept{prefixable_unit}@ auto U>
struct kilo_ : prefixed_unit<"k", mag_power<10, 3>, U> {};

template<@\libconcept{prefixable_unit}@ auto U>
constexpr kilo_<U> kilo;

inline constexpr auto kilogram = kilo<si::gram>;
\end{codeblock}
\end{example}

\rSec4[qty.common.unit]{Common}

\indexlibraryglobal{common_unit}
\begin{codeblock}
namespace std {

template<@\libconcept{unit}@ U1, @\libconcept{unit}@ U2, @\libconcept{unit}@... Rest>
struct common_unit final : decltype(@\exposidnc{get-common-scaled-unit}@(U1{}, U2{}, Rest{}...))::@\exposidnc{base-type}@
{
  using @\exposidnc{base-type}@ = common_unit;       // \expos
  static constexpr auto @\exposidnc{common-unit}@ =  // \expos
    @\exposidnc{get-common-scaled-unit}@(U1{}, U2{}, Rest{}...);
};

}
\end{codeblock}

\pnum
\tcode{common_unit} is used by the library
to encapsulate a conversion factor between units\irefiev{112-01-33}
common to the operands of quantity addition.
\begin{example}
The result of \tcode{1 * km + 1 * mi}
has a common unit $[8/125] \txtrm{ m}$
encapsulated by \tcode{common_unit<mi, km>}.
\end{example}
A program that instantiates a specialization of \tcode{common_unit}
that is not a possible result of the library specifications
is ill-formed, no diagnostic required.

\begin{itemdecl}
template<@\libconcept{unit}@ U1, @\libconcept{unit}@ U2, @\libconcept{unit}@... Rest>
consteval @\libconcept{unit}@ auto @\exposidnc{get-common-scaled-unit}@(U1, U2, Rest... rest)  // \expos
  requires @\seebelownc@;
\end{itemdecl}

\begin{itemdescr}
\pnum
\effects
Equivalent to:
\begin{codeblock}
constexpr auto res = [] {
  constexpr auto canonical_lhs = @\exposidnc{get-canonical-unit}@(U1{});
  constexpr auto canonical_rhs = @\exposidnc{get-canonical-unit}@(U2{});
  constexpr auto common_mag = @\exposidnc{common-magnitude}@(canonical_lhs.mag, canonical_rhs.mag);
  if constexpr (common_mag == mag<1>)
    return canonical_lhs.reference_unit;
  else
    return scaled_unit<common_mag, decltype(auto(canonical_lhs.reference_unit))>{};
}();
if constexpr (sizeof...(rest) == 0)
  return res;
else
  return @\exposidnc{get-common-scaled-unit}@(res, rest...);
\end{codeblock}

\pnum
\remarks
The expression in the \fakegrammarterm{requires-clause} is equivalent to:
\begin{codeblock}
(convertible(U1{}, U2{}) && (sizeof...(Rest) == 0 || requires {
   @\exposidnc{get-common-scaled-unit}@(@\exposidnc{get-common-scaled-unit}@(u1, u2), rest...);
 }))
\end{codeblock}
\end{itemdescr}

\rSec4[qty.derived.unit]{Derived}

\indexlibraryglobal{derived_unit}
\begin{codeblock}
namespace std {

template<typename... Expr>
struct @\exposidnc{derived-unit-impl}@ :              // \expos
    @\exposidnc{unit-interface}@,
    @\exposidnc{expr-fractions}@<struct one, Expr...> {
  using @\exposidnc{base-type}@ = @\exposidnc{derived-unit-impl}@;  // \expos
};

template<@\exposconceptnc{SymbolicConstant}@... Expr>
struct derived_unit final : @\exposidnc{derived-unit-impl}@<Expr...> {};

}
\end{codeblock}

\pnum
\tcode{derived_unit} is used by the library
to represent a derived unit\irefiev{112-01-19}.
\begin{example}
\begin{codeblock}
using namespace si::unit_symbols;
int x = m * m;  // error: cannot construct from \tcode{derived_unit<\exposidnc{power}<si::metre, 2>>}
int y = m * s;  // error: cannot construct from \tcode{derived_unit<si::metre, si::second>}
int z = m / s;  // error: cannot construct from \tcode{derived_unit<si::metre, \exposidnc{per}<si::second>>}
\end{codeblock}
\end{example}
A program that instantiates a specialization of \tcode{derived_unit}
that is not a possible result of the library specifications
is ill-formed, no diagnostic required.

\rSec4[qty.unit.one]{One}

\indexlibraryglobal{one}
\begin{codeblock}
namespace std {

struct one final : @\exposidnc{derived-unit-impl}@<> {};

}
\end{codeblock}

\pnum
\tcode{one} represents the base unit\irefiev{112-01-18} of a quantity of dimension one\irefiev{112-01-13}.

\rSec3[qty.unit.ops]{Operations}

\indexlibraryglobalexpos{unit-interface}
\begin{codeblock}
namespace std {

struct @\exposidnc{unit-interface}@ {  // \expos
  template<@\libconcept{unit_magnitude}@ M, @\libconcept{unit}@ U>
  friend consteval @\libconcept{unit}@ auto operator*(M, U u);

  friend consteval @\libconcept{unit}@ auto operator*(@\libconcept{unit}@ auto, @\libconcept{unit_magnitude}@ auto) = delete;

  template<@\libconcept{unit_magnitude}@ M, @\libconcept{unit}@ U>
  friend consteval @\libconcept{unit}@ auto operator/(M mag, U u);

  template<@\libconcept{unit}@ Lhs, @\libconcept{unit}@ Rhs>
  friend consteval @\libconcept{unit}@ auto operator*(Lhs lhs, Rhs rhs);

  template<@\libconcept{unit}@ Lhs, @\libconcept{unit}@ Rhs>
  friend consteval @\libconcept{unit}@ auto operator/(Lhs lhs, Rhs rhs);

  // \ref{qty.unit.cmp}, comparison

  template<@\libconcept{unit}@ Lhs, @\libconcept{unit}@ Rhs>
  friend consteval bool operator==(Lhs, Rhs);

  template<@\libconcept{unit}@ Lhs, @\libconcept{unit}@ Rhs>
  friend consteval bool @\exposid{equivalent}@(Lhs lhs, Rhs rhs);  // \expos
};

}
\end{codeblock}

\indexlibrarymember{operator*}{unit_magnitude}
\indexlibrarymember{operator*}{unit}
\begin{itemdecl}
template<@\libconcept{unit_magnitude}@ M, @\libconcept{unit}@ U>
friend consteval @\libconcept{unit}@ auto operator*(M, U u);
\end{itemdecl}

\begin{itemdescr}
\pnum
\effects
Equivalent to:
\begin{codeblock}
if constexpr (std::is_same_v<M, decltype(auto(mag<1>))>)
  return u;
else if constexpr (@\exposidnc{is-specialization-of}@<U, scaled_unit>()) {
  if constexpr (M{} * U::@\exposidnc{mag}@ == mag<1>)
    return U::@\exposidnc{reference-unit}@;
  else
    return scaled_unit<M{} * U::@\exposidnc{mag}@, decltype(auto(U::@\exposidnc{reference-unit}@))>{};
} else
  return scaled_unit<M{}, U>{};
\end{codeblock}
\end{itemdescr}

\begin{itemdecl}
friend consteval @\libconcept{unit}@ auto operator*(@\libconcept{unit}@ auto, @\libconcept{unit_magnitude}@ auto) = delete;
\end{itemdecl}

\begin{itemdescr}
\pnum
\recommended
Suggest swapping the operands.
\end{itemdescr}

\indexlibrarymember{operator/}{unit_magnitude}
\indexlibrarymember{operator/}{unit}
\begin{itemdecl}
template<@\libconcept{unit_magnitude}@ M, @\libconcept{unit}@ U>
friend consteval @\libconcept{unit}@ auto operator/(M mag, U u);
\end{itemdecl}

\begin{itemdescr}
\pnum
\returns
\tcode{mag * inverse(u)}.
\end{itemdescr}

\indexlibrarymember{operator*}{unit}
\begin{itemdecl}
template<@\libconcept{unit}@ Lhs, @\libconcept{unit}@ Rhs>
friend consteval @\libconcept{unit}@ auto operator*(Lhs lhs, Rhs rhs);
\end{itemdecl}

\begin{itemdescr}
\pnum
\returns
\tcode{\exposidnc{expr-multiply}<derived_unit, struct one>(lhs, rhs)}.
\end{itemdescr}

\indexlibrarymember{operator/}{unit}
\begin{itemdecl}
template<@\libconcept{unit}@ Lhs, @\libconcept{unit}@ Rhs>
friend consteval @\libconcept{unit}@ auto operator/(Lhs lhs, Rhs rhs);
\end{itemdecl}

\begin{itemdescr}
\pnum
\returns
\tcode{\exposidnc{expr-divide}<derived_unit, struct one>(lhs, rhs)}.
\end{itemdescr}

\indexlibrarymember{inverse}{unit}
\begin{itemdecl}
consteval @\libconcept{unit}@ auto inverse(@\libconcept{unit}@ auto u);
\end{itemdecl}

\begin{itemdescr}
\pnum
\returns
\tcode{one / u}.
\end{itemdescr}

\indexlibrarymember{pow}{unit}
\begin{itemdecl}
template<std::intmax_t Num, std::intmax_t Den = 1, @\libconcept{unit}@ U>
  requires(Den != 0)
consteval @\libconcept{unit}@ auto pow(U u);
\end{itemdecl}

\begin{itemdescr}
\pnum
\returns
\tcode{\exposidnc{expr-pow}<Num, Den, derived_unit, struct one>(u)}.
\end{itemdescr}

\indexlibrarymember{sqrt}{unit}
\begin{itemdecl}
consteval @\libconcept{unit}@ auto sqrt(@\libconcept{unit}@ auto u);
\end{itemdecl}

\begin{itemdescr}
\pnum
\returns
\tcode{pow<1, 2>(u)}.
\end{itemdescr}

\indexlibrarymember{cbrt}{unit}
\begin{itemdecl}
consteval @\libconcept{unit}@ auto cbrt(@\libconcept{unit}@ auto u);
\end{itemdecl}

\begin{itemdescr}
\pnum
\returns
\tcode{pow<1, 3>(u)}.
\end{itemdescr}

\indexlibrarymember{square}{unit}
\begin{itemdecl}
consteval @\libconcept{unit}@ auto square(@\libconcept{unit}@ auto u);
\end{itemdecl}

\begin{itemdescr}
\pnum
\returns
\tcode{pow<2>(u)}.
\end{itemdescr}

\indexlibrarymember{cubic}{unit}
\begin{itemdecl}
consteval @\libconcept{unit}@ auto cubic(@\libconcept{unit}@ auto u);
\end{itemdecl}

\begin{itemdescr}
\pnum
\returns
\tcode{pow<3>(u)}.
\end{itemdescr}

\rSec3[qty.unit.cmp]{Comparison}

\indexlibrarymember{operator==}{unit}
\begin{itemdecl}
template<@\libconcept{unit}@ Lhs, @\libconcept{unit}@ Rhs>
friend consteval bool operator==(Lhs, Rhs);
\end{itemdecl}

\begin{itemdescr}
\pnum
\returns
\tcode{std::is_same_v<Lhs, Rhs>}.
\end{itemdescr}

\indexlibrarymemberexpos{unit}{equivalent}
\begin{itemdecl}
template<@\libconcept{unit}@ Lhs, @\libconcept{unit}@ Rhs>
friend consteval bool @\exposid{equivalent}@(Lhs lhs, Rhs rhs);
\end{itemdecl}

\begin{itemdescr}
\pnum
\effects
Equivalent to:
\begin{codeblock}
const auto lhs_canonical = @\exposidnc{get-canonical-unit}@(lhs);
const auto rhs_canonical = @\exposidnc{get-canonical-unit}@(rhs);
return lhs_canonical.mag == rhs_canonical.mag &&
       lhs_canonical.reference_unit == rhs_canonical.reference_unit;
\end{codeblock}
\end{itemdescr}

\indexlibrarymember{convertible}{unit}
\begin{itemdecl}
template<@\libconcept{unit}@ From, @\libconcept{unit}@ To>
consteval bool convertible(From from, To to);
\end{itemdecl}

\begin{itemdescr}
\pnum
\effects
Equivalent to:
\begin{codeblock}
if constexpr (std::is_same_v<From, To>)
  return true;
else if constexpr (@\exposconceptnc{PotentiallyConvertibleTo}@<From, To>)
  return std::is_same_v<decltype(@\exposidnc{get-canonical-unit}@(from).reference_unit),
                        decltype(@\exposidnc{get-canonical-unit}@(to).reference_unit)>;
else
  return false;
\end{codeblock}
\end{itemdescr}

\rSec3[qty.unit.obs]{Observers}

\indexlibrarymember{get_quantity_spec}{unit}
\begin{itemdecl}
consteval @\libconcept{quantity_spec}@ auto get_quantity_spec(@\libconcept{unit}@ auto u);
\end{itemdecl}

\begin{itemdescr}
\pnum
\returns
\tcode{kind_of<\exposidnc{get-associated-quantity}(u)>}.
\end{itemdescr}

\indexlibrarymember{get_unit}{unit}
\begin{itemdecl}
consteval @\libconcept{unit}@ auto get_unit(@\libconcept{unit}@ auto u);
\end{itemdecl}

\begin{itemdescr}
\pnum
\returns
\tcode{u}.
\end{itemdescr}

\indexlibrarymemberexpos{unit}{get-common-unit}
\begin{itemdecl}
consteval @\libconcept{unit}@ auto @\exposidnc{get-common-unit}@(@\libconcept{unit}@ auto... us)  // \expos
  requires @\seebelownc@;
\end{itemdecl}

\begin{itemdescr}
\pnum
Let
\begin{itemize}
\item
\tcode{u1} be \tcode{us...[0]},
\item
\tcode{u2} be \tcode{us...[1]}, and
\item
\tcode{rest} be a pack denoting the elements of \tcode{us} without \tcode{u1} and \tcode{u2}.
\end{itemize}

\pnum
The \defn{members} of a unit \tcode{u} are
\tcode{Us{}...} if \tcode{decltype(u)} is \tcode{common_unit<Us...>}, and \tcode{u} otherwise.

\pnum
Two units \tcode{a} and \tcode{b} are \defnx{compatible}{compatible, units} if
\tcode{convertible(a, b)} is \tcode{true} and either
\tcode{\exposid{explicitly-convertible}(get_quantity_spec(a), get_quantity_spec(b))} or
\tcode{\exposid{explicitly-convertible}(get_quantity_spec(b), get_quantity_spec(a))} is \tcode{true}.

\pnum
For a unit \tcode{u}, let $m(\tcode{u})$ be \tcode{\exposid{get-canonical-unit}(u).mag}.
A unit \tcode{a} \defn{divides} a unit \tcode{b} if
\tcode{\exposid{common-magnitude}($m(\tcode{a})$, $m(\tcode{b})$) == $m(\tcode{a})$};
\tcode{a} \defn{strictly divides} \tcode{b} if \tcode{a} divides \tcode{b}
and \tcode{\exposid{equivalent}(a, b)} is \tcode{false}.
\begin{note}
For magnitudes that are rational numbers, \tcode{a} divides \tcode{b} when
\tcode{b} is an integer multiple of \tcode{a}.
\tcode{si::radian} divides \tcode{mag<2> * pi * si::radian}, a revolution,
because the exponent of $\pi$ is compared like any other.
\end{note}

\pnum
A unit is \defnx{anonymous}{anonymous unit} if its type is a specialization of
\tcode{scaled_unit} or of \tcode{derived_unit}.
Of two equivalent units \tcode{a} and \tcode{b}, \tcode{a} is \defnx{preferred}{preferred, unit} if
\begin{itemize}
\item
\tcode{decltype(a)} is derived from \tcode{decltype(b)::\exposid{base-type}} and
\tcode{decltype(b)} is not derived from \tcode{decltype(a)::\exposidnc{base-type}}, or
\item
neither is derived from the other's \exposidnc{base-type}, \tcode{b} is anonymous, and \tcode{a} is not, or
\item
both are anonymous and \tcode{a} has fewer unit factors than \tcode{b}, where
the unit factors of a specialization of \tcode{derived_unit} are the elements of its
\exposid{num} and \exposidnc{den}, those of a specialization of \tcode{scaled_unit} are
those of its reference unit, and any other unit is one unit factor.
\end{itemize}
\begin{example}
\tcode{si::hertz} is preferred to \tcode{one / si::second},
\tcode{si::kilogram} to \tcode{mag<1000> * si::gram}, and
\tcode{si::joule} to \tcode{si::kilogram * square(si::metre) / square(si::second)}, and
\tcode{si::newton * si::metre} to the latter as well.
Neither \tcode{iec::byte} nor \tcode{iec::octet} is preferred to the other,
nor \tcode{si::newton * si::metre} to \tcode{si::watt * si::second}.
\end{example}

\pnum
The \defn{common unit} of \tcode{u1} and \tcode{u2} is determined as follows.
Let $M$ be the set of the members of \tcode{u1} and of \tcode{u2},
without every unit that another unit of the set strictly divides, and,
of two equivalent units of which one is preferred, without the other, and,
of two units whose magnitudes have the same absolute value and opposite signs,
without the one whose magnitude is negative.
Two units \tcode{a} and \tcode{b} are ordered by $m(\tcode{a})$ and $m(\tcode{b})$;
two equivalent units by their symbols, compared as \tcode{symbol_text} compares,
when both have a symbol and the symbols differ;
otherwise by their definitions:
two units each defined as a scaled unit, named or not
(\tcode{imperial::link} and \tcode{usc::link} share the symbol \tcode{"li"}),
by their magnitudes and then their reference units,
two specializations of \tcode{derived_unit} by their unit factors in turn,
each by the base unit in this order and then by its exponent,
and a unit defined as a scaled unit against one that is not, by its magnitude against \tcode{mag<1>}.
$M$ is ordered by this order.
\begin{itemize}
\item
If $M$ has one element, the common unit is that element.
\item
Otherwise, if every element of $M$ is a specialization of \tcode{scaled_unit},
the common unit is \tcode{\exposid{get-common-scaled-unit}(M...)}.
\item
Otherwise, the common unit is \tcode{common_unit<M...>{}}.
\end{itemize}
\begin{note}
The common unit depends on the set of members only.
Every member of a \tcode{common_unit} is an integer multiple of its \exposid{common-unit}
when the magnitudes are rational,
and no member is strictly divided by another.
\end{note}
\begin{example}
\begin{codeblock}
static_assert(@\exposid{get-common-unit}@(si::kilo<si::metre>, si::milli<si::metre>) == si::milli<si::metre>);
static_assert(@\exposid{get-common-unit}@(si::kilo<si::metre>, si::milli<si::metre>, yard_pound::mile) ==
              si::milli<si::metre>);
static_assert(is_same_v<decltype(@\exposid{get-common-unit}@(si::kilo<si::metre>, yard_pound::mile)),
                        common_unit<si::kilo_<struct si::metre>, struct yard_pound::mile>>);
static_assert(is_same_v<decltype(@\exposid{get-common-unit}@(yard_pound::mile, si::kilo<si::metre>)),
                        common_unit<si::kilo_<struct si::metre>, struct yard_pound::mile>>);
static_assert(@\exposid{get-common-unit}@(mag_ratio<1, 20> * si::kelvin, mag_ratio<1, 9> * si::kelvin) ==
              mag_ratio<1, 180> * si::kelvin);
\end{codeblock}
\tcode{si::hertz} and \tcode{si::becquerel} are not compatible,
because \tcode{isq::frequency} and \tcode{isq::activity} are not convertible in either direction.
\end{example}

\pnum
\returns
\tcode{u1} if \tcode{sizeof...(us) == 1};
otherwise, the common unit of \tcode{u1} and \tcode{u2} if \tcode{sizeof...(us) == 2};
otherwise, \tcode{\exposid{get-common-unit}(\exposid{get-common-unit}(u1, u2), rest...)},
with which every element of \tcode{us} is compatible.

\pnum
\tcode{\exposid{compatible-units}(u1, u2)} is \tcode{true}
if every member of \tcode{u1} is compatible with every member of \tcode{u2}, and \tcode{false} otherwise.
\tcode{\exposid{compatible-with-all}(c, us...)} is \tcode{true}
if every element of \tcode{us} is compatible with \tcode{c}, and \tcode{false} otherwise.

\pnum
\remarks
The expression in the \fakegrammarterm{requires-clause} is equivalent to:
\begin{codeblock}
(sizeof...(us) != 0 &&
 (sizeof...(us) == 1 ||
  (sizeof...(us) == 2 && @\exposid{compatible-units}@(u1, u2)) ||
  (requires { @\exposid{get-common-unit}@(@\exposid{get-common-unit}@(u1, u2), rest...); } &&
   @\exposid{compatible-with-all}@(@\exposid{get-common-unit}@(@\exposid{get-common-unit}@(u1, u2), rest...), us...))))
\end{codeblock}
\end{itemdescr}

\rSec3[assoc.qty]{Associated quantity}

\indexlibrarymemberexpos{unit}{get-associated-quantity}
\begin{itemdecl}
template<@\libconcept{unit}@ U>
consteval auto @\exposidnc{get-associated-quantity}@(U u);  // \expos
\end{itemdecl}

\begin{itemdescr}
\pnum
\returns
\begin{itemize}
\item
If \tcode{U} is of the form \tcode{common_unit<Us...>},
returns
\begin{codeblock}
@\exposidnc{remove-kind}@(get_common_quantity_spec(kind_of<@\exposidnc{get-associated-quantity}@(Us{})>...))
\end{codeblock}
\begin{note}
The members' kinds take part, so that a common unit of an angle unit and a unit of a pure number
is a unit of the angle: \tcode{kind_of<dimensionless>} converts implicitly to
\tcode{kind_of<isq::angular_measure>}, while \tcode{dimensionless} does not convert implicitly to
\tcode{isq::angular_measure}.
\end{note}
\item
Otherwise, if \tcode{U::\exposidnc{quantity-spec}} is a valid expression,
returns
\begin{codeblock}
@\exposidnc{remove-kind}@(U::@\exposidnc{quantity-spec}@)
\end{codeblock}
\item
Otherwise, if \tcode{U::\exposidnc{reference-unit}} is a valid expression,
returns
\begin{codeblock}
@\exposidnc{get-associated-quantity}@(U::@\exposidnc{reference-unit}@)
\end{codeblock}
\item
Otherwise, if
\tcode{\exposidnc{is-derived-from-specialization-of}<U, \exposidnc{expr-fractions}>()}
is \tcode{true},
returns
\begin{codeblock}
@\exposidnc{expr-map}@<@\exposidnc{to-quantity-spec}@, derived_quantity_spec, struct dimensionless>(u)
\end{codeblock}
where \exposid{to-quantity-spec} is defined as follows:
\begin{codeblock}
template<@\libconcept{unit}@ U>
using @\exposidnc{to-quantity-spec}@ = decltype(@\exposidnc{get-associated-quantity}@(U{}));  // \expos
\end{codeblock}
\end{itemize}
\end{itemdescr}

\rSec3[qty.unit.sym.fmt]{Symbol formatting}

\indexlibrarymember{unit_symbol_to}{unit}
\begin{itemdecl}
template<typename CharT = char, std::@\stdconcept{output_iterator}@<CharT> Out, @\libconcept{unit}@ U>
constexpr Out unit_symbol_to(Out out, U u, const unit_symbol_formatting& fmt = {});
\end{itemdecl}

\begin{itemdescr}
\pnum
\effects
TBD.

\pnum
\returns
TBD.
\end{itemdescr}

\indexlibrarymember{unit_symbol}{unit}
\begin{itemdecl}
template<unit_symbol_formatting fmt = {}, typename CharT = char, @\libconcept{unit}@ U>
consteval std::string_view unit_symbol(U);
\end{itemdecl}

\begin{itemdescr}
\pnum
\effects
Equivalent to:
\begin{codeblock}
TBD.
\end{codeblock}
\end{itemdescr}

\rSec2[qty.ref.concepts]{Concepts}

\begin{itemdecl}
template<typename T>
concept @\defexposconceptnc{BareReference}@ =   // \expos
  @\libconcept{unit}@<T> || (@\exposidnc{is-specialization-of}@<T, reference>());

template<typename T>
concept @\defexposconceptnc{PointReference}@ =  // \expos
  @\exposidnc{is-specialization-of}@<T, point_>() && !@\libconcept{quantity_spec}@<T>;

template<typename T>
concept @\deflibconcept{quantity_reference}@ = @\exposconceptnc{BareReference}@<T> || @\exposconceptnc{PointReference}@<T>;
\end{itemdecl}

\pnum
A type \tcode{T} that satisfies \tcode{quantity_reference}
represents the reference of a quantity\irefiev{112-01-01}.

\pnum
A type that models \exposconceptnc{PointReference} is a specialization \tcode{point_<R, PO>}\iref{qty.point}
and represents the reference of a point,
a quantity measured from the point origin \tcode{PO}\iref{qty.pt.orig}.

\begin{itemdecl}
template<typename T, auto QS>
concept @\deflibconcept{quantity_reference_of}@ = @\libconcept{quantity_reference}@<T> && @\libconcept{quantity_spec_of}@<decltype(get_quantity_spec(T{})), QS>;
\end{itemdecl}

\rSec2[qty.ref.syn]{Class template \tcode{reference}}

\indexlibrarymember{operator*}{reference}
\indexlibrarymember{operator*}{unit}
\indexlibrarymember{operator/}{reference}
\indexlibrarymember{operator/}{unit}
\indexlibrarymember{inverse}{reference}
\indexlibrarymember{pow}{reference}
\indexlibrarymember{sqrt}{reference}
\indexlibrarymember{cbrt}{reference}
\begin{codeblock}
namespace std {

template<@\libconcept{quantity_spec}@ auto Q, @\libconcept{unit}@ auto U>
using @\exposidnc{reference-t}@ = reference<decltype(Q), decltype(U)>;  // \expos

template<@\libconcept{quantity_spec}@ Q, @\libconcept{unit}@ U>
struct @\libglobal{reference}@ {
  // \ref{qty.ref.ops}, operations

  template<typename Q2, typename U2>
  friend consteval auto operator*(reference, reference<Q2, U2>)
    -> @\exposidnc{reference-t}@<Q{} * Q2{}, U{} * U2{}>;

  template<@\libconcept{unit}@ U2>
  friend consteval auto operator*(reference, U2)
    -> @\exposidnc{reference-t}@<Q{} * get_quantity_spec(U2{}), U{} * U2{}>;

  template<@\libconcept{unit}@ U1>
  friend consteval auto operator*(U1, reference)
    -> @\exposidnc{reference-t}@<get_quantity_spec(U1{}) * Q{}, U1{} * U{}>;

  template<typename Q2, typename U2>
  friend consteval auto operator/(reference, reference<Q2, U2>)
    -> @\exposidnc{reference-t}@<Q{} / Q2{}, U{} / U2{}>;

  template<@\libconcept{unit}@ U2>
  friend consteval auto operator/(reference, U2)
    -> @\exposidnc{reference-t}@<Q{} / get_quantity_spec(U2{}), U{} / U2{}>;

  template<@\libconcept{unit}@ U1>
  friend consteval auto operator/(U1, reference)
    -> @\exposidnc{reference-t}@<get_quantity_spec(U1{}) / Q{}, U1{} / U{}>;

  friend consteval auto inverse(reference) -> @\exposidnc{reference-t}@<inverse(Q{}), inverse(U{})>;

  template<std::intmax_t Num, std::intmax_t Den = 1>
    requires(Den != 0)
  friend consteval auto pow(reference) -> @\exposidnc{reference-t}@<pow<Num, Den>(Q{}), pow<Num, Den>(U{})>;
  friend consteval auto sqrt(reference) -> @\exposidnc{reference-t}@<sqrt(Q{}), sqrt(U{})>;
  friend consteval auto cbrt(reference) -> @\exposidnc{reference-t}@<cbrt(Q{}), cbrt(U{})>;

  // \ref{qty.ref.cmp}, comparison

  template<typename Q2, typename U2>
  friend consteval bool operator==(reference, reference<Q2, U2>);

  template<@\libconcept{unit}@ U2>
  friend consteval bool operator==(reference, U2 u2);

  template<typename Q2, typename U2>
  friend consteval bool convertible(reference, reference<Q2, U2>);

  template<@\libconcept{unit}@ U2>
  friend consteval bool convertible(reference, U2 u2);

  template<@\libconcept{unit}@ U1>
  friend consteval bool convertible(U1 u1, reference);
};

}
\end{codeblock}

\pnum
\tcode{reference<Q, U>} represents the reference of a quantity\irefiev{112-01-01}.
The unit of measurement \tcode{U}\irefiev{112-01-14}
is used to measure a value of the quantity \tcode{Q}\irefiev{112-01-28}.
\begin{note}
\tcode{reference} is typically implicitly instantiated
when specifying that a unit measures a more specific quantity.
\begin{example}
\begin{codeblock}
using namespace si::unit_symbols;
auto x = 1 * m;                 // measures a length
auto y = 1 * isq::width[m];     // measures a width
auto z = 1 * isq::diameter[m];  // measures a diameter
\end{codeblock}
\end{example}
\end{note}

\rSec2[qty.ref.ops]{Operations}

\pnum
Each member function with a \fakegrammarterm{trailing-return-type}
of \tcode{-> \exposidnc{reference-t}<T...>}
returns \tcode{\{\}}.
It does not participate in overload resolution unless
every expression in \tcode{T...} is a valid constant expression.
\begin{note}
An operation on quantity specifications that is not formed\iref{qty.spec.ops},
such as the product of two vector quantities, is not formed on references either.
\end{note}

\indexlibrarymember{operator*}{Representation}
\indexlibrarymember{operator*}{quantity_reference}
\begin{itemdecl}
template<typename FwdRep, @\exposconceptnc{BareReference}@ R,
         @\libconcept{representation_of}@<get_quantity_spec(R{})> Rep = std::remove_cvref_t<FwdRep>>
  requires(!@\exposconceptnc{OffsetUnit}@<decltype(get_unit(R{}))>)
constexpr quantity<R{}, Rep> operator*(FwdRep&& lhs, R r);
\end{itemdecl}

\begin{itemdescr}
\pnum
\effects
Equivalent to:
\tcode{return quantity\{std::forward<FwdRep>(lhs), r\};}
\end{itemdescr}

\indexlibrarymember{operator*}{Representation}
\indexlibrarymember{operator*}{quantity_reference}
\begin{itemdecl}
template<typename FwdRep, @\exposconceptnc{BareReference}@ R,
         @\libconcept{representation_of}@<get_quantity_spec(R{})> Rep = std::remove_cvref_t<FwdRep>>
  requires(!@\exposconceptnc{OffsetUnit}@<decltype(get_unit(R{}))>)
constexpr @\libconcept{delta_quantity}@ auto operator/(FwdRep&& lhs, R);
\end{itemdecl}

\begin{itemdescr}
\pnum
\effects
Equivalent to:
\tcode{return quantity\{std::forward<FwdRep>(lhs), inverse(R\{\})\};}
\end{itemdescr}

\indexlibrarymember{operator*}{quantity}
\indexlibrarymember{operator*}{quantity_reference}
\indexlibrarymember{operator/}{quantity}
\indexlibrarymember{operator/}{quantity_reference}
\begin{itemdecl}
template<typename FwdQ, @\exposconceptnc{BareReference}@ R, @\libconcept{delta_quantity}@ Q = std::remove_cvref_t<FwdQ>>
constexpr @\libconcept{delta_quantity}@ auto operator*(FwdQ&& q, R);
template<typename FwdQ, @\exposconceptnc{BareReference}@ R, @\libconcept{delta_quantity}@ Q = std::remove_cvref_t<FwdQ>>
constexpr @\libconcept{delta_quantity}@ auto operator/(FwdQ&& q, R);
\end{itemdecl}

\begin{itemdescr}
\pnum
Let \tcode{@} be the \fakegrammarterm{operator}.

\pnum
\effects
Equivalent to:
\begin{codeblock}
return quantity{std::forward<FwdQ>(q).@\exposidnc{numerical-value}@, Q::reference @\atsign@ R{}};
\end{codeblock}
\end{itemdescr}

\begin{itemdecl}
template<@\exposconceptnc{BareReference}@ R, typename Rep>
  requires @\libconcept{representation_of}@<std::remove_cvref_t<Rep>, get_quantity_spec(R{})>
constexpr auto operator*(R, Rep&&) = delete;

template<@\exposconceptnc{BareReference}@ R, typename Rep>
  requires @\libconcept{representation_of}@<std::remove_cvref_t<Rep>, get_quantity_spec(R{})>
constexpr auto operator/(R, Rep&&) = delete;

template<@\exposconceptnc{BareReference}@ R, typename Q>
  requires @\libconcept{delta_quantity}@<std::remove_cvref_t<Q>>
constexpr auto operator*(R, Q&&) = delete;

template<@\exposconceptnc{BareReference}@ R, typename Q>
  requires @\libconcept{delta_quantity}@<std::remove_cvref_t<Q>>
constexpr auto operator/(R, Q&&) = delete;
\end{itemdecl}

\begin{itemdescr}
\pnum
\recommended
Suggest swapping the operands.
\end{itemdescr}

\rSec2[qty.ref.cmp]{Comparison}

\indexlibrarymember{operator==}{reference}
\begin{itemdecl}
template<typename Q2, typename U2>
friend consteval bool operator==(reference, reference<Q2, U2>);
\end{itemdecl}

\begin{itemdescr}
\pnum
\returns
\tcode{Q\{\} == Q2\{\} \&\& U\{\} == U2\{\}}.
\end{itemdescr}

\indexlibrarymember{operator==}{reference}
\indexlibrarymember{operator==}{unit}
\begin{itemdecl}
template<@\libconcept{unit}@ U2>
friend consteval bool operator==(reference, U2 u2);
\end{itemdecl}

\begin{itemdescr}
\pnum
\returns
\tcode{Q\{\} == get_quantity_spec(u2) \&\& U\{\} == u2}.
\end{itemdescr}

\indexlibrarymember{convertible}{reference}
\begin{itemdecl}
template<typename Q2, typename U2>
friend consteval bool convertible(reference, reference<Q2, U2>);
\end{itemdecl}

\begin{itemdescr}
\pnum
\returns
\tcode{\exposid{implicitly-convertible}(Q\{\}, Q2\{\}) \&\& convertible(U\{\}, U2\{\})}.
\end{itemdescr}

\indexlibrarymember{convertible}{reference}
\indexlibrarymember{convertible}{unit}
\begin{itemdecl}
template<@\libconcept{unit}@ U2>
friend consteval bool convertible(reference, U2 u2);
\end{itemdecl}

\begin{itemdescr}
\pnum
\returns
\tcode{\exposid{implicitly-convertible}(Q\{\}, get_quantity_spec(u2)) \&\& convertible(U\{\}, u2)}.
\end{itemdescr}

\indexlibrarymember{convertible}{reference}
\indexlibrarymember{convertible}{unit}
\begin{itemdecl}
template<@\libconcept{unit}@ U1>
friend consteval bool convertible(U1 u1, reference);
\end{itemdecl}

\begin{itemdescr}
\pnum
\returns
\tcode{\exposid{implicitly-convertible}(get_quantity_spec(u1), Q\{\}) \&\& convertible(u1, U\{\})}.
\end{itemdescr}

\rSec2[qty.ref.obs]{Observers}

\indexlibrarymember{get_quantity_spec}{reference}
\begin{itemdecl}
template<typename Q, typename U>
consteval @\libconcept{quantity_spec}@ auto get_quantity_spec(reference<Q, U>);
\end{itemdecl}

\begin{itemdescr}
\pnum
\returns
\tcode{Q\{\}}.
\end{itemdescr}

\indexlibrarymember{get_unit}{reference}
\begin{itemdecl}
template<typename Q, typename U>
consteval @\libconcept{unit}@ auto get_unit(reference<Q, U>);
\end{itemdecl}

\begin{itemdescr}
\pnum
\returns
\tcode{U\{\}}.
\end{itemdescr}

\indexlibrarymember{get_common_reference}{unit}
\begin{itemdecl}
consteval @\libconcept{unit}@ auto @\libglobal{get_common_reference}@(@\libconcept{unit}@ auto u1, @\libconcept{unit}@ auto u2,
                                                       @\libconcept{unit}@ auto... rest)
  requires @\seebelownc@;
\end{itemdecl}

\begin{itemdescr}
\pnum
\returns
\tcode{\exposid{get-common-unit}(u1, u2, rest...)}.

\pnum
\remarks
The expression in the \fakegrammarterm{requires-clause} is equivalent to:
\begin{codeblock}
requires {
  get_common_quantity_spec(get_quantity_spec(u1), get_quantity_spec(u2),
                           get_quantity_spec(rest)...);
  { @\exposid{get-common-unit}@(u1, u2, rest...) } -> @\libconcept{unit}@;
}
\end{codeblock}
\end{itemdescr}

\indexlibrarymember{get_common_reference}{reference}
\begin{itemdecl}
template<@\exposconceptnc{BareReference}@ R1, @\exposconceptnc{BareReference}@ R2, @\exposconceptnc{BareReference}@... Rest>
consteval @\libconcept{quantity_reference}@ auto @\libglobal{get_common_reference}@(R1 r1, R2 r2, Rest... rest)
  requires @\seebelownc@;
\end{itemdecl}

\begin{itemdescr}
\pnum
\returns
\begin{codeblock}
@\exposidnc{reference-t}@<get_common_quantity_spec(get_quantity_spec(R1{}), get_quantity_spec(R2{}),
                                     get_quantity_spec(rest)...),
            @\exposid{get-common-unit}@(get_unit(R1{}), get_unit(R2{}), get_unit(rest)...)>{};
\end{codeblock}

\pnum
\remarks
The expression in the \fakegrammarterm{requires-clause} is equivalent to:
\begin{codeblock}
requires {
  get_common_quantity_spec(get_quantity_spec(r1), get_quantity_spec(r2),
                           get_quantity_spec(rest)...);
  @\exposid{get-common-unit}@(get_unit(r1), get_unit(r2), get_unit(rest)...);
}
\end{codeblock}
\end{itemdescr}

\rSec1[qty.rep]{Representation}

\rSec2[qty.rep.general]{General}

\pnum
Subclause \ref{qty.rep} specifies the components
used to constrain the numerical value of a quantity\irefiev{112-01-29}.

\rSec2[qty.rep.traits]{Traits}

\rSec3[qty.fp.traits]{Floating-point}

\begin{itemdecl}
template<typename T>
struct @\exposidnc{actual-value-type}@ : @\exposidnc{cond-value-type}@<T> {}; // see \refcpp{readable.traits}

template<typename T>
  requires(!std::is_pointer_v<T> && !std::is_array_v<T>) &&
          requires { typename std::indirectly_readable_traits<T>::value_type; }
struct @\exposidnc{actual-value-type}@<T> : std::indirectly_readable_traits<T> {};

template<typename T>
using @\exposidnc{actual-value-type-t}@ = @\exposidnc{actual-value-type}@<T>::value_type;

template<typename Rep>
constexpr bool @\libglobal{treat_as_floating_point}@ =
  std::chrono::treat_as_floating_point_v<@\exposidnc{actual-value-type-t}@<Rep>>;
\end{itemdecl}

\begin{itemdescr}
\pnum
\tcode{quantity} uses \tcode{treat_as_floating_point}
to help determine whether implicit conversions are allowed among its specializations.

\pnum
\remarks
Pursuant to \refcpp{namespace.std}\iref{spec.ext},
users may specialize \tcode{treat_as_floating_point}
for cv-unqualified program-defined types.
Such specializations shall be usable in constant expressions\irefcpp{expr.const}
and have type \tcode{const bool}.
\end{itemdescr}

\rSec3[qty.char.traits]{Quantity character}

\indexlibraryglobal{disable_scalar}
\indexlibraryglobal{disable_complex}
\indexlibraryglobal{disable_vector}
\begin{itemdecl}
template<typename T>
constexpr bool disable_scalar = false;
template<typename T>
constexpr bool disable_complex = false;
template<typename T>
constexpr bool disable_vector = false;
\end{itemdecl}

\begin{itemdescr}
\pnum
Some quantities are defined as having a numerical value\irefiev{112-01-29} of a specific set\irefiev{102-01-02}.
The representation concepts use these traits
to help determine the sets \tcode{T} represents.

\pnum
\remarks
Pursuant to \refcpp{namespace.std}\iref{spec.ext},
users may specialize these templates
for cv-unqualified program-defined types.
Such specializations shall be usable in constant expressions\irefcpp{expr.const}
and have type \tcode{const bool}.

\pnum
\begin{note}
These templates prevent use of representation types with the library
that satisfy but do not in fact model their corresponding concept.
\end{note}
\end{itemdescr}

\rSec3[qty.val.traits]{Values}

\pnum
\tcode{quantity} uses \tcode{representation_values}
to construct special values of its representation type.

\begin{codeblock}
namespace std {

template<typename Rep>
struct @\libglobal{representation_values}@ : std::chrono::duration_values<Rep> {
  static constexpr Rep one() noexcept;
};

}
\end{codeblock}

\pnum
The requirements on \tcode{std::chrono::duration_values<Rep>}\irefcpp{time.traits.duration.values}
also apply to \tcode{representation_values<Rep>}.

\ednote{
The reference implementation provides \tcode{zero()} and \tcode{one()} only for a scalar
representation type. \tcode{Rep(0)} isn't a zero vector for every vector type,
so a vector or a tensor type specializes \tcode{representation_values} to provide \tcode{zero()}.
This is to be reflected here together with the wording of the tag type for zero (see P4406R0).
}

\begin{itemdecl}
static constexpr Rep one() noexcept;
\end{itemdecl}

\begin{itemdescr}
\pnum
\returns
\tcode{Rep(1)}.

\pnum
\remarks
The value returned shall be the neutral element for multiplication\irefiev{102-01-19}.
\end{itemdescr}

\rSec2[qty.rep.cpos]{Customization point objects}

\rSec3[qty.rep.cpos.general]{General}

\pnum
Within subclause \ref{qty.rep.cpos},
\term{reified object} is as defined in \refcpp{range.access.general}.

\rSec3[qty.real.cpo]{\tcode{math::real}}

\pnum
The name \tcode{math::real} denotes a customization point object\irefcpp{customization.point.object}.

\pnum
Given a subexpression \tcode{E} with type \tcode{T},
let \tcode{t} be an lvalue that denotes the reified object for \tcode{E}.
Then:
\begin{itemize}
\item
If \tcode{T} does not model \exposconcept{WeaklyRegular},
\tcode{math::real(E)} is ill-formed.
\item
If \tcode{auto(t.real())} is a valid expression whose type models \exposconcept{Scalar},
\tcode{math::real(E)} is expression-equivalent to \tcode{auto(t.real())}.
\item
Otherwise, if \tcode{T} is a class or enumeration type and
\tcode{auto(real(t))} is a valid expression whose type models \exposconcept{Scalar}
where the meaning of \tcode{real} is established as-if by performing argument-dependent lookup only\irefcpp{basic.lookup.argdep},
then \tcode{math::real(E)} is expression-equivalent to that expression.
\item
Otherwise, \tcode{math::real(E)} is ill-formed.
\end{itemize}

\rSec3[qty.imag.cpo]{\tcode{math::imag}}

\pnum
The name \tcode{math::imag} denotes a customization point object\irefcpp{customization.point.object}.

\pnum
Given a subexpression \tcode{E} with type \tcode{T},
let \tcode{t} be an lvalue that denotes the reified object for \tcode{E}.
Then:
\begin{itemize}
\item
If \tcode{T} does not model \exposconcept{WeaklyRegular},
\tcode{math::imag(E)} is ill-formed.
\item
If \tcode{auto(t.imag())} is a valid expression whose type models \exposconcept{Scalar},
\tcode{math::imag(E)} is expression-equivalent to \tcode{auto(t.imag())}.
\item
Otherwise, if \tcode{T} is a class or enumeration type and
\tcode{auto(imag(t))} is a valid expression whose type models \exposconcept{Scalar}
where the meaning of \tcode{imag} is established as-if by performing argument-dependent lookup only\irefcpp{basic.lookup.argdep},
then \tcode{math::imag(E)} is expression-equivalent to that expression.
\item
Otherwise, \tcode{math::imag(E)} is ill-formed.
\end{itemize}

\rSec3[qty.modulus.cpo]{\tcode{math::modulus}}

\pnum
The name \tcode{math::modulus} denotes a customization point object\irefcpp{customization.point.object}.

\pnum
Given a subexpression \tcode{E} with type \tcode{T},
let \tcode{t} be an lvalue that denotes the reified object for \tcode{E}.
Then:
\begin{itemize}
\item
If \tcode{T} does not model \exposconcept{WeaklyRegular},
\tcode{math::modulus(E)} is ill-formed.
\item
If \tcode{auto(t.modulus())} is a valid expression whose type models \exposconcept{Scalar},
\tcode{math::modulus(E)} is expression-equivalent to \tcode{auto(t.modulus())}.
\item
Otherwise, if \tcode{T} is a class or enumeration type and
\tcode{auto(modulus(t))} is a valid expression whose type models \exposconcept{Scalar}
where the meaning of \tcode{modulus} is established as-if by performing argument-dependent lookup only\irefcpp{basic.lookup.argdep},
then \tcode{math::modulus(E)} is expression-equivalent to that expression.
\item
If \tcode{auto(t.abs())} is a valid expression whose type models \exposconcept{Scalar},
\tcode{math::modulus(E)} is expres\-sion-equivalent to \tcode{auto(t.abs())}.
\item
Otherwise, if \tcode{T} is a class or enumeration type and
\tcode{auto(abs(t))} is a valid expression whose type models \exposconcept{Scalar}
where the meaning of \tcode{abs} is established as-if by performing argument-dependent lookup only\irefcpp{basic.lookup.argdep},
then \tcode{math::modulus(E)} is expression-equivalent to that expression.
\item
Otherwise, \tcode{math::modulus(E)} is ill-formed.
\end{itemize}

\rSec3[qty.mag.cpo]{\tcode{math::magnitude}}

\pnum
The name \tcode{math::magnitude} denotes a customization point object\irefcpp{customization.point.object}.

\pnum
Given a subexpression \tcode{E} with type \tcode{T},
let \tcode{t} be an lvalue that denotes the reified object for \tcode{E}.
Then:
\begin{itemize}
\item
If \tcode{T} does not model \exposconcept{WeaklyRegular},
\tcode{math::magnitude(E)} is ill-formed.
\item
If \tcode{auto(t.magnitude())} is a valid expression whose type models \exposconcept{Scalar},
\tcode{math::magnitude(E)} is expression-equivalent to \tcode{auto(t.magnitude())}.
\item
Otherwise, if \tcode{T} is a class or enumeration type and
\tcode{auto(magnitude(t))} is a valid expression whose type models \exposconcept{Scalar}
where the meaning of \tcode{magnitude} is established as-if by performing argument-dependent lookup only\irefcpp{basic.lookup.argdep},
then \tcode{math::magnitude(E)} is expression-equivalent to that expression.
\item
Otherwise, if \tcode{auto(t.abs())} is a valid expression whose type models \exposconcept{Scalar},
\tcode{math::magnitude(\brk{}E)} is expression-equivalent to \tcode{auto(t.abs())}.
\item
Otherwise, if \tcode{T} is a class or enumeration type and
\tcode{auto(abs(t))} is a valid expression whose type models \exposconcept{Scalar}
where the meaning of \tcode{magnitude} is established as-if by performing argument-dependent lookup only\irefcpp{basic.lookup.argdep},
then \tcode{math::magnitude(E)} is expression-equivalent to that expression.
\item
Otherwise, \tcode{math::magnitude(E)} is ill-formed.
\end{itemize}

\rSec2[qty.rep.concepts]{Concepts}

\begin{itemdecl}
template<typename T>
concept @\defexposconceptnc{WeaklyRegular}@ = std::@\stdconcept{copyable}@<T> && std::@\stdconcept{equality_comparable}@<T>;  // \expos

template<typename T>
concept @\defexposconceptnc{Scalar}@ = (!disable_scalar<T>) && @\exposconceptnc{WeaklyRegular}@<T> && requires(T a, T b) {  // \expos
  // scalar operations
  { -a } -> std::common_with<T>;
  { a + b } -> std::common_with<T>;
  { a - b } -> std::common_with<T>;
  { a * b } -> std::common_with<T>;
  { a / b } -> std::common_with<T>;
};
\end{itemdecl}

\pnum
TBD.

\begin{itemdecl}
template<typename T>
using @\exposidnc{value-type-t}@ = @\exposidnc{actual-value-type-t}@<T>;  // \expos, see \ref{qty.fp.traits}

template<typename T>
concept @\defexposconceptnc{Complex}@ =  // \expos
  (!disable_complex<T>) && @\exposconceptnc{WeaklyRegular}@<T> && @\exposconceptnc{Scalar}@<@\exposidnc{value-type-t}@<T>> &&
  std::constructible_from<T, @\exposidnc{value-type-t}@<T>, @\exposidnc{value-type-t}@<T>> &&
  requires(T a, T b, @\exposidnc{value-type-t}@<T> s) {
    // complex operations
    { -a } -> std::common_with<T>;
    { a + b } -> std::common_with<T>;
    { a - b } -> std::common_with<T>;
    { a * b } -> std::common_with<T>;
    { a / b } -> std::common_with<T>;
    { a * s } -> std::common_with<T>;
    { s * a } -> std::common_with<T>;
    { a / s } -> std::common_with<T>;
    math::real(a);
    math::imag(a);
    math::modulus(a);
  };
\end{itemdecl}

\pnum
TBD.

\begin{itemdecl}
template<typename T>
concept @\defexposconceptnc{Vector}@ =  // \expos
  (!disable_vector<T>) && @\exposconceptnc{WeaklyRegular}@<T> && @\exposconceptnc{Scalar}@<@\exposidnc{value-type-t}@<T>> &&
  requires(T a, T b, @\exposidnc{value-type-t}@<T> s) {
    // vector operations
    { -a } -> std::common_with<T>;
    { a + b } -> std::common_with<T>;
    { a - b } -> std::common_with<T>;
    { a * s } -> std::common_with<T>;
    { s * a } -> std::common_with<T>;
    { a / s } -> std::common_with<T>;
    math::magnitude(a);
  };
\end{itemdecl}

\pnum
TBD.

\begin{itemdecl}
template<typename T>
using @\exposidnc{scaling-factor-type-t}@ =  // \expos
  std::conditional_t<treat_as_floating_point<T>, long double, std::intmax_t>;

template<typename T>
concept @\defexposconceptnc{ScalarRepresentation}@ =  // \expos
  (!@\exposidnc{is-specialization-of}@<T, quantity>()) && @\exposconceptnc{Scalar}@<T> &&
  requires(T a, T b, @\exposidnc{scaling-factor-type-t}@<T> f) {
    // scaling
    { a * f } -> std::common_with<T>;
    { f * a } -> std::common_with<T>;
    { a / f } -> std::common_with<T>;
  };
\end{itemdecl}

\pnum
TBD.

\begin{itemdecl}
template<typename T>
concept @\defexposconceptnc{ComplexRepresentation}@ =  // \expos
  (!@\exposidnc{is-specialization-of}@<T, quantity>()) && @\exposconceptnc{Complex}@<T> &&
  requires(T a, T b, @\exposidnc{scaling-factor-type-t}@<T> f) {
    // scaling
    { a * T(f) } -> std::common_with<T>;
    { T(f) * a } -> std::common_with<T>;
    { a / T(f) } -> std::common_with<T>;
  };
\end{itemdecl}

\pnum
TBD.

\begin{itemdecl}
template<typename T>
concept @\defexposconceptnc{VectorRepresentation}@ =  // \expos
  (!@\exposidnc{is-specialization-of}@<T, quantity>()) && @\exposconceptnc{Vector}@<T>;
\end{itemdecl}

\pnum
TBD.

\begin{itemdecl}
template<typename T>
concept @\deflibconcept{Representation}@ = @\exposconceptnc{ScalarRepresentation}@<T> || @\exposconceptnc{ComplexRepresentation}@<T> ||
                         @\exposconceptnc{VectorRepresentation}@<T>;
\end{itemdecl}

\pnum
A type \tcode{T} models \libconcept{Representation} if
it represents the numerical value of a quantity\irefiev{112-01-29}.

\begin{itemdecl}
template<typename T, quantity_character Ch>
concept @\defexposconceptnc{IsOfCharacter}@ = (Ch == quantity_character::real_scalar && @\exposconceptnc{Scalar}@<T>) ||  // \expos
                        (Ch == quantity_character::complex_scalar && @\exposconceptnc{Complex}@<T>) ||
                        (Ch == quantity_character::vector && @\exposconceptnc{Vector}@<T>);

template<typename T, auto V>
concept @\deflibconcept{representation_of}@ =
  @\libconcept{Representation}@<T> && ((@\libconcept{quantity_spec}@<decltype(V)> &&
                         (@\exposconceptnc{QuantityKindSpec}@<decltype(@\exposidnc{remove-point}@(V))> ||
                          @\exposconceptnc{IsOfCharacter}@<T, @\exposidnc{remove-point}@(V).character>)) ||
                        (std::@\stdconcept{same_as}@<quantity_character, decltype(V)> && @\exposconceptnc{IsOfCharacter}@<T, V>));
\end{itemdecl}

\pnum
A type \tcode{T} models \tcode{\libconcept{representation_of}<V>} if
\tcode{T} models \libconcept{Representation} and
\begin{itemize}
\item
\tcode{V} is a kind of quantity or the quantity specification of a point on the scale of one, or
\item
if \tcode{V} is a quantity, or the quantity specification of a point on its scale,
then \tcode{T} represents a value of its character, or
\item
if \tcode{V} is a quantity character, then \tcode{T} represents a value of \tcode{V}.
\end{itemize}

\rSec1[qty]{Quantity}

\rSec2[qty.general]{General}

\pnum
Subclause \ref{qty} describes the class template \tcode{quantity}
that represents the value of a quantity\irefiev{112-01-28},
and the point origins and point references it uses.

\pnum
A \defn{delta} is the value of a quantity
that is an element of a vector space\irefiev{102-03-01,102-03-04}.
A \defn{point} is the value of a quantity
that is an element of an affine space\irefiev{102-03-02,102-04-01}.
A specialization \tcode{quantity<R, Rep>} represents a point
if \tcode{decltype(R)} models \exposconcept{PointReference}\iref{qty.ref.concepts},
and a delta otherwise.

\rSec2[qty.pt.orig]{Point origin}

\rSec3[qty.pt.orig.general]{General}

\pnum
This subclause specifies the components
for defining the origin of an affine space.
An \defn{origin} is a point from which measurements\irefiev{112-04-01} take place.

\rSec3[qty.pt.orig.concepts]{Concepts}

\begin{itemdecl}
template<typename T>
concept @\defexposconceptnc{PointOrigin}@ =                     // \expos
  @\exposconceptnc{SymbolicConstant}@<T> && std::@\stdconcept{derived_from}@<T, @\exposidnc{point-origin-interface}@>;

template<typename T, auto QS>
concept @\deflibconcept{point_origin_for}@ =
  @\exposconceptnc{PointOrigin}@<T> &&
  @\libconcept{quantity_spec_of}@<decltype(@\exposidnc{remove-point}@(QS)), @\exposidnc{remove-point}@(T::@\exposidnc{quantity-spec}@)>;

template<typename T, auto V>
concept @\defexposconceptnc{SameAbsolutePointOriginAs}@ =  // \expos
  @\exposconceptnc{PointOrigin}@<T> && @\exposconceptnc{PointOrigin}@<decltype(V)> && @\exposidnc{same-absolute-point-origins}@(T{}, V);
\end{itemdecl}

\pnum
The exposition-only concept \exposconceptnc{PointOrigin} is satisfied by the types of point origins.

\pnum
\begin{note}
\tcode{QS} names either the quantity of the deltas measured from the origin, \tcode{delta} of it,
or the quantity of the points,
so \tcode{\libconcept{point_origin_for}<T, isq::altitude>},
\tcode{\libconcept{point_origin_for}<T, delta<isq::altitude>>}, and
\tcode{\libconcept{point_origin_for}<T, point<isq::altitude>>} are equivalent.
\end{note}

\rSec3[qty.pt.orig.types]{Types}

\rSec4[qty.abs.pt.orig]{Absolute}

\begin{codeblock}
namespace std {

template<@\libconcept{quantity_spec}@ auto QS>
struct @\libglobal{absolute_point_origin}@ : @\exposidnc{point-origin-interface}@ {
  static constexpr @\libconcept{quantity_spec}@ auto @\exposidnc{quantity-spec}@ = @\exposidnc{as-point}@(QS);  // \expos
};

}
\end{codeblock}

\pnum
An \defnadj{absolute}{origin} is an origin
chosen by convention and not defined in terms of another origin.
A specialization of \tcode{absolute_point_origin} is used as a base type when defining an absolute origin.
\tcode{QS} is the quantity the origin represents.
\begin{note}
An origin is a point,
so its quantity specification \exposid{quantity-spec} is the one of a point,
whether \tcode{QS} is \tcode{isq::altitude} or \tcode{point<isq::altitude>}.
\end{note}

\pnum
The quantity specification \tcode{get_quantity_spec(po)}\iref{qty.pt.orig.ops}
of an absolute origin \tcode{po} of type \tcode{PO} derived from \tcode{absolute_point_origin<QS>} is
\tcode{\exposid{normalized-point-spec}(\exposid{remove-point}(PO::\exposid{quantity-spec}), po)}\iref{point.qty}:
\tcode{point<isq::altitude, mean_sea_level>} for an origin
\tcode{struct mean_sea_level : absolute_point_origin<isq::altitude> \{\}},
and \tcode{point<S>} for a natural point origin of the scale \tcode{S}.

\rSec4[qty.rel.pt.orig]{Relative}

\begin{codeblock}
namespace std {

template<@\libconcept{point_quantity}@ auto QP>
struct @\libglobal{relative_point_origin}@ : @\exposidnc{point-origin-interface}@ {
  static constexpr @\libconcept{point_quantity}@ auto @\exposidnc{quantity-point}@ = QP;       // \expos
  static constexpr @\libconcept{quantity_spec}@ auto @\exposidnc{quantity-spec}@ = @\seebelownc@;  // \expos
  static constexpr @\exposconceptnc{PointOrigin}@ auto @\exposidnc{absolute-point-origin}@ =      // \expos
    QP.absolute_point_origin;
};

}
\end{codeblock}

\pnum
A \defnadj{relative}{origin} is an origin
of a subspace\irefiev{102-03-03}.
A specialization of \tcode{relative_point_origin} is used as a base type when defining a relative origin \placeholder{O}.
\placeholder{O} is offset from \tcode{QP.absolute_point_origin}
by \tcode{QP.quantity_from(QP.absolute_point_origin)}.

\pnum
The member \exposid{quantity-spec} is equal to
\tcode{QP.point_origin.\exposidnc{quantity-spec}} if
\begin{codeblock}
@\exposconceptnc{QuantityKindSpec}@<decltype(@\exposidnc{remove-point}@(QP.quantity_spec))>
\end{codeblock}
is satisfied, and
to \tcode{\exposid{as-point}(\exposid{remove-point}(QP.quantity_spec))} otherwise.

\pnum
The quantity specification \tcode{get_quantity_spec(po)}\iref{qty.pt.orig.ops}
of a relative origin \tcode{po} of type \tcode{PO} derived from \tcode{relative_point_origin<QP>} is
\tcode{\exposid{normalized-point-spec}(\exposid{remove-point}(PO::\exposid{quantity-spec}), PO::\exposid{absolute-point-origin})}.
\begin{note}
The origin takes the more specific of the two scales:
given \tcode{struct msl : absolute_point_origin<isq::altitude> \{\}},
the quantity specification of a relative origin defined by \tcode{msl + 42 * m} is
\tcode{point<isq::altitude, msl>}, and
that of one defined by \tcode{msl + isq::height(42 * m)} is \tcode{point<isq::height, msl>}.
\end{note}

\rSec4[qty.natural.pt.orig]{Natural}

\begin{codeblock}
namespace std {

template<@\libconcept{quantity_spec}@ auto QS>
  requires(!@\exposconceptnc{PointQuantitySpec}@<decltype(QS)>) && (!@\exposconceptnc{DeltaQuantitySpec}@<decltype(QS)>)
struct @\libglobal{natural_point_origin_}@ : absolute_point_origin<QS> {};

}
\end{codeblock}

\pnum
\tcode{natural_point_origin_<QS>} represents an origin
chosen by convention as the value $0$ of the quantity \tcode{QS}.
A \defnadj{natural}{point origin} is a point origin
of a class type derived from a specialization of \tcode{natural_point_origin_}.
\begin{note}
\tcode{natural_point_origin<isq::height>}, \tcode{natural_point_origin<delta<isq::height>>},
and \tcode{natural_point_origin<point<isq::height>>} denote the same origin.
A class derived from \tcode{natural_point_origin_<QS>} is the natural point origin of \tcode{QS}
under a name of its own.
\end{note}

\pnum
Two natural point origins compare equal
when the scale of one of them is implicitly convertible to the scale of the other\iref{qty.pt.orig.ops}:
the natural point origins of a quantity, of its ancestors and descendants, and of its kind count as one.
\begin{note}
\tcode{natural_point_origin<kind_of<isq::length>>}, \tcode{natural_point_origin<isq::length>}, and
\tcode{natural_point_origin<isq::height>} are equal to one another, while
\tcode{natural_point_origin<isq::height>} and \tcode{natural_point_origin<isq::width>} are not:
the natural origins of two siblings are the zeros of different scales,
along axes that can point in different directions.
The relation is not transitive.
\end{note}

\rSec3[qty.pt.orig.ops]{Operations}

\begin{codeblock}
namespace std {

struct @\exposidnc{point-origin-interface}@ {
  template<@\exposconceptnc{PointOrigin}@ PO>
  friend consteval @\libconcept{quantity_spec}@ auto get_quantity_spec(PO po);

  template<@\exposconceptnc{PointOrigin}@ PO, typename FwdQ, @\libconcept{delta_quantity}@ Q = std::remove_cvref_t<FwdQ>>
    requires @\libconcept{point_origin_for}@<PO, Q::quantity_spec>
  friend constexpr @\libconcept{point_quantity}@ auto operator+(PO, FwdQ&& q);
  template<typename FwdQ, @\exposconceptnc{PointOrigin}@ PO, @\libconcept{delta_quantity}@ Q = std::remove_cvref_t<FwdQ>>
    requires @\libconcept{point_origin_for}@<PO, Q::quantity_spec>
  friend constexpr @\libconcept{point_quantity}@ auto operator+(FwdQ&& q, PO);

  template<@\exposconceptnc{PointOrigin}@ PO, @\libconcept{delta_quantity}@ Q>
    requires @\libconcept{point_origin_for}@<PO, Q::quantity_spec>
  friend constexpr @\libconcept{point_quantity}@ auto operator-(PO po, const Q& q)
    requires requires { -q; };

  template<@\exposconceptnc{PointOrigin}@ PO1, @\exposconceptnc{SameAbsolutePointOriginAs}@<PO1{}> PO2>
    requires @\seebelownc@
  friend constexpr @\libconcept{delta_quantity}@ auto operator-(PO1 po1, PO2 po2);

  template<@\exposconceptnc{PointOrigin}@ PO1, @\exposconceptnc{PointOrigin}@ PO2>
  friend consteval bool operator==(PO1 po1, PO2 po2);

  template<@\libconcept{point_quantity}@ QP, @\exposconceptnc{PointOrigin}@ PO2>
    requires @\libconcept{quantity_of}@<QP, PO2{}> &&
             @\libconcept{quantity_reference_of}@<decltype(auto(QP::reference)), get_quantity_spec(PO2{})>
  friend constexpr @\libconcept{delta_quantity}@ auto operator-(const QP& qp, PO2 po);
  template<@\exposconceptnc{PointOrigin}@ PO2, @\libconcept{point_quantity}@ QP>
    requires @\libconcept{quantity_of}@<QP, PO2{}> &&
             @\libconcept{quantity_reference_of}@<decltype(auto(QP::reference)), get_quantity_spec(PO2{})>
  friend constexpr @\libconcept{delta_quantity}@ auto operator-(PO2 po, const QP& qp);
};

}
\end{codeblock}

\pnum
Every point origin is derived from \exposid{point-origin-interface},
which provides the observer and the operators of a point origin.
The affine space operators with a point origin operand are hidden friends of \exposid{point-origin-interface},
not of \exposid{quantity-interface}\iref{qty.syn},
so that only an expression with a point origin operand considers them.

\indexlibrarymemberexpos{get_quantity_spec}{point-origin-interface}
\begin{itemdecl}
template<@\exposconceptnc{PointOrigin}@ PO>
friend consteval @\libconcept{quantity_spec}@ auto get_quantity_spec(PO po);
\end{itemdecl}

\begin{itemdescr}
\pnum
\returns
The quantity specification of \tcode{po}\iref{qty.abs.pt.orig,qty.rel.pt.orig},
a point quantity specification that names the absolute point origin of \tcode{po}
unless that is the natural point origin of its scale.
\end{itemdescr}

\indexlibrarymemberexpos{operator+}{point-origin-interface}
\indexlibrarymemberexpos{operator+}{quantity}
\begin{itemdecl}
template<@\exposconceptnc{PointOrigin}@ PO, typename FwdQ, @\libconcept{delta_quantity}@ Q = std::remove_cvref_t<FwdQ>>
  requires @\libconcept{point_origin_for}@<PO, Q::quantity_spec>
friend constexpr @\libconcept{point_quantity}@ auto operator+(PO, FwdQ&& q);
template<typename FwdQ, @\exposconceptnc{PointOrigin}@ PO, @\libconcept{delta_quantity}@ Q = std::remove_cvref_t<FwdQ>>
  requires @\libconcept{point_origin_for}@<PO, Q::quantity_spec>
friend constexpr @\libconcept{point_quantity}@ auto operator+(FwdQ&& q, PO);
\end{itemdecl}

\begin{itemdescr}
\pnum
\effects
Equivalent to:
\tcode{return quantity\{std::forward<FwdQ>(q), PO\{\}\};}
\end{itemdescr}

\indexlibrarymemberexpos{operator-}{point-origin-interface}
\indexlibrarymemberexpos{operator-}{quantity}
\begin{itemdecl}
template<@\exposconceptnc{PointOrigin}@ PO, @\libconcept{delta_quantity}@ Q>
  requires @\libconcept{point_origin_for}@<PO, Q::quantity_spec>
friend constexpr @\libconcept{point_quantity}@ auto operator-(PO po, const Q& q)
  requires requires { -q; };
\end{itemdecl}

\begin{itemdescr}
\pnum
\effects
Equivalent to:
\tcode{return po + -q;}
\end{itemdescr}

\indexlibrarymemberexpos{operator-}{point-origin-interface}
\begin{itemdecl}
template<@\exposconceptnc{PointOrigin}@ PO1, @\exposconceptnc{SameAbsolutePointOriginAs}@<PO1{}> PO2>
  requires @\seebelownc@
friend constexpr @\libconcept{delta_quantity}@ auto operator-(PO1 po1, PO2 po2);
\end{itemdecl}

\begin{itemdescr}
\pnum
\effects
Equivalent to:
\begin{codeblock}
if constexpr (@\exposidnc{is-derived-from-specialization-of}@<PO1, absolute_point_origin>()) {
  return po1 - po2.@\exposidnc{quantity-point}@;
} else if constexpr (@\exposidnc{is-derived-from-specialization-of}@<PO2, absolute_point_origin>()) {
  return po1.@\exposidnc{quantity-point}@ - po2;
} else {
  return po1.@\exposidnc{quantity-point}@ - po2.@\exposidnc{quantity-point}@;
}
\end{codeblock}

\pnum
\remarks
The expression in the \fakegrammarterm{requires-clause} is equivalent to:
\begin{codeblock}
(@\libconcept{point_origin_for}@<PO2, PO1::@\exposidnc{quantity-spec}@> ||
 @\libconcept{point_origin_for}@<PO1, PO2::@\exposidnc{quantity-spec}@>) &&
  (@\exposidnc{is-derived-from-specialization-of}@<PO1, relative_point_origin>() ||
   @\exposidnc{is-derived-from-specialization-of}@<PO2, relative_point_origin>())
\end{codeblock}
\end{itemdescr}

\indexlibrarymemberexpos{operator==}{point-origin-interface}
\begin{itemdecl}
template<@\exposconceptnc{PointOrigin}@ PO1, @\exposconceptnc{PointOrigin}@ PO2>
friend consteval bool operator==(PO1 po1, PO2 po2);
\end{itemdecl}

\begin{itemdescr}
\pnum
\effects
Equivalent to:
\begin{codeblock}
if constexpr (@\exposidnc{is-derived-from-specialization-of}@<PO1, absolute_point_origin>() &&
              @\exposidnc{is-derived-from-specialization-of}@<PO2, absolute_point_origin>())
  return std::is_same_v<PO1, PO2> ||
         (@\exposidnc{is-derived-from-specialization-of}@<PO1, natural_point_origin_>() &&
          @\exposidnc{is-derived-from-specialization-of}@<PO2, natural_point_origin_>() &&
          (@\exposid{implicitly-convertible}@(@\exposidnc{remove-point}@(po1.@\exposidnc{quantity-spec}@),
                                  @\exposidnc{remove-point}@(po2.@\exposidnc{quantity-spec}@)) ||
           @\exposid{implicitly-convertible}@(@\exposidnc{remove-point}@(po2.@\exposidnc{quantity-spec}@),
                                  @\exposidnc{remove-point}@(po1.@\exposidnc{quantity-spec}@))));
else if constexpr (@\exposidnc{is-derived-from-specialization-of}@<PO1, relative_point_origin>() &&
                   @\exposidnc{is-derived-from-specialization-of}@<PO2, relative_point_origin>())
  return PO1::@\exposidnc{quantity-point}@ == PO2::@\exposidnc{quantity-point}@;
else if constexpr (@\exposidnc{is-derived-from-specialization-of}@<PO1, relative_point_origin>())
  return @\exposidnc{same-absolute-point-origins}@(po1, po2) &&
         is_eq_zero(PO1::@\exposidnc{quantity-point}@.quantity_from(
           PO1::@\exposidnc{quantity-point}@.absolute_point_origin));
else if constexpr (@\exposidnc{is-derived-from-specialization-of}@<PO2, relative_point_origin>())
  return @\exposidnc{same-absolute-point-origins}@(po1, po2) &&
         is_eq_zero(PO2::@\exposidnc{quantity-point}@.quantity_from(
           PO2::@\exposidnc{quantity-point}@.absolute_point_origin));
\end{codeblock}
\end{itemdescr}

\indexlibrarymemberexpos{operator-}{point-origin-interface}
\indexlibrarymember{operator-}{quantity}
\begin{itemdecl}
template<@\libconcept{point_quantity}@ QP, @\exposconceptnc{PointOrigin}@ PO2>
  requires @\libconcept{quantity_of}@<QP, PO2{}> &&
           @\libconcept{quantity_reference_of}@<decltype(auto(QP::reference)), get_quantity_spec(PO2{})>
friend constexpr @\libconcept{delta_quantity}@ auto operator-(const QP& qp, PO2 po);
template<@\exposconceptnc{PointOrigin}@ PO2, @\libconcept{point_quantity}@ QP>
  requires @\libconcept{quantity_of}@<QP, PO2{}> &&
           @\libconcept{quantity_reference_of}@<decltype(auto(QP::reference)), get_quantity_spec(PO2{})>
friend constexpr @\libconcept{delta_quantity}@ auto operator-(PO2 po, const QP& qp);
\end{itemdecl}

\begin{itemdescr}
\pnum
Let \tcode{d} be
\tcode{typename QP::quantity_type\{qp.\exposidnc{numerical-value}, QP::reference.\exposidnc{delta-reference}\}},
the delta of \tcode{qp} from \tcode{qp.point_origin}.

\pnum
\effects
For the first signature,
equivalent to:
\begin{codeblock}
if constexpr (QP::point_origin == po)
  return d;
else if constexpr (@\exposidnc{is-derived-from-specialization-of}@<PO2, absolute_point_origin>()) {
  return d + (qp.point_origin - qp.absolute_point_origin);
} else if constexpr (QP::point_origin == po.@\exposidnc{quantity-point}@.point_origin) {
  return d - po.@\exposidnc{quantity-point}@.quantity_from(po.@\exposidnc{quantity-point}@.point_origin);
} else {
  return d - po.@\exposidnc{quantity-point}@.quantity_from(po.@\exposidnc{quantity-point}@.point_origin) +
         (qp.point_origin - po.@\exposidnc{quantity-point}@.point_origin);
}
\end{codeblock}
For the second signature,
equivalent to: \tcode{return -(qp - po);}
\end{itemdescr}

\rSec3[qty.pt.orig.utils]{Utilities}

\rSec4[qty.same.abs.pt.origs]{Same absolute}

\indexlibrarymemberexpos{point_origin}{same-absolute-point-origins}
\begin{itemdecl}
template<@\exposconceptnc{PointOrigin}@ PO1, @\exposconceptnc{PointOrigin}@ PO2>
consteval bool @\exposidnc{same-absolute-point-origins}@(PO1 po1, PO2 po2);  // \expos
\end{itemdecl}

\begin{itemdescr}
\pnum
\effects
Equivalent to:
\begin{codeblock}
if constexpr (@\exposidnc{is-derived-from-specialization-of}@<PO1, absolute_point_origin>() &&
              @\exposidnc{is-derived-from-specialization-of}@<PO2, absolute_point_origin>())
  return po1 == po2;
else if constexpr (@\exposidnc{is-derived-from-specialization-of}@<PO1, relative_point_origin>() &&
                   @\exposidnc{is-derived-from-specialization-of}@<PO2, relative_point_origin>())
  return po1.@\exposidnc{absolute-point-origin}@ == po2.@\exposidnc{absolute-point-origin}@;
else if constexpr (@\exposidnc{is-derived-from-specialization-of}@<PO1, relative_point_origin>())
  return po1.@\exposidnc{absolute-point-origin}@ == po2;
else if constexpr (@\exposidnc{is-derived-from-specialization-of}@<PO2, relative_point_origin>())
  return po1 == po2.@\exposidnc{absolute-point-origin}@;
else
  return false;
\end{codeblock}
\end{itemdescr}

\indexlibrarymemberexpos{point_origin}{get-absolute-point-origin}
\begin{itemdecl}
template<@\exposconceptnc{PointOrigin}@ PO>
consteval @\exposconceptnc{PointOrigin}@ auto @\exposidnc{get-absolute-point-origin}@(PO po);  // \expos
\end{itemdecl}

\begin{itemdescr}
\pnum
\returns
\tcode{po.\exposidnc{absolute-point-origin}}
if \tcode{\exposidnc{is-derived-from-specialization-of}<PO, relative_point_origin>()} is \tcode{true},
and \tcode{po} otherwise.
\end{itemdescr}

\rSec4[qty.def.pt.orig]{Default}

\indexlibraryglobal{default_point_origin}
\begin{itemdecl}
template<@\exposconceptnc{BareReference}@ R>
consteval @\libconcept{point_origin_for}@<get_quantity_spec(R{})> auto default_point_origin(R);
\end{itemdecl}

\begin{itemdescr}
\pnum
\effects
Equivalent to:
\begin{codeblock}
if constexpr (requires { get_unit(R{}).@\exposidnc{point-origin}@; })
  return get_unit(R{}).@\exposidnc{point-origin}@;
else
  return natural_point_origin<get_quantity_spec(R{})>;
\end{codeblock}
\end{itemdescr}

\rSec2[qty.point]{Point}

\pnum
The variable template \tcode{point} names
the quantity specification of a point, \tcode{point<QS>} and \tcode{point<QS, PO>}\iref{point.qty},
and, for a unit \tcode{U}, the reference of a point measured from an origin,
\tcode{point<U, PO>} and \tcode{point<U>}.
A point quantity specification bound to a unit, \tcode{point<QS>[U]} or \tcode{point<QS, PO>[U]}\iref{qty.spec.ops},
is also the reference of a point.
A \defnadj{point}{reference} is a specialization of \tcode{point_}
that models \exposconcept{PointReference}\iref{qty.ref.concepts}.
\begin{example}
\begin{codeblock}
using namespace si::unit_symbols;
inline constexpr struct mean_sea_level final : absolute_point_origin<isq::altitude> {} mean_sea_level;

quantity<point<si::metre>> q1 = point<si::metre>(42);
quantity q2 = point<isq::altitude, mean_sea_level>[m](42);
static_assert(@\libconcept{quantity_of}@<decltype(q2), point<isq::altitude, mean_sea_level>>);
static_assert(!@\libconcept{quantity_of}@<decltype(q2), point<isq::altitude>>);
\end{codeblock}
\end{example}

\indexlibraryglobal{point}
\begin{itemdecl}
template<auto Arg, auto... Rest>
  requires @\seebelownc@
constexpr auto point = @\seebelownc@;
\end{itemdecl}

\begin{itemdescr}
\pnum
The value is
\begin{itemize}
\item
\tcode{point_<decltype(Arg)>\{\}}, the quantity specification of a point\iref{point.qty}
measured from \tcode{natural_point_origin<Arg>},
if \tcode{decltype(Arg)} models \libconcept{quantity_spec} and \tcode{sizeof...(Rest) == 0} is \tcode{true};
\item
otherwise, \tcode{\exposid{normalized-point-spec}(Arg, Rest...)}\iref{point.qty},
the quantity specification of a point measured from \tcode{Rest...},
if \tcode{decltype(Arg)} models \libconcept{quantity_spec};
\item
otherwise, \tcode{\exposid{make-point-reference}(get_quantity_spec(Arg), Arg, Rest...)},
the reference of a point of the kind associated with the unit \tcode{Arg}, in \tcode{Arg},
measured from \tcode{Rest...} if it is given and from \tcode{default_point_origin(Arg)} otherwise.
\end{itemize}

\pnum
\remarks
The expression in the \fakegrammarterm{requires-clause} is equivalent to:
\begin{codeblock}
(sizeof...(Rest) <= 1) &&
((@\libconcept{quantity_spec}@<decltype(Arg)> && (!@\exposconceptnc{PointQuantitySpec}@<decltype(Arg)>) &&
  (!@\exposconceptnc{DeltaQuantitySpec}@<decltype(Arg)>) && (... && @\libconcept{point_origin_for}@<decltype(Rest), Arg>)) ||
 (@\libconcept{unit}@<decltype(Arg)> && (... && @\libconcept{point_origin_for}@<decltype(Rest), get_quantity_spec(Arg)>)))
\end{codeblock}

\pnum
\begin{note}
\tcode{point<>} takes a unit or a quantity specification, and nothing else.
The test is on the value of \tcode{Arg}:
\tcode{kind_of<isq::length>[m]} is the unit \tcode{m}, so \tcode{point<kind_of<isq::length>[m]>} is \tcode{point<m>},
while \tcode{point<isq::height[m]>}, \tcode{point<isq::altitude[m], mean_sea_level>},
\tcode{point<delta<isq::height>>}, and \tcode{point<point<isq::height>>} are ill-formed.
\end{note}
\end{itemdescr}

\indexlibrarymemberexpos{quantity_spec}{make-point-reference}
\begin{itemdecl}
template<@\libconcept{quantity_spec}@ QS, @\libconcept{unit}@ U, @\exposconceptnc{PointOrigin}@... PO>
  requires(sizeof...(PO) <= 1) && (... && @\libconcept{point_origin_for}@<PO, QS{}>)
consteval @\libconcept{quantity_reference}@ auto @\exposidnc{make-point-reference}@(QS, U, PO...);  // \expos
\end{itemdecl}

\begin{itemdescr}
\pnum
Let \tcode{r} be \tcode{\exposid{make-reference}(QS\{\}, U\{\})}\iref{qty.spec.utils}.

\pnum
\returns
\tcode{point_<decltype(r), PO...>\{\}} if \tcode{sizeof...(PO) == 1} is \tcode{true},
and \tcode{point_<decltype(r), decltype(default_point_origin(r))>\{\}} otherwise.
\end{itemdescr}

\indexlibrarymemberexpos{quantity_spec}{make-point-spec-reference}
\begin{itemdecl}
template<@\exposconceptnc{PointQuantitySpec}@ QS, @\libconcept{unit}@ U>
consteval @\libconcept{quantity_reference}@ auto @\exposidnc{make-point-spec-reference}@(QS, U);  // \expos
\end{itemdecl}

\begin{itemdescr}
\pnum
Let \tcode{S} be \tcode{QS::\exposid{quantity-spec}} and
\tcode{O} be \tcode{\exposid{point-spec-origin}(QS\{\})}\iref{point.qty}.

\pnum
\returns
\tcode{\exposid{make-point-reference}(S, U\{\})} if \tcode{QS} is \tcode{point_<decltype(S)>},
and \tcode{\exposid{make-point-reference}(S, U\{\}, O)} otherwise.
\begin{note}
\tcode{point<QS>[U]} is measured from \tcode{default_point_origin(QS[U])}, which compares equal to
\tcode{natural_point_origin<QS>}\iref{qty.spec.ops}.
\end{note}
\end{itemdescr}

\indexlibraryglobal{point_}
\begin{codeblock}
namespace std {

template<typename T>
concept @\defexposconceptnc{DefaultOriginPointReference}@ =                       // \expos
  @\exposconceptnc{PointReference}@<T> &&
  (T::@\exposidnc{point-origin}@ == default_point_origin(T::@\exposidnc{delta-reference}@));

template<@\exposconceptnc{BareReference}@ R, @\exposconceptnc{PointOrigin}@ PO>
struct point_<R, PO> {
  static constexpr R @\exposidnc{delta-reference}@{};  // \expos
  static constexpr PO @\exposidnc{point-origin}@{};     // \expos

  friend consteval @\libconcept{quantity_spec}@ auto get_quantity_spec(point_);
  friend consteval @\libconcept{unit}@ auto get_unit(point_);

  template<typename R2, typename PO2>
  friend consteval bool operator==(point_, point_<R2, PO2>);
  template<@\exposconceptnc{BareReference}@ R2>
  friend consteval bool operator==(point_, R2);

  template<typename FwdRep,
           @\libconcept{representation_of}@<get_quantity_spec(R{})> Rep = std::remove_cvref_t<FwdRep>>
  constexpr quantity<point_{}, Rep> operator()(FwdRep&& value) const;
};

}
\end{codeblock}

\pnum
\tcode{point_<R, PO>} represents the reference of a point
measured from the origin \tcode{PO}.
\tcode{R} is the reference of the delta from \tcode{PO}.

\indexlibrarymember{get_quantity_spec}{point_}
\begin{itemdecl}
friend consteval @\libconcept{quantity_spec}@ auto get_quantity_spec(point_);
\end{itemdecl}

\begin{itemdescr}
\pnum
\returns
\tcode{\exposid{normalized-point-spec}(get_quantity_spec(R\{\}), PO\{\})}\iref{point.qty}.

\pnum
\begin{note}
The quantity specification of a point reference names its origin
whenever that is not the natural point origin of its scale, also when the unit brings it.
A point and a delta never share a quantity specification.
\end{note}
\begin{example}
\tcode{get_quantity_spec(point<isq::altitude>[m])} is \tcode{point<isq::altitude>},
\tcode{get_quantity_spec(point<isq::altitude, mean_sea_level>[m])} is
\tcode{point<isq::altitude, mean_sea_level>}, and
\tcode{get_quantity_spec(point<deg_C>)} is
\tcode{point<kind_of<isq::thermodynamic_temperature>, si::ice_point>}.
\end{example}
\end{itemdescr}

\indexlibrarymember{get_unit}{point_}
\begin{itemdecl}
friend consteval @\libconcept{unit}@ auto get_unit(point_);
\end{itemdecl}

\begin{itemdescr}
\pnum
\returns
\tcode{get_unit(R\{\})}.
\end{itemdescr}

\indexlibrarymember{operator==}{point_}
\begin{itemdecl}
template<typename R2, typename PO2>
friend consteval bool operator==(point_, point_<R2, PO2>);
template<@\exposconceptnc{BareReference}@ R2>
friend consteval bool operator==(point_, R2);
\end{itemdecl}

\begin{itemdescr}
\pnum
\returns
\tcode{R\{\} == R2\{\} \&\& PO\{\} == PO2\{\}} for the first signature, and
\tcode{false} for the second signature.
\end{itemdescr}

\indexlibrarymember{operator()}{point_}
\begin{itemdecl}
template<typename FwdRep,
         @\libconcept{representation_of}@<get_quantity_spec(R{})> Rep = std::remove_cvref_t<FwdRep>>
constexpr quantity<point_{}, Rep> operator()(FwdRep&& value) const;
\end{itemdecl}

\begin{itemdescr}
\pnum
\effects
Equivalent to:
\tcode{return quantity<point_\{\}, Rep>\{std::forward<FwdRep>(value), *this\};}
\end{itemdescr}

\rSec2[qty.like]{Interoperability}

\pnum
The interfaces specified in this subclause
are used by \tcode{quantity}
to specify conversions with other types representing deltas and points.
\begin{note}
\ref{qty.chrono} implements them for \tcode{std::chrono::duration} and \tcode{std::chrono::time_point}.
\end{note}

\begin{itemdecl}
template<typename T>
concept @\defexposconceptnc{quantity-like}@ =                                       // \expos
  (!@\exposidnc{is-specialization-of}@<std::remove_cv_t<T>, quantity>()) &&
  requires(const T& qty, const quantity_like_traits<T>::rep& num) {
    { quantity_like_traits<T>::to_numerical_value(qty) } ->
      std::@\stdconcept{same_as}@<typename quantity_like_traits<T>::rep>;
    { quantity_like_traits<T>::from_numerical_value(num) } -> std::@\stdconcept{same_as}@<T>;
    requires std::same_as<decltype(quantity_like_traits<T>::explicit_import), const bool>;
    requires std::same_as<decltype(quantity_like_traits<T>::explicit_export), const bool>;
    typename std::bool_constant<quantity_like_traits<T>::explicit_import>;
    typename std::bool_constant<quantity_like_traits<T>::explicit_export>;
    typename quantity<quantity_like_traits<T>::reference, typename quantity_like_traits<T>::rep>;
  };
\end{itemdecl}

\pnum
In the following descriptions, let
\begin{itemize}
\item
\tcode{Traits} be \tcode{quantity_like_traits},
\item
\tcode{Q} be a type for which \tcode{Traits<Q>} is specialized,
\item
\tcode{qty} be an lvalue of type \tcode{const Q}, and
\item
\tcode{num} be an lvalue of type \tcode{const Traits<Q>::rep}.
\end{itemize}

\pnum
\tcode{Q} models \exposconceptnc{quantity-like} if and only if:
\begin{itemize}
\item
\tcode{Traits<Q>::to_numerical_value(qty)} returns the numerical value\irefiev{112-01-29} of \tcode{qty}.
\item
\tcode{Traits<Q>::from_numerical_value(num)} returns a \tcode{Q} with numerical value \tcode{num}.
\item
If \tcode{decltype(auto(Traits<Q>::reference))} models \exposconcept{PointReference},
both numerical values are offset from \tcode{Traits<Q>::reference.\exposidnc{point-origin}}.
\end{itemize}
\begin{note}
\tcode{Q} corresponds to a point if \tcode{Traits<Q>::reference} is a point reference\iref{qty.point},
and to a delta otherwise.
\end{note}

\pnum
If the following expression is \tcode{true}, the specified conversion will be explicit.
\begin{itemize}
\item
\tcode{Traits<Q>::explicit_import} for the conversion from \tcode{Q} to this library's type.
\item
\tcode{Traits<Q>::explicit_export} for the conversion from this library's type to \tcode{Q}.
\end{itemize}

\rSec2[qty.syn]{Class template \tcode{quantity}}

\begin{codeblock}
namespace std {

template<typename T>
concept @\deflibconcept{delta_quantity}@ =
  (@\exposidnc{is-specialization-of}@<std::remove_cv_t<T>, quantity>()) &&     // \expos
  (!@\exposconceptnc{PointReference}@<decltype(auto(T::reference))>);

template<typename T>
concept @\deflibconcept{point_quantity}@ =
  (@\exposidnc{is-specialization-of}@<std::remove_cv_t<T>, quantity>()) &&     // \expos
  @\exposconceptnc{PointReference}@<decltype(auto(T::reference))>;

template<typename Q, auto V>
concept @\deflibconcept{quantity_of}@ =
  (@\libconcept{delta_quantity}@<Q> || @\libconcept{point_quantity}@<Q>) &&
  ((@\libconcept{quantity_spec}@<decltype(V)> && @\libconcept{quantity_spec_of}@<decltype(auto(Q::quantity_spec)), V>) ||
   (@\libconcept{point_quantity}@<Q> &&
    @\exposconceptnc{SameAbsolutePointOriginAs}@<decltype(auto(Q::absolute_point_origin)), V>) ||
   (std::@\stdconcept{same_as}@<quantity_character, decltype(V)> && (Q::quantity_spec.character == V) &&
    @\libconcept{representation_of}@<typename Q::rep, V>));

template<@\libconcept{unit}@ UFrom, @\libconcept{unit}@ UTo>
consteval bool @\exposidnc{integral-conversion-factor}@(UFrom from, UTo to)              // \expos
{
  return @\exposidnc{is-integral}@(@\exposidnc{get-canonical-unit}@(from).mag / @\exposidnc{get-canonical-unit}@(to).mag);
}

template<typename T>
concept @\defexposconceptnc{IsFloatingPoint}@ = treat_as_floating_point<T>;                      // \expos

template<typename FromRep, typename ToRep, auto FromUnit = one, auto ToUnit = one>
concept @\defexposconceptnc{ValuePreservingTo}@ =                                                // \expos
  @\libconcept{Representation}@<std::remove_cvref_t<FromRep>> && @\libconcept{Representation}@<ToRep> &&
  @\libconcept{unit}@<decltype(FromUnit)> && @\libconcept{unit}@<decltype(ToUnit)> && std::@\stdconcept{assignable_from}@<ToRep&, FromRep> &&
  (@\exposidnc{IsFloatingPoint}@<ToRep> || (!@\exposconceptnc{IsFloatingPoint}@<std::remove_cvref_t<FromRep>> &&
                              (@\exposidnc{integral-conversion-factor}@(FromUnit, ToUnit))));

template<typename QFrom, typename QTo>
concept @\defexposconceptnc{QuantityConvertibleTo}@ =                                            // \expos
  @\libconcept{delta_quantity}@<QFrom> && @\libconcept{delta_quantity}@<QTo> &&
  @\exposconceptnc{QuantitySpecConvertibleTo}@<QFrom::quantity_spec, QTo::quantity_spec> &&
  @\exposconceptnc{UnitConvertibleTo}@<QFrom::unit, QTo::unit> &&
  @\exposidnc{ValuePreservingTo}@<typename QFrom::rep, typename QTo::rep, QFrom::unit, QTo::unit> &&
  requires(QFrom q) { @\exposidnc{sudo-cast}@<QTo>(q); };                                // see \ref{qty.non.mem.conv}

template<auto QS, typename Func, typename T, typename U>
concept @\defexposconceptnc{InvokeResultOf}@ =                                                   // \expos
  @\libconcept{quantity_spec}@<decltype(QS)> && std::@\stdconcept{regular_invocable}@<Func, T, U> &&
  @\libconcept{representation_of}@<std::invoke_result_t<Func, T, U>, QS>;

template<typename Func, typename Q1, typename Q2,
         auto QS = std::invoke_result_t<Func, decltype(auto(Q1::quantity_spec)),
                                        decltype(auto(Q2::quantity_spec))>{}>
concept @\defexposconceptnc{InvocableQuantities}@ =                                              // \expos
  @\libconcept{quantity_spec}@<decltype(QS)> && @\libconcept{delta_quantity}@<Q1> && @\libconcept{delta_quantity}@<Q2> &&
  @\exposconceptnc{InvokeResultOf}@<QS, Func, typename Q1::rep, typename Q2::rep>;

template<auto R1, auto R2>
concept @\defexposconceptnc{HaveCommonReference}@ = requires { get_common_reference(R1, R2); };  // \expos

template<typename Func, @\libconcept{delta_quantity}@ Q1, @\libconcept{delta_quantity}@ Q2>
using @\exposidnc{common-quantity-for}@ =                                                // \expos
  quantity<get_common_reference(Q1::reference, Q2::reference),
           std::invoke_result_t<Func, typename Q1::rep, typename Q2::rep>>;

template<typename Common, auto R, typename Rep>
  constexpr decltype(auto) @\exposidnc{to-common-quantity}@(const quantity<R, Rep>& q);         // \expos

template<typename Func, typename Q1, typename Q2>
concept @\defexposconceptnc{CommonlyInvocableQuantities}@ =                                      // \expos
  @\libconcept{delta_quantity}@<Q1> && @\libconcept{delta_quantity}@<Q2> &&
  @\exposconceptnc{HaveCommonReference}@<Q1::reference, Q2::reference> &&
  std::@\stdconcept{convertible_to}@<Q1, @\exposidnc{common-quantity-for}@<Func, Q1, Q2>> &&
  std::@\stdconcept{convertible_to}@<Q2, @\exposidnc{common-quantity-for}@<Func, Q1, Q2>> &&
  @\exposconceptnc{InvocableQuantities}@<Func, Q1, Q2,
                      get_common_quantity_spec(Q1::quantity_spec, Q2::quantity_spec)>;

template<auto R1, auto R2, typename Rep1, typename Rep2>
concept @\defexposconceptnc{SameValueAs}@ =                                                      // \expos
  (@\exposid{equivalent}@(get_unit(R1), get_unit(R2))) && std::@\stdconcept{convertible_to}@<Rep1, Rep2>;

template<typename R2, auto R>
concept @\defexposconceptnc{SameKindReference}@ =                                                // \expos
  @\libconcept{quantity_reference}@<R2> &&
  (@\exposconceptnc{PointReference}@<R2> == @\exposconceptnc{PointReference}@<decltype(R)>) &&
  (!@\exposconceptnc{PointReference}@<R2> || R2::@\exposidnc{point-origin}@ == R.@\exposidnc{point-origin}@);

template<typename T>
using @\exposidnc{quantity-like-type}@ =                                                 // \expos
  quantity<quantity_like_traits<T>::reference, typename quantity_like_traits<T>::rep>;

template<typename T, typename U, typename TT = std::remove_reference_t<T>>
concept @\defexposconceptnc{Mutable}@ = (!std::is_const_v<TT>) && std::@\stdconcept{same_as}@<std::remove_volatile_t<TT>, U>; // \expos

struct @\exposidnc{quantity-interface}@ {  // \expos
  // \ref{qty.arith.ops}, arithmetic operations

  template<auto R1, typename Rep1, auto R2, typename Rep2>
    requires @\exposconceptnc{CommonlyInvocableQuantities}@<std::plus<>, quantity<R1, Rep1>, quantity<R2, Rep2>>
  friend constexpr @\libconcept{delta_quantity}@ auto operator+(const quantity<R1, Rep1>& lhs,
                                                 const quantity<R2, Rep2>& rhs);
  template<auto R1, typename Rep1, auto R2, typename Rep2>
    requires @\exposconceptnc{CommonlyInvocableQuantities}@<std::minus<>, quantity<R1, Rep1>, quantity<R2, Rep2>>
  friend constexpr @\libconcept{delta_quantity}@ auto operator-(const quantity<R1, Rep1>& lhs,
                                                 const quantity<R2, Rep2>& rhs);
  template<auto R1, typename Rep1, auto R2, typename Rep2>
    requires(!treat_as_floating_point<Rep1>) && (!treat_as_floating_point<Rep2>) &&
            @\exposconceptnc{CommonlyInvocableQuantities}@<std::modulus<>, quantity<R1, Rep1>, quantity<R2, Rep2>>
  friend constexpr @\libconcept{delta_quantity}@ auto operator%(const quantity<R1, Rep1>& lhs,
                                                 const quantity<R2, Rep2>& rhs);

  template<auto R1, typename Rep1, @\libconcept{Representation}@ Value>
    requires(!@\exposconceptnc{PointReference}@<decltype(R1)>) && (quantity<R1, Rep1>::unit == std::one) &&
            @\exposconceptnc{InvokeResultOf}@<quantity<R1, Rep1>::quantity_spec, std::plus<>, Rep1, const Value&>
  friend constexpr @\libconcept{delta_quantity}@ auto operator+(const quantity<R1, Rep1>& lhs, const Value& rhs);
  template<auto R1, typename Rep1, @\libconcept{Representation}@ Value>
    requires(!@\exposconceptnc{PointReference}@<decltype(R1)>) && (quantity<R1, Rep1>::unit == std::one) &&
            @\exposconceptnc{InvokeResultOf}@<quantity<R1, Rep1>::quantity_spec, std::minus<>, Rep1, const Value&>
  friend constexpr @\libconcept{delta_quantity}@ auto operator-(const quantity<R1, Rep1>& lhs, const Value& rhs);
  template<auto R1, typename Rep1, @\libconcept{Representation}@ Value>
    requires(!@\exposconceptnc{PointReference}@<decltype(R1)>) && (quantity<R1, Rep1>::unit == std::one) &&
            @\exposconceptnc{InvokeResultOf}@<quantity<R1, Rep1>::quantity_spec, std::modulus<>, Rep1, const Value&>
  friend constexpr @\libconcept{delta_quantity}@ auto operator%(const quantity<R1, Rep1>& lhs, const Value& rhs);

  template<auto R1, typename Rep1, @\libconcept{Representation}@ Value>
    requires(!@\exposconceptnc{PointReference}@<decltype(R1)>) && (quantity<R1, Rep1>::unit == std::one) &&
            @\exposconceptnc{InvokeResultOf}@<quantity<R1, Rep1>::quantity_spec, std::plus<>, Rep1, const Value&>
  friend constexpr @\libconcept{delta_quantity}@ auto operator+(const Value& lhs, const quantity<R1, Rep1>& rhs);
  template<auto R1, typename Rep1, @\libconcept{Representation}@ Value>
    requires(!@\exposconceptnc{PointReference}@<decltype(R1)>) && (quantity<R1, Rep1>::unit == std::one) &&
            @\exposconceptnc{InvokeResultOf}@<quantity<R1, Rep1>::quantity_spec, std::minus<>, Rep1, const Value&>
  friend constexpr @\libconcept{delta_quantity}@ auto operator-(const Value& lhs, const quantity<R1, Rep1>& rhs);
  template<auto R1, typename Rep1, @\libconcept{Representation}@ Value>
    requires(!@\exposconceptnc{PointReference}@<decltype(R1)>) && (quantity<R1, Rep1>::unit == std::one) &&
            @\exposconceptnc{InvokeResultOf}@<quantity<R1, Rep1>::quantity_spec, std::modulus<>, Rep1, const Value&>
  friend constexpr @\libconcept{delta_quantity}@ auto operator%(const Value& lhs, const quantity<R1, Rep1>& rhs);

  template<auto R1, typename Rep1, auto R2, typename Rep2>
    requires @\exposconceptnc{InvocableQuantities}@<std::multiplies<>, quantity<R1, Rep1>, quantity<R2, Rep2>>
  friend constexpr @\libconcept{delta_quantity}@ auto operator*(const quantity<R1, Rep1>& lhs,
                                                 const quantity<R2, Rep2>& rhs);
  template<auto R1, typename Rep1, auto R2, typename Rep2>
    requires @\exposconceptnc{InvocableQuantities}@<std::divides<>, quantity<R1, Rep1>, quantity<R2, Rep2>>
  friend constexpr @\libconcept{delta_quantity}@ auto operator/(const quantity<R1, Rep1>& lhs,
                                                 const quantity<R2, Rep2>& rhs);

  template<auto R1, typename Rep1, typename Value>
    requires(!@\exposconceptnc{PointReference}@<decltype(R1)>) &&
            (!@\libconcept{delta_quantity}@<Value>) && (!@\libconcept{quantity_reference}@<Value>) &&
            @\exposconceptnc{InvokeResultOf}@<quantity<R1, Rep1>::quantity_spec, std::multiplies<>, Rep1, const Value&>
  friend constexpr @\libconcept{quantity_of}@<quantity<R1, Rep1>::quantity_spec> auto operator*(const quantity<R1, Rep1>& lhs,
                                                                                 const Value& rhs);
  template<auto R1, typename Rep1, typename Value>
    requires(!@\exposconceptnc{PointReference}@<decltype(R1)>) &&
            (!@\libconcept{delta_quantity}@<Value>) && (!@\libconcept{quantity_reference}@<Value>) &&
            @\exposconceptnc{InvokeResultOf}@<quantity<R1, Rep1>::quantity_spec, std::divides<>, Rep1, const Value&>
  friend constexpr @\libconcept{quantity_of}@<quantity<R1, Rep1>::quantity_spec> auto operator/(const quantity<R1, Rep1>& lhs,
                                                                                 const Value& rhs);

  template<typename Value, auto R1, typename Rep1>
    requires(!@\exposconceptnc{PointReference}@<decltype(R1)>) &&
            (!@\libconcept{delta_quantity}@<Value>) && (!@\libconcept{quantity_reference}@<Value>) &&
            @\exposconceptnc{InvokeResultOf}@<quantity<R1, Rep1>::quantity_spec, std::multiplies<>, const Value&, Rep1>
  friend constexpr @\libconcept{quantity_of}@<quantity<R1, Rep1>::quantity_spec> auto operator*(const Value& lhs,
                                                                                 const quantity<R1, Rep1>& rhs);
  template<typename Value, auto R1, typename Rep1>
    requires(!@\exposconceptnc{PointReference}@<decltype(R1)>) &&
            (!@\libconcept{delta_quantity}@<Value>) && (!@\libconcept{quantity_reference}@<Value>) &&
            @\exposconceptnc{ScalarOperand}@<decltype(auto(quantity<R1, Rep1>::quantity_spec))> &&
            @\exposconceptnc{InvokeResultOf}@<quantity<R1, Rep1>::quantity_spec, std::divides<>, const Value&, Rep1>
  friend constexpr @\libconcept{delta_quantity}@ auto operator/(const Value& lhs, const quantity<R1, Rep1>& rhs);

  template<auto R1, typename Rep1, @\libconcept{delta_quantity}@ Q2>
    requires @\seebelownc@
  friend constexpr @\libconcept{point_quantity}@ auto operator+(const quantity<R1, Rep1>& qp, const Q2& q);
  template<auto R1, typename Rep1, @\libconcept{delta_quantity}@ Q2>
    requires @\seebelownc@
  friend constexpr @\libconcept{point_quantity}@ auto operator+(const Q2& q, const quantity<R1, Rep1>& qp);
  template<auto R1, typename Rep1, @\libconcept{delta_quantity}@ Q2>
    requires @\seebelownc@
  friend constexpr @\libconcept{point_quantity}@ auto operator-(const quantity<R1, Rep1>& qp, const Q2& q);

  template<auto R1, typename Rep1, @\libconcept{point_quantity}@ QP2>
    requires @\seebelownc@
  friend constexpr @\libconcept{delta_quantity}@ auto operator-(const quantity<R1, Rep1>& lhs, const QP2& rhs);


  // \ref{qty.cmp}, comparison

  template<auto R1, typename Rep1, auto R2, typename Rep2>
    requires @\seebelownc@
  friend constexpr bool operator==(const quantity<R1, Rep1>& lhs,
                                   const quantity<R2, Rep2>& rhs);
  template<auto R1, typename Rep1, auto R2, typename Rep2>
    requires @\seebelownc@
  friend constexpr auto operator<=>(const quantity<R1, Rep1>& lhs,
                                    const quantity<R2, Rep2>& rhs);

  template<auto R1, typename Rep1, @\libconcept{Representation}@ Value>
    requires @\seebelownc@
  friend constexpr bool operator==(const quantity<R1, Rep1>& lhs, const Value& rhs);
  template<auto R1, typename Rep1, @\libconcept{Representation}@ Value>
    requires @\seebelownc@
  friend constexpr auto operator<=>(const quantity<R1, Rep1>& lhs, const Value& rhs);

  template<auto R1, typename Rep1, @\libconcept{point_quantity}@ QP2>
    requires @\seebelownc@
  friend constexpr bool operator==(const quantity<R1, Rep1>& lhs, const QP2& rhs);
  template<auto R1, typename Rep1, @\libconcept{point_quantity}@ QP2>
    requires @\seebelownc@
  friend constexpr auto operator<=>(const quantity<R1, Rep1>& lhs, const QP2& rhs);
};

template<@\libconcept{quantity_reference}@ auto R, @\libconcept{representation_of}@<get_quantity_spec(R)> Rep = double>
class @\libglobal{quantity}@ : public @\exposidnc{quantity-interface}@ {
public:
  Rep @\exposidnc{numerical-value}@;  // \expos

  // member types and values
  static constexpr @\libconcept{quantity_reference}@ auto reference = R;
  static constexpr std::@\libconcept{quantity_spec}@ auto quantity_spec = get_quantity_spec(reference);
  static constexpr std::@\libconcept{dimension}@ auto dimension = quantity_spec.dimension;
  static constexpr std::@\libconcept{unit}@ auto unit = get_unit(reference);
  using rep = Rep;

  // members of a point, present only if \tcode{decltype(R)} models \exposconcept{PointReference}
  static constexpr @\exposconceptnc{PointOrigin}@ auto point_origin = R.@\exposidnc{point-origin}@;
  static constexpr @\exposconceptnc{PointOrigin}@ auto absolute_point_origin = @\seebelownc@;
  using quantity_type = quantity<R.@\exposidnc{delta-reference}@, Rep>;

  // \ref{qty.static}, static member functions
  static constexpr quantity zero() noexcept
    requires @\seebelownc@;
  static constexpr quantity one() noexcept
    requires @\seebelownc@;
  static constexpr quantity min() noexcept
    requires @\seebelownc@;
  static constexpr quantity max() noexcept
    requires @\seebelownc@;

  // \ref{qty.cons}, constructors and assignment

  quantity() = default;
  quantity(const quantity&) = default;
  quantity(quantity&&) = default;
  ~quantity() = default;

  template<typename FwdValue, @\libconcept{quantity_reference}@ R2>
    requires @\exposconceptnc{SameValueAs}@<R2{}, R, std::remove_cvref_t<FwdValue>, Rep> &&
             @\exposconceptnc{SameKindReference}@<R2, R>
  constexpr quantity(FwdValue&& v, R2);

  template<typename FwdValue, @\libconcept{quantity_reference}@ R2, typename Value = std::remove_cvref_t<FwdValue>>
    requires(!@\exposconceptnc{SameValueAs}@<R2{}, R, Value, Rep>) &&
            @\exposconceptnc{QuantityConvertibleTo}@<quantity<R2{}, Value>, quantity>
  constexpr quantity(FwdValue&& v, R2);

  template<@\exposconceptnc{ValuePreservingTo}@<Rep> FwdValue>
    requires(!@\exposconceptnc{PointReference}@<decltype(R)>) && (unit == std::one)
  constexpr quantity(FwdValue&& v);

  template<@\exposconceptnc{QuantityConvertibleTo}@<quantity> Q>
  constexpr explicit(@\seebelownc@) quantity(const Q& q);

  template<@\exposconceptnc{quantity-like}@ Q>
    requires @\seebelownc@
  constexpr explicit(@\seebelownc@) quantity(const Q& q);

  template<typename FwdQ, typename Q = std::remove_cvref_t<FwdQ>>
    requires @\seebelownc@
  constexpr explicit quantity(FwdQ&& q);

  template<typename FwdQ, @\exposconceptnc{PointOrigin}@ PO2, typename Q = std::remove_cvref_t<FwdQ>>
    requires @\seebelownc@
  constexpr quantity(FwdQ&& q, PO2);

  template<auto R2, typename Rep2>
    requires @\seebelownc@
  constexpr explicit(@\seebelownc@) quantity(const quantity<R2, Rep2>& qp);

  quantity& operator=(const quantity&) = default;
  quantity& operator=(quantity&&) = default;

  template<@\exposconceptnc{ValuePreservingTo}@<Rep> FwdValue>
    requires(!@\exposconceptnc{PointReference}@<decltype(R)>) && (unit == std::one)
  constexpr quantity& operator=(FwdValue&& v);

  // \ref{qty.conv}, conversions

  template<@\exposconceptnc{UnitCompatibleWith}@<unit, quantity_spec> ToU>
    requires @\seebelownc@
  constexpr @\libconcept{quantity_of}@<quantity_spec> auto in(ToU) const;

  template<@\libconcept{representation_of}@<quantity_spec> ToRep>
    requires @\seebelownc@
  constexpr @\libconcept{quantity_of}@<quantity_spec> auto in() const;

  template<@\libconcept{representation_of}@<quantity_spec> ToRep,
           @\exposconceptnc{UnitCompatibleWith}@<unit, quantity_spec> ToU>
    requires @\seebelownc@
  constexpr @\libconcept{quantity_of}@<quantity_spec> auto in(ToU) const;

  template<@\exposconceptnc{UnitCompatibleWith}@<unit, quantity_spec> ToU>
    requires @\seebelownc@
  constexpr @\libconcept{quantity_of}@<quantity_spec> auto force_in(ToU) const;

  template<@\libconcept{representation_of}@<quantity_spec> ToRep>
    requires @\seebelownc@
  constexpr @\libconcept{quantity_of}@<quantity_spec> auto force_in() const;

  template<@\libconcept{representation_of}@<quantity_spec> ToRep,
           @\exposconceptnc{UnitCompatibleWith}@<unit, quantity_spec> ToU>
    requires @\seebelownc@
  constexpr @\libconcept{quantity_of}@<quantity_spec> auto force_in(ToU) const;

  template<@\exposconceptnc{PointOrigin}@ NewPO>
    requires @\exposconceptnc{PointReference}@<decltype(R)> &&
             @\exposconceptnc{SameAbsolutePointOriginAs}@<NewPO, absolute_point_origin>
  constexpr @\libconcept{quantity_of}@<(NewPO{})> auto point_for(NewPO new_origin) const;

  constexpr @\libconcept{delta_quantity}@ auto magnitude() const
    requires @\seebelownc@;

  // \ref{qty.obs}, observers

  template<std::@\libconcept{unit}@ U>
    requires(!@\exposconceptnc{PointReference}@<decltype(R)>) && (@\exposid{equivalent}@(U{}, unit))
  constexpr rep& numerical_value_ref_in(U) & noexcept;
  template<std::@\libconcept{unit}@ U>
    requires(!@\exposconceptnc{PointReference}@<decltype(R)> ||
             @\exposconceptnc{DefaultOriginPointReference}@<decltype(R)>) && (@\exposid{equivalent}@(U{}, unit))
  constexpr const rep& numerical_value_ref_in(U) const & noexcept;
  template<std::@\libconcept{unit}@ U>
    requires(@\exposid{equivalent}@(U{}, unit))
  void numerical_value_ref_in(U) const && = delete;

  template<@\exposconceptnc{UnitCompatibleWith}@<unit, quantity_spec> U>
    requires @\seebelownc@
  constexpr rep numerical_value_in(U) const noexcept;

  template<@\exposconceptnc{UnitCompatibleWith}@<unit, quantity_spec> U>
    requires @\seebelownc@
  constexpr rep force_numerical_value_in(U) const noexcept;

  template<@\libconcept{representation_of}@<quantity_spec> ToRep, @\exposconceptnc{UnitCompatibleWith}@<unit, quantity_spec> U>
    requires @\seebelownc@
  constexpr ToRep numerical_value_in(U) const;

  template<@\libconcept{representation_of}@<quantity_spec> ToRep, @\exposconceptnc{UnitCompatibleWith}@<unit, quantity_spec> U>
    requires @\seebelownc@
  constexpr ToRep force_numerical_value_in(U) const;

  template<@\exposconceptnc{PointOrigin}@ PO2>
    requires @\exposconceptnc{PointReference}@<decltype(R)> && requires(const quantity qp) { qp - PO2{}; }
  constexpr @\libconcept{delta_quantity}@ auto quantity_from(PO2) const &;
  template<@\exposconceptnc{PointOrigin}@ PO2>
    requires @\exposconceptnc{PointReference}@<decltype(R)> && requires(const quantity qp) { qp - PO2{}; }
  constexpr @\libconcept{delta_quantity}@ auto quantity_from(PO2) &&;

  template<typename QP>
    requires @\exposconceptnc{PointReference}@<decltype(R)> && @\libconcept{quantity_of}@<QP, absolute_point_origin>
  constexpr @\libconcept{delta_quantity}@ auto quantity_from(const QP& rhs) const;

  constexpr @\libconcept{delta_quantity}@ auto quantity_from_zero() const
    requires @\exposconceptnc{DefaultOriginPointReference}@<decltype(R)>;

  // \ref{qty.conv.ops}, conversion operations

  template<typename V_, std::@\stdconcept{constructible_from}@<Rep> Value = std::remove_cvref_t<V_>>
    requires(!@\exposconceptnc{PointReference}@<decltype(R)>) && (unit == std::one)
  explicit operator V_() const & noexcept;

  template<typename Q_, @\exposconceptnc{quantity-like}@ Q = std::remove_cvref_t<Q_>>
    requires @\seebelownc@
  constexpr explicit(@\seebelownc@) operator Q_() const noexcept(@\seebelownc@);

  // \ref{qty.unary.ops}, unary operations

  constexpr @\libconcept{quantity_of}@<quantity_spec> auto operator+() const
    requires @\seebelownc@;
  constexpr @\libconcept{quantity_of}@<quantity_spec> auto operator-() const
    requires @\seebelownc@;

  template<@\exposconceptnc{Mutable}@<quantity> Q>
  friend constexpr decltype(auto) operator++(Q&& q)
    requires @\seebelownc@;
  template<@\exposconceptnc{Mutable}@<quantity> Q>
  friend constexpr decltype(auto) operator--(Q&& q)
    requires @\seebelownc@;

  constexpr @\libconcept{quantity_of}@<quantity_spec> auto operator++(int)
    requires @\seebelownc@;
  constexpr @\libconcept{quantity_of}@<quantity_spec> auto operator--(int)
    requires @\seebelownc@;

  // \ref{qty.assign.ops}, compound assignment operations

  template<@\exposconceptnc{Mutable}@<quantity> Q, auto R2, typename Rep2>
    requires @\seebelownc@
  friend constexpr decltype(auto) operator+=(Q&& lhs, const quantity<R2, Rep2>& rhs);
  template<@\exposconceptnc{Mutable}@<quantity> Q, auto R2, typename Rep2>
    requires @\seebelownc@
  friend constexpr decltype(auto) operator-=(Q&& lhs, const quantity<R2, Rep2>& rhs);
  template<@\exposconceptnc{Mutable}@<quantity> Q, auto R2, typename Rep2>
    requires @\seebelownc@
  friend constexpr decltype(auto) operator%=(Q&& lhs, const quantity<R2, Rep2>& rhs);

  template<@\exposconceptnc{Mutable}@<quantity> Q, @\exposconceptnc{ValuePreservingTo}@<Rep> Value>
    requires @\seebelownc@
  friend constexpr decltype(auto) operator*=(Q&& lhs, const Value& rhs);
  template<@\exposconceptnc{Mutable}@<quantity> Q, @\exposconceptnc{ValuePreservingTo}@<Rep> Value>
    requires @\seebelownc@
  friend constexpr decltype(auto) operator/=(Q&& lhs, const Value& rhs);

  template<@\exposconceptnc{Mutable}@<quantity> Q, @\libconcept{quantity_of}@<dimensionless> Q2>
    requires @\seebelownc@
  friend constexpr decltype(auto) operator*=(Q&& lhs, const Q2& rhs);
  template<@\exposconceptnc{Mutable}@<quantity> Q, @\libconcept{quantity_of}@<dimensionless> Q2>
    requires @\seebelownc@
  friend constexpr decltype(auto) operator/=(Q&& lhs, const Q2& rhs);

  // \ref{qty.val.cmp}, value comparison
  friend constexpr bool is_eq_zero(const quantity& q) requires @\seebelownc@;
  friend constexpr bool is_neq_zero(const quantity& q) requires @\seebelownc@;
  friend constexpr bool is_lt_zero(const quantity& q) requires @\seebelownc@;
  friend constexpr bool is_gt_zero(const quantity& q) requires @\seebelownc@;
  friend constexpr bool is_lteq_zero(const quantity& q) requires @\seebelownc@;
  friend constexpr bool is_gteq_zero(const quantity& q) requires @\seebelownc@;
};

template<@\libconcept{Representation}@ Value, @\libconcept{quantity_reference}@ R>
quantity(Value, R) -> quantity<R{}, Value>;

template<@\libconcept{Representation}@ Value>
quantity(Value) -> quantity<one, Value>;

template<@\exposconceptnc{quantity-like}@ Q>
explicit(quantity_like_traits<Q>::explicit_import) quantity(Q)
  -> quantity<quantity_like_traits<Q>::reference, typename quantity_like_traits<Q>::rep>;

template<@\libconcept{delta_quantity}@ Q, @\libconcept{point_origin_for}@<Q::quantity_spec> PO>
quantity(Q, PO) -> quantity<@\exposidnc{make-point-reference}@(Q::quantity_spec, Q::unit, PO{}), typename Q::rep>;

}
\end{codeblock}

\pnum
\tcode{quantity<R, Rep>} is a structural type\irefcppx{temp.param}{term.structural.type}
if \tcode{Rep} is a structural type.

\pnum
If \tcode{decltype(R)} models \exposconcept{PointReference},
\exposid{numerical-value} is the numerical value\irefiev{112-01-29} in \tcode{unit}
of the delta of the point from \tcode{point_origin}.

\pnum
The member \tcode{absolute_point_origin} is equal to \tcode{point_origin} if
\begin{codeblock}
@\exposidnc{is-derived-from-specialization-of}@<decltype(auto(point_origin)), std::absolute_point_origin>()
\end{codeblock}
is \tcode{true}, and
to \tcode{point_origin.\exposidnc{quantity-point}.absolute_point_origin} otherwise.

\pnum
\begin{note}
\tcode{\libconcept{quantity_of}<Q, V>} is satisfied
by a delta whose quantity specification converts implicitly to \tcode{V},
by a point whose quantity specification converts implicitly to \tcode{V}
when \tcode{V} is the quantity specification of a point,
that is, by a point measured from an origin compatible with the origin of \tcode{V}\iref{qty.spec.conv},
by a point of the same absolute point origin as \tcode{V}
when \tcode{V} is a point origin, and
by a delta or a point of the character \tcode{V}
when \tcode{V} is a quantity character.
A delta never satisfies \tcode{\libconcept{quantity_of}<Q, point<QS>>}
and a point never satisfies \tcode{\libconcept{quantity_of}<Q, QS>}.
\tcode{\libconcept{quantity_of}<Q, V>} subsumes
\tcode{\libconcept{delta_quantity}<Q> || \libconcept{point_quantity}<Q>}.
\end{note}

\indexlibraryglobalexpos{to-common-quantity}
\begin{itemdecl}
template<typename Common, auto R, typename Rep>
  constexpr decltype(auto) @\exposidnc{to-common-quantity}@(const quantity<R, Rep>& q);         // \expos
\end{itemdecl}

\begin{itemdescr}
\pnum
\returns
\tcode{q} if
\tcode{std::is_same_v<Rep, typename Common::rep> \&\& \exposid{equivalent}(get_unit(R), Common::unit)}
is \tcode{true};
otherwise, \tcode{Common(q)}.
\begin{note}
An operand that already has the representation type and unit of \tcode{Common}
is not copied.
\end{note}
\end{itemdescr}

\rSec2[qty.static]{Static member functions}

\indexlibrarymember{zero}{quantity}
\indexlibrarymember{one}{quantity}
\indexlibrarymember{min}{quantity}
\indexlibrarymember{max}{quantity}
\begin{itemdecl}
static constexpr quantity zero() noexcept
  requires @\seebelownc@;
static constexpr quantity one() noexcept
  requires @\seebelownc@;
static constexpr quantity min() noexcept
  requires @\seebelownc@;
static constexpr quantity max() noexcept
  requires @\seebelownc@;
\end{itemdecl}

\begin{itemdescr}
\pnum
Let \placeholder{F} be one of \tcode{zero}, \tcode{one}, \tcode{min}, and \tcode{max}.

\pnum
\returns
\tcode{\{quantity_type::\placeholdernc{F}(), point_origin\}}
if \tcode{decltype(R)} models \exposconcept{PointReference}, and
\tcode{\{representation_values<rep>::\placeholdernc{F}(), R\}} otherwise.

\pnum
\remarks
Let \placeholder{C} be
\begin{itemize}
\item
\tcode{(!\exposconceptnc{PointReference}<decltype(R)> || \exposconceptnc{DefaultOriginPointReference}<decltype(R)>)}
if \placeholder{F} is \tcode{zero},
\item
\tcode{(!\exposconceptnc{PointReference}<decltype(R)>)}
if \placeholder{F} is \tcode{one}, and
\item
\tcode{true} otherwise.
\end{itemize}
The expression in the \fakegrammarterm{requires-clause} is equivalent to:
\begin{codeblock}
@\placeholder{C}@ && requires { representation_values<rep>::@\placeholdernc{F}@(); }
\end{codeblock}
\end{itemdescr}

\rSec2[qty.cons]{Constructors and assignment}

\indexlibraryctor{quantity}
\begin{itemdecl}
template<typename FwdValue, @\libconcept{quantity_reference}@ R2>
  requires @\exposconceptnc{SameValueAs}@<R2{}, R, std::remove_cvref_t<FwdValue>, Rep> &&
           @\exposconceptnc{SameKindReference}@<R2, R>
constexpr quantity(FwdValue&& v, R2);

template<@\exposconceptnc{ValuePreservingTo}@<Rep> FwdValue>
  requires(!@\exposconceptnc{PointReference}@<decltype(R)>) && (unit == std::one)
constexpr quantity(FwdValue&& v);
\end{itemdecl}

\begin{itemdescr}
\pnum
\effects
Initializes \exposidnc{numerical-value} with \tcode{std::forward<FwdValue>(v)}.

\pnum
\begin{note}
\exposconcept{SameKindReference} makes \tcode{quantity\{42, m\}} a delta
and \tcode{quantity\{42, point<m>\}} a point,
and keeps a number with the reference of a delta from initializing a point
and a number with the reference of a point from initializing a delta.
\end{note}
\end{itemdescr}

\indexlibraryctor{quantity}
\begin{itemdecl}
template<typename FwdValue, @\libconcept{quantity_reference}@ R2, typename Value = std::remove_cvref_t<FwdValue>>
  requires(!@\exposconceptnc{SameValueAs}@<R2{}, R, Value, Rep>) &&
          @\exposconceptnc{QuantityConvertibleTo}@<quantity<R2{}, Value>, quantity>
constexpr quantity(FwdValue&& v, R2);
\end{itemdecl}

\begin{itemdescr}
\pnum
\effects
Equivalent to
\tcode{quantity(quantity<R2\{\}, Value>\{std::forward<FwdValue>(v), R2\{\}\})}.
\end{itemdescr}

\indexlibraryctor{quantity}
\begin{itemdecl}
template<@\exposconceptnc{QuantityConvertibleTo}@<quantity> Q>
constexpr explicit(!std::@\stdconcept{convertible_to}@<typename Q::rep, Rep>) quantity(const Q& q);
\end{itemdecl}

\begin{itemdescr}
\pnum
\effects
Equivalent to \tcode{\exposidnc{sudo-cast}<quantity>(q)}\iref{qty.non.mem.conv}.
\end{itemdescr}

\indexlibraryctor{quantity}
\begin{itemdecl}
template<@\exposconceptnc{quantity-like}@ Q>
  requires @\seebelownc@
constexpr explicit(@\seebelownc@) quantity(const Q& q);
\end{itemdecl}

\begin{itemdescr}
\pnum
\effects
Equivalent to:
\begin{codeblock}
quantity(@\exposidnc{quantity-like-type}@<Q>{quantity_like_traits<Q>::to_numerical_value(q),
                                quantity_like_traits<Q>::reference})
\end{codeblock}

\pnum
\remarks
The expression in the \fakegrammarterm{requires-clause} is equivalent to:
\begin{codeblock}
(@\libconcept{point_quantity}@<@\exposidnc{quantity-like-type}@<Q>> == @\exposconceptnc{PointReference}@<decltype(R)>) &&
  std::@\stdconcept{constructible_from}@<quantity, @\exposidnc{quantity-like-type}@<Q>>
\end{codeblock}
The expression inside \tcode{explicit} is equivalent to:
\begin{codeblock}
quantity_like_traits<Q>::explicit_import ||
  !std::@\stdconcept{convertible_to}@<@\exposidnc{quantity-like-type}@<Q>, quantity>
\end{codeblock}
\end{itemdescr}

\indexlibraryctor{quantity}
\begin{itemdecl}
template<typename FwdQ, typename Q = std::remove_cvref_t<FwdQ>>
  requires @\seebelownc@
constexpr explicit quantity(FwdQ&& q);
\end{itemdecl}

\begin{itemdescr}
\pnum
\effects
Initializes \exposidnc{numerical-value} with
\tcode{quantity_type(std::forward<FwdQ>(q)).\exposidnc{numerical-value}}.

\pnum
\remarks
The expression in the \fakegrammarterm{requires-clause} is equivalent to:
\begin{codeblock}
@\exposconceptnc{DefaultOriginPointReference}@<decltype(R)> &&
  @\libconcept{quantity_of}@<Q, get_quantity_spec(R.@\exposidnc{delta-reference}@)> &&
  std::@\stdconcept{constructible_from}@<quantity_type, FwdQ>
\end{codeblock}
\begin{note}
This constructor creates a point at the delta \tcode{q} from the default point origin
of its reference.
\end{note}
\end{itemdescr}

\indexlibraryctor{quantity}
\begin{itemdecl}
template<typename FwdQ, @\exposconceptnc{PointOrigin}@ PO2, typename Q = std::remove_cvref_t<FwdQ>>
  requires @\seebelownc@
constexpr quantity(FwdQ&& q, PO2);
\end{itemdecl}

\begin{itemdescr}
\pnum
\effects
If \tcode{PO2\{\} == point_origin} is \tcode{true},
initializes \exposidnc{numerical-value} with
\tcode{quantity_type(std::forward<FwdQ>(q)).\exposidnc{numerical-value}}.
Otherwise, equivalent to:
\begin{codeblock}
quantity(quantity<@\exposidnc{make-point-reference}@(Q::quantity_spec, Q::unit, PO2{}), typename Q::rep>{
  std::forward<FwdQ>(q), PO2{}})
\end{codeblock}

\pnum
\remarks
The expression in the \fakegrammarterm{requires-clause} is equivalent to:
\begin{codeblock}
@\exposconceptnc{PointReference}@<decltype(R)> && @\libconcept{delta_quantity}@<Q> &&
  @\libconcept{point_origin_for}@<PO2, Q::quantity_spec> &&
  std::@\stdconcept{constructible_from}@<quantity_type, FwdQ> &&
  @\exposconceptnc{SameAbsolutePointOriginAs}@<PO2, point_origin>
\end{codeblock}
\end{itemdescr}

\indexlibraryctor{quantity}
\begin{itemdecl}
template<auto R2, typename Rep2>
  requires @\seebelownc@
constexpr explicit(@\seebelownc@) quantity(const quantity<R2, Rep2>& qp);
\end{itemdecl}

\begin{itemdescr}
\pnum
Let \tcode{QP} be \tcode{quantity<R2, Rep2>}.

\pnum
\effects
If \tcode{point_origin == QP::point_origin} is \tcode{true},
initializes \exposidnc{numerical-value} with
\tcode{quantity_type(qp.quantity_from(point_origin)).\exposidnc{numerical-value}}.
Otherwise, initializes \exposidnc{numerical-value} with
\tcode{quantity_type(qp - point_origin).\exposidnc{numerical-value}}.

\pnum
\remarks
The expression in the \fakegrammarterm{requires-clause} is equivalent to:
\begin{codeblock}
@\exposconceptnc{PointReference}@<decltype(R)> && @\libconcept{quantity_of}@<QP, absolute_point_origin> &&
  std::@\stdconcept{constructible_from}@<quantity_type, typename QP::quantity_type>
\end{codeblock}
The expression inside \tcode{explicit} is equivalent to:
\begin{codeblock}
!std::@\stdconcept{convertible_to}@<typename QP::quantity_type, quantity_type>
\end{codeblock}
\end{itemdescr}

\indexlibrarymember{operator=}{quantity}
\begin{itemdecl}
template<@\exposconceptnc{ValuePreservingTo}@<Rep> FwdValue>
  requires(!@\exposconceptnc{PointReference}@<decltype(R)>) && (unit == std::one)
constexpr quantity& operator=(FwdValue&& v);
\end{itemdecl}

\begin{itemdescr}
\pnum
\effects
Equivalent to \tcode{\exposidnc{numerical-value} = std::forward<FwdValue>(v)}.

\pnum
\returns
\tcode{*this}.
\end{itemdescr}

\rSec2[qty.conv]{Conversions}

\indexlibrarymember{in}{quantity}
\indexlibrarymember{force_in}{quantity}
\begin{itemdecl}
template<@\exposconceptnc{UnitCompatibleWith}@<unit, quantity_spec> ToU>
  requires @\seebelownc@
constexpr @\libconcept{quantity_of}@<quantity_spec> auto in(ToU) const;

template<@\libconcept{representation_of}@<quantity_spec> ToRep>
  requires @\seebelownc@
constexpr @\libconcept{quantity_of}@<quantity_spec> auto in() const;

template<@\libconcept{representation_of}@<quantity_spec> ToRep,
         @\exposconceptnc{UnitCompatibleWith}@<unit, quantity_spec> ToU>
  requires @\seebelownc@
constexpr @\libconcept{quantity_of}@<quantity_spec> auto in(ToU) const;

template<@\exposconceptnc{UnitCompatibleWith}@<unit, quantity_spec> ToU>
  requires @\seebelownc@
constexpr @\libconcept{quantity_of}@<quantity_spec> auto force_in(ToU) const;

template<@\libconcept{representation_of}@<quantity_spec> ToRep>
  requires @\seebelownc@
constexpr @\libconcept{quantity_of}@<quantity_spec> auto force_in() const;

template<@\libconcept{representation_of}@<quantity_spec> ToRep,
         @\exposconceptnc{UnitCompatibleWith}@<unit, quantity_spec> ToU>
  requires @\seebelownc@
constexpr @\libconcept{quantity_of}@<quantity_spec> auto force_in(ToU) const;
\end{itemdecl}

\begin{itemdescr}
\pnum
Let \tcode{\placeholder{result-expr}} be, for each signature,
\begin{itemize}
\item
\tcode{quantity<\exposidnc{make-reference}(quantity_spec, ToU\{\}), Rep>\{*this\}} for the first,
\item
\tcode{quantity<reference, ToRep>\{*this\}} for the second,
\item
\tcode{quantity<\exposidnc{make-reference}(quantity_spec, ToU\{\}), ToRep>\{*this\}} for the third,
\item
\tcode{value_cast<ToU\{\}>(*this)} for the fourth,
\item
\tcode{value_cast<ToRep>(*this)} for the fifth, and
\item
\tcode{value_cast<ToU\{\}, ToRep>(*this)} for the sixth.
\end{itemize}

\pnum
If \tcode{decltype(R)} models \exposconcept{PointReference}, let
\begin{itemize}
\item
\tcode{PO2} be \tcode{ToU::\exposidnc{point-origin}}
if the signature has a \tcode{ToU} parameter,
\tcode{\exposconceptnc{DefaultOriginPointReference}<decltype(R)>} is satisfied, and
\tcode{\exposconceptnc{OffsetUnit}<ToU>} is satisfied,
and \tcode{point_origin} otherwise, and
\item
\tcode{\placeholder{converted-quantity-expr}} be an expression denoting
the call to the member function of the same name
with the same template arguments and function arguments
on \tcode{quantity_from(PO2)}.
\end{itemize}

\pnum
\effects
If \tcode{decltype(R)} models \exposconcept{PointReference}, equivalent to:
\begin{codeblock}
return std::quantity{@\placeholdernc{converted-quantity-expr}@, PO2};
\end{codeblock}
Otherwise, equivalent to: \tcode{return \placeholdernc{result-expr};}

\pnum
\remarks
If \tcode{decltype(R)} models \exposconcept{PointReference},
the expression in the \fakegrammarterm{requires-clause} is equivalent to:
\begin{codeblock}
requires { @\placeholdernc{converted-quantity-expr}@; }
\end{codeblock}
Otherwise, it is equivalent to:
\begin{itemize}
\item
\tcode{\exposconceptnc{QuantityConvertibleTo}<quantity, decltype(\placeholdernc{result-expr})>}
for the first three signatures, and
\item
\tcode{requires \{ \placeholdernc{result-expr}; \}} for the other signatures.
\end{itemize}

\pnum
\begin{note}
A point measured from the default point origin of its reference
is converted to the default point origin of the new unit,
so \tcode{point<deg_C>(20.).in(K)} is measured from the origin of \tcode{K}.
Every other point keeps its point origin.
\end{note}
\end{itemdescr}

\indexlibrarymember{point_for}{quantity}
\begin{itemdecl}
template<@\exposconceptnc{PointOrigin}@ NewPO>
  requires @\exposconceptnc{PointReference}@<decltype(R)> &&
           @\exposconceptnc{SameAbsolutePointOriginAs}@<NewPO, absolute_point_origin>
constexpr @\libconcept{quantity_of}@<(NewPO{})> auto point_for(NewPO new_origin) const;
\end{itemdecl}

\begin{itemdescr}
\pnum
\effects
Equivalent to:
\begin{codeblock}
if constexpr (std::is_same_v<NewPO, decltype(auto(point_origin))>)
  return *this;
else
  return std::quantity{*this - new_origin, new_origin};
\end{codeblock}
\end{itemdescr}

\indexlibrarymember{magnitude}{quantity}
\begin{itemdecl}
constexpr @\libconcept{delta_quantity}@ auto magnitude() const
  requires @\seebelownc@;
\end{itemdecl}

\begin{itemdescr}
\pnum
\returns
\tcode{std::quantity\{math::magnitude(\exposid{numerical-value}), \exposid{make-reference}(magnitude(quantity_spec), unit)\}}\iref{qty.mag.cpo,qty.spec.ops}.

\pnum
\remarks
The expression in the \fakegrammarterm{requires-clause} is equivalent to:
\begin{codeblock}
(!@\exposconceptnc{PointReference}@<decltype(R)>) && @\exposconceptnc{NonScalarOperand}@<decltype(auto(quantity_spec))> &&
  requires(const Rep& v) { math::magnitude(v); }
\end{codeblock}

\pnum
\begin{note}
The result is a scalar quantity of \tcode{magnitude(quantity_spec)},
which converts implicitly to the nearest scalar quantity above \tcode{quantity_spec} in its hierarchy\iref{magnitude.qty}:
the magnitude of a \tcode{quantity<isq::velocity[m / s], V>} is a
\tcode{quantity<magnitude(isq::velocity)[m / s], double>}, for a vector type \tcode{V} whose magnitude is a \tcode{double},
and it initializes a \tcode{quantity<isq::speed[m / s]>}, while the vector itself does not.
For a quantity of a kind, \tcode{magnitude(quantity_spec)} is \tcode{quantity_spec}.
A quantity of a quantity specification that admits only a scalar character has no magnitude.
\end{note}
\end{itemdescr}

\rSec2[qty.obs]{Observers}

\indexlibrarymember{numerical_value_ref_in}{quantity}
\begin{itemdecl}
template<std::@\libconcept{unit}@ U>
  requires(!@\exposconceptnc{PointReference}@<decltype(R)>) && (@\exposid{equivalent}@(U{}, unit))
constexpr rep& numerical_value_ref_in(U) & noexcept;
template<std::@\libconcept{unit}@ U>
  requires(!@\exposconceptnc{PointReference}@<decltype(R)> ||
           @\exposconceptnc{DefaultOriginPointReference}@<decltype(R)>) && (@\exposid{equivalent}@(U{}, unit))
constexpr const rep& numerical_value_ref_in(U) const & noexcept;
\end{itemdecl}

\begin{itemdescr}
\pnum
\returns
\exposidnc{numerical-value}.

\pnum
\begin{note}
Only a point measured from the default point origin of its reference
has a unique numerical value, the one on the scale of its unit.
The numerical value of any other point is obtained from its delta from an explicit origin,
\tcode{quantity_from(origin).numerical_value_in(u)}.
\end{note}
\end{itemdescr}

\indexlibrarymember{numerical_value_in}{quantity}
\begin{itemdecl}
template<@\exposconceptnc{UnitCompatibleWith}@<unit, quantity_spec> U>
  requires @\seebelownc@
constexpr rep numerical_value_in(U) const noexcept;
\end{itemdecl}

\begin{itemdescr}
\pnum
\effects
Equivalent to:
\tcode{return (*this).in(U\{\}).\exposidnc{numerical-value};}

\pnum
\remarks
The expression in the \fakegrammarterm{requires-clause} is equivalent to:
\begin{codeblock}
(!@\exposconceptnc{PointReference}@<decltype(R)> || @\exposconceptnc{DefaultOriginPointReference}@<decltype(R)>) &&
  requires(const quantity q) { q.in(U{}); }
\end{codeblock}
\end{itemdescr}

\indexlibrarymember{force_numerical_value_in}{quantity}
\begin{itemdecl}
template<@\exposconceptnc{UnitCompatibleWith}@<unit, quantity_spec> U>
  requires @\seebelownc@
constexpr rep force_numerical_value_in(U) const noexcept;
\end{itemdecl}

\begin{itemdescr}
\pnum
\effects
Equivalent to:
\tcode{return (*this).force_in(U\{\}).\exposidnc{numerical-value};}

\pnum
\remarks
The expression in the \fakegrammarterm{requires-clause} is equivalent to:
\begin{codeblock}
(!@\exposconceptnc{PointReference}@<decltype(R)> || @\exposconceptnc{DefaultOriginPointReference}@<decltype(R)>) &&
  requires(const quantity q) { q.force_in(U{}); }
\end{codeblock}
\end{itemdescr}

\indexlibrarymember{numerical_value_in}{quantity}
\indexlibrarymember{force_numerical_value_in}{quantity}
\begin{itemdecl}
template<@\libconcept{representation_of}@<quantity_spec> ToRep, @\exposconceptnc{UnitCompatibleWith}@<unit, quantity_spec> U>
  requires @\seebelownc@
constexpr ToRep numerical_value_in(U) const;

template<@\libconcept{representation_of}@<quantity_spec> ToRep, @\exposconceptnc{UnitCompatibleWith}@<unit, quantity_spec> U>
  requires @\seebelownc@
constexpr ToRep force_numerical_value_in(U) const;
\end{itemdecl}

\begin{itemdescr}
\pnum
Let \tcode{\placeholder{conv}} be \tcode{in} for the first signature and \tcode{force_in} for the second.

\pnum
\effects
Equivalent to:
\tcode{return (*this).template \placeholder{conv}<ToRep>(U\{\}).\exposidnc{numerical-value};}

\pnum
\remarks
The expression in the \fakegrammarterm{requires-clause} is equivalent to:
\begin{codeblock}
(!@\exposconceptnc{PointReference}@<decltype(R)> || @\exposconceptnc{DefaultOriginPointReference}@<decltype(R)>) &&
  requires(const quantity q) { q.template @\placeholdernc{conv}@<ToRep>(U{}); }
\end{codeblock}

\pnum
\begin{note}
These return the numerical value in another representation type,
with the rules of \tcode{in<ToRep>(U\{\})} and \tcode{force_in<ToRep>(U\{\})}\iref{qty.conv}:
\tcode{(1567 * m).numerical_value_in<int>(km)} is ill-formed,
while \tcode{(1567 * m).force_numerical_value_in<int>(km)} is \tcode{1}.
\end{note}
\end{itemdescr}

\indexlibrarymember{quantity_from}{quantity}
\begin{itemdecl}
template<@\exposconceptnc{PointOrigin}@ PO2>
  requires @\exposconceptnc{PointReference}@<decltype(R)> && requires(const quantity qp) { qp - PO2{}; }
constexpr @\libconcept{delta_quantity}@ auto quantity_from(PO2 rhs) const &;

template<typename QP>
  requires @\exposconceptnc{PointReference}@<decltype(R)> && @\libconcept{quantity_of}@<QP, absolute_point_origin>
constexpr @\libconcept{delta_quantity}@ auto quantity_from(const QP& rhs) const;
\end{itemdecl}

\begin{itemdescr}
\pnum
\effects
Equivalent to:
\tcode{return *this - rhs;}
\end{itemdescr}

\indexlibrarymember{quantity_from}{quantity}
\begin{itemdecl}
template<@\exposconceptnc{PointOrigin}@ PO2>
  requires @\exposconceptnc{PointReference}@<decltype(R)> && requires(const quantity qp) { qp - PO2{}; }
constexpr @\libconcept{delta_quantity}@ auto quantity_from(PO2 rhs) &&;
\end{itemdecl}

\begin{itemdescr}
\pnum
\effects
Equivalent to:
\begin{codeblock}
if constexpr (PO2{} == point_origin)
  return quantity_type{std::move(@\exposidnc{numerical-value}@), R.@\exposidnc{delta-reference}@};
else
  return *this - rhs;
\end{codeblock}
\end{itemdescr}

\indexlibrarymember{quantity_from_zero}{quantity}
\begin{itemdecl}
constexpr @\libconcept{delta_quantity}@ auto quantity_from_zero() const
  requires @\exposconceptnc{DefaultOriginPointReference}@<decltype(R)>;
\end{itemdecl}

\begin{itemdescr}
\pnum
\effects
Equivalent to:
\tcode{return quantity_from(point_origin);}
\end{itemdescr}

\rSec2[qty.conv.ops]{Conversion operations}

\indexlibrarymisc{\idxcode{quantity}}{\idxcode{operator constructible_from<Rep>}}
\begin{itemdecl}
template<typename V_, std::@\stdconcept{constructible_from}@<Rep> Value = std::remove_cvref_t<V_>>
  requires(!@\exposconceptnc{PointReference}@<decltype(R)>) && (unit == std::one)
explicit operator V_() const & noexcept;
\end{itemdecl}

\begin{itemdescr}
\pnum
\returns
\exposidnc{numerical-value}.
\end{itemdescr}

\indexlibrarymisc{\idxcode{quantity}}{\idxcode{operator \exposconceptnc{quantity-like}}}
\begin{itemdecl}
template<typename Q_, @\exposconceptnc{quantity-like}@ Q = std::remove_cvref_t<Q_>>
  requires @\seebelownc@
constexpr explicit(@\seebelownc@) operator Q_() const noexcept(@\seebelownc@);
\end{itemdecl}

\begin{itemdescr}
\pnum
Let
\begin{itemize}
\item
\tcode{L} be \tcode{\exposidnc{quantity-like-type}<Q>},
\item
\tcode{u} be \tcode{get_unit(quantity_like_traits<Q>::reference)}, and
\item
\tcode{\placeholder{number-expr}} be
\tcode{L\{*this\}.quantity_from(L::point_origin).numerical_value_in(u)}
if \tcode{L} models \libconcept{point_quantity}, and
\tcode{numerical_value_in(u)} otherwise.
\end{itemize}

\pnum
\effects
Equivalent to:
\tcode{return quantity_like_traits<Q>::from_numerical_value(\placeholdernc{number-expr});}

\pnum
\remarks
The expression in the \fakegrammarterm{requires-clause} is equivalent to:
\begin{codeblock}
(@\libconcept{point_quantity}@<L> == @\exposconceptnc{PointReference}@<decltype(R)>) && std::@\stdconcept{constructible_from}@<L, quantity>
\end{codeblock}
The expression inside \tcode{explicit} is equivalent to:
\begin{codeblock}
quantity_like_traits<Q>::explicit_export || !std::@\stdconcept{convertible_to}@<quantity, L>
\end{codeblock}
The exception specification is equivalent to:
\begin{codeblock}
noexcept(quantity_like_traits<Q>::from_numerical_value(@\exposidnc{numerical-value}@)) &&
  std::is_nothrow_copy_constructible_v<rep>
\end{codeblock}
\end{itemdescr}

\rSec2[qty.unary.ops]{Unary operations}

\pnum
In the following descriptions,
let \tcode{@} be the \fakegrammarterm{operator}.

\indexlibrarymember{operator+}{quantity}
\indexlibrarymember{operator-}{quantity}
\begin{itemdecl}
constexpr @\libconcept{quantity_of}@<quantity_spec> auto operator+() const
  requires @\seebelownc@;
constexpr @\libconcept{quantity_of}@<quantity_spec> auto operator-() const
  requires @\seebelownc@;
\end{itemdecl}

\begin{itemdescr}
\pnum
\effects
Equivalent to:
\tcode{return std::quantity\{@\exposidnc{numerical-value}, reference\};}

\pnum
\remarks
The expression in the \fakegrammarterm{requires-clause} is equivalent to:
\begin{codeblock}
(!@\exposconceptnc{PointReference}@<decltype(R)>) && requires(const rep v) {
  { @\atsign@v } -> std::@\stdconcept{common_with}@<rep>;
}
\end{codeblock}
\end{itemdescr}

\indexlibrarymember{operator++}{quantity}
\indexlibrarymember{operator--}{quantity}
\begin{itemdecl}
template<@\exposconceptnc{Mutable}@<quantity> Q>
friend constexpr decltype(auto) operator++(Q&& q)
  requires @\seebelownc@;
template<@\exposconceptnc{Mutable}@<quantity> Q>
friend constexpr decltype(auto) operator--(Q&& q)
  requires @\seebelownc@;
\end{itemdecl}

\begin{itemdescr}
\pnum
\effects
Equivalent to
\tcode{@q.\exposidnc{numerical-value}}.

\pnum
\returns
\tcode{std::forward<Q>(q)}.

\pnum
\remarks
The expression in the \fakegrammarterm{requires-clause} is equivalent to:
\begin{codeblock}
requires(rep& v) {
  { @\atsign@v } -> std::@\stdconcept{same_as}@<rep&>;
}
\end{codeblock}
\begin{note}
A point is moved by a delta of one \tcode{unit}.
\end{note}
\end{itemdescr}

\indexlibrarymember{operator++}{quantity}
\indexlibrarymember{operator--}{quantity}
\begin{itemdecl}
constexpr @\libconcept{quantity_of}@<quantity_spec> auto operator++(int)
  requires @\seebelownc@;
constexpr @\libconcept{quantity_of}@<quantity_spec> auto operator--(int)
  requires @\seebelownc@;
\end{itemdecl}

\begin{itemdescr}
\pnum
\effects
Equivalent to:
\tcode{return std::quantity\{\exposid{numerical-value}@, reference\};}

\pnum
\remarks
The expression in the \fakegrammarterm{requires-clause} is equivalent to:
\begin{codeblock}
requires(rep& v) {
  { v@\atsign@ } -> std::@\stdconcept{common_with}@<rep>;
}
\end{codeblock}
\end{itemdescr}

\rSec2[qty.assign.ops]{Compound assignment operations}

\pnum
In the following descriptions,
let \tcode{@} be the \fakegrammarterm{operator}.

\indexlibrarymember{operator+=}{quantity}
\indexlibrarymember{operator-=}{quantity}
\indexlibrarymember{operator\%=}{quantity}
\begin{itemdecl}
template<@\exposconceptnc{Mutable}@<quantity> Q, auto R2, typename Rep2>
  requires @\seebelownc@
friend constexpr decltype(auto) operator+=(Q&& lhs, const quantity<R2, Rep2>& rhs);
template<@\exposconceptnc{Mutable}@<quantity> Q, auto R2, typename Rep2>
  requires @\seebelownc@
friend constexpr decltype(auto) operator-=(Q&& lhs, const quantity<R2, Rep2>& rhs);
template<@\exposconceptnc{Mutable}@<quantity> Q, auto R2, typename Rep2>
  requires @\seebelownc@
friend constexpr decltype(auto) operator%=(Q&& lhs, const quantity<R2, Rep2>& rhs);
\end{itemdecl}

\begin{itemdescr}
\pnum
\expects
If \tcode{@} is \tcode{\%=}, then \tcode{is_neq_zero(rhs)} is \tcode{true}.

\pnum
\effects
Equivalent to
\tcode{lhs.\exposid{numerical-value} @ rhs.in(lhs.unit).\exposid{numerical-value}}.

\pnum
\returns
\tcode{std::forward<Q>(lhs)}.

\pnum
\remarks
Let
\begin{itemize}
\item
\placeholder{D} be \tcode{quantity_type}
if \tcode{decltype(R)} models \exposconcept{PointReference},
and \tcode{quantity} otherwise, and
\item
\placeholder{C} be
\tcode{(!treat_as_floating_point<rep>) \&\& (!\exposconceptnc{PointReference}<decltype(R)>)}
if \tcode{@} is \tcode{\%=}, and
\tcode{true} otherwise.
\end{itemize}
The expression in the \fakegrammarterm{requires-clause} is equivalent to:
\begin{codeblock}
@\exposconceptnc{QuantityConvertibleTo}@<quantity<R2, Rep2>, @\placeholder{D}@> && @\placeholder{C}@ &&
requires(rep& a, const Rep2 b) {
  { a @\atsign@ b } -> std::@\stdconcept{same_as}@<rep&>;
}
\end{codeblock}
\begin{note}
A point is moved by a delta,
and the delta is in the quantity of the deltas measured from its point origin.
\end{note}

\pnum
\recommended
If \tcode{\exposid{equivalent}(unit, get_unit(rhs.reference))} is \tcode{true},
then the expression \tcode{rhs.in(lhs.unit)} is replaced with \tcode{rhs}.
\end{itemdescr}

\indexlibrarymember{operator*=}{quantity}
\indexlibrarymember{operator/=}{quantity}
\begin{itemdecl}
template<@\exposconceptnc{Mutable}@<quantity> Q, @\exposconceptnc{ValuePreservingTo}@<Rep> Value>
  requires @\seebelownc@
friend constexpr decltype(auto) operator*=(Q&& lhs, const Value& rhs);
template<@\exposconceptnc{Mutable}@<quantity> Q, @\exposconceptnc{ValuePreservingTo}@<Rep> Value>
  requires @\seebelownc@
friend constexpr decltype(auto) operator/=(Q&& lhs, const Value& rhs);
\end{itemdecl}

\begin{itemdescr}
\pnum
\expects
If \tcode{@} is \tcode{/=}, then \tcode{rhs != representation_values<Value>::zero()} is \tcode{true}.

\pnum
\effects
Equivalent to
\tcode{lhs.\exposid{numerical-value} @ rhs}.

\pnum
\returns
\tcode{std::forward<Q>(lhs)}.

\pnum
\remarks
The expression in the \fakegrammarterm{requires-clause} is equivalent to:
\begin{codeblock}
(!@\exposconceptnc{PointReference}@<decltype(R)>) && (!@\libconcept{delta_quantity}@<Value>) &&
requires(rep& a, const Value b) {
  { a @\atsign@ b } -> std::@\stdconcept{same_as}@<rep&>;
}
\end{codeblock}
\end{itemdescr}

\indexlibrarymember{operator*=}{quantity}
\indexlibrarymember{operator/=}{quantity}
\begin{itemdecl}
template<@\exposconceptnc{Mutable}@<quantity> Q, @\libconcept{quantity_of}@<dimensionless> Q2>
  requires @\seebelownc@
friend constexpr decltype(auto) operator*=(Q&& lhs, const Q2& rhs);
template<@\exposconceptnc{Mutable}@<quantity> Q, @\libconcept{quantity_of}@<dimensionless> Q2>
  requires @\seebelownc@
friend constexpr decltype(auto) operator/=(Q&& lhs, const Q2& rhs);
\end{itemdecl}

\begin{itemdescr}
\pnum
\effects
Equivalent to:
\tcode{return std::forward<Q>(lhs) @ rhs.\exposid{numerical-value};}

\pnum
\remarks
The expression in the \fakegrammarterm{requires-clause} is equivalent to:
\begin{codeblock}
(!@\exposconceptnc{PointReference}@<decltype(R)>) &&
(Q2::unit == std::one) && @\exposconceptnc{ValuePreservingTo}@<typename Q2::rep, Rep> &&
requires(rep& a, const Q2::rep b) {
  { a @\atsign@ b } -> std::@\stdconcept{same_as}@<rep&>;
}
\end{codeblock}
\end{itemdescr}

\rSec2[qty.arith.ops]{Arithmetic operations}

\pnum
In the following descriptions,
let \tcode{@} be the \fakegrammarterm{operator}.

\indexlibrarymember{operator+}{quantity}
\indexlibrarymember{operator-}{quantity}
\indexlibrarymember{operator\%}{quantity}
\begin{itemdecl}
template<auto R1, typename Rep1, auto R2, typename Rep2>
  requires @\exposconceptnc{CommonlyInvocableQuantities}@<std::plus<>, quantity<R1, Rep1>, quantity<R2, Rep2>>
friend constexpr @\libconcept{delta_quantity}@ auto operator+(const quantity<R1, Rep1>& lhs,
                                               const quantity<R2, Rep2>& rhs);
template<auto R1, typename Rep1, auto R2, typename Rep2>
  requires @\exposconceptnc{CommonlyInvocableQuantities}@<std::minus<>, quantity<R1, Rep1>, quantity<R2, Rep2>>
friend constexpr @\libconcept{delta_quantity}@ auto operator-(const quantity<R1, Rep1>& lhs,
                                               const quantity<R2, Rep2>& rhs);
template<auto R1, typename Rep1, auto R2, typename Rep2>
  requires(!treat_as_floating_point<Rep1>) && (!treat_as_floating_point<Rep2>) &&
          @\exposconceptnc{CommonlyInvocableQuantities}@<std::modulus<>, quantity<R1, Rep1>, quantity<R2, Rep2>>
friend constexpr @\libconcept{delta_quantity}@ auto operator%(const quantity<R1, Rep1>& lhs,
                                               const quantity<R2, Rep2>& rhs);
\end{itemdecl}

\begin{itemdescr}
\pnum
Let \placeholder{F} be the first argument to \exposconceptnc{CommonlyInvocableQuantities}.

\pnum
\expects
If \tcode{@} is \tcode{\%}, then \tcode{is_neq_zero(rhs)} is \tcode{true}.

\pnum
\effects
Equivalent to:
\begin{codeblock}
using ret = @\exposidnc{common-quantity-for}@<@\placeholdernc{F}@, quantity<R1, Rep1>, quantity<R2, Rep2>>;
const auto& ret_lhs = @\exposidnc{to-common-quantity}@<ret>(lhs);
const auto& ret_rhs = @\exposidnc{to-common-quantity}@<ret>(rhs);
return std::quantity{
  ret_lhs.numerical_value_ref_in(ret::unit) @\atsign@ ret_rhs.numerical_value_ref_in(ret::unit),
  ret::reference};
\end{codeblock}
\end{itemdescr}

\indexlibrarymember{operator+}{quantity}
\indexlibrarymember{operator-}{quantity}
\indexlibrarymember{operator\%}{quantity}
\begin{itemdecl}
template<auto R1, typename Rep1, @\libconcept{Representation}@ Value>
  requires(!@\exposconceptnc{PointReference}@<decltype(R1)>) && (quantity<R1, Rep1>::unit == std::one) &&
          @\exposconceptnc{InvokeResultOf}@<quantity<R1, Rep1>::quantity_spec, std::plus<>, Rep1, const Value&>
friend constexpr @\libconcept{delta_quantity}@ auto operator+(const quantity<R1, Rep1>& lhs, const Value& rhs);
template<auto R1, typename Rep1, @\libconcept{Representation}@ Value>
  requires(!@\exposconceptnc{PointReference}@<decltype(R1)>) && (quantity<R1, Rep1>::unit == std::one) &&
          @\exposconceptnc{InvokeResultOf}@<quantity<R1, Rep1>::quantity_spec, std::minus<>, Rep1, const Value&>
friend constexpr @\libconcept{delta_quantity}@ auto operator-(const quantity<R1, Rep1>& lhs, const Value& rhs);
template<auto R1, typename Rep1, @\libconcept{Representation}@ Value>
  requires(!@\exposconceptnc{PointReference}@<decltype(R1)>) && (quantity<R1, Rep1>::unit == std::one) &&
          @\exposconceptnc{InvokeResultOf}@<quantity<R1, Rep1>::quantity_spec, std::modulus<>, Rep1, const Value&>
friend constexpr @\libconcept{delta_quantity}@ auto operator%(const quantity<R1, Rep1>& lhs, const Value& rhs);
\end{itemdecl}

\begin{itemdescr}
\pnum
\effects
Equivalent to:
\tcode{return lhs @ std::quantity\{rhs\};}
\end{itemdescr}

\indexlibrarymember{operator+}{quantity}
\indexlibrarymember{operator-}{quantity}
\indexlibrarymember{operator\%}{quantity}
\begin{itemdecl}
template<auto R1, typename Rep1, @\libconcept{Representation}@ Value>
  requires(!@\exposconceptnc{PointReference}@<decltype(R1)>) && (quantity<R1, Rep1>::unit == std::one) &&
          @\exposconceptnc{InvokeResultOf}@<quantity<R1, Rep1>::quantity_spec, std::plus<>, Rep1, const Value&>
friend constexpr @\libconcept{delta_quantity}@ auto operator+(const Value& lhs, const quantity<R1, Rep1>& rhs);
template<auto R1, typename Rep1, @\libconcept{Representation}@ Value>
  requires(!@\exposconceptnc{PointReference}@<decltype(R1)>) && (quantity<R1, Rep1>::unit == std::one) &&
          @\exposconceptnc{InvokeResultOf}@<quantity<R1, Rep1>::quantity_spec, std::minus<>, Rep1, const Value&>
friend constexpr @\libconcept{delta_quantity}@ auto operator-(const Value& lhs, const quantity<R1, Rep1>& rhs);
template<auto R1, typename Rep1, @\libconcept{Representation}@ Value>
  requires(!@\exposconceptnc{PointReference}@<decltype(R1)>) && (quantity<R1, Rep1>::unit == std::one) &&
          @\exposconceptnc{InvokeResultOf}@<quantity<R1, Rep1>::quantity_spec, std::modulus<>, Rep1, const Value&>
friend constexpr @\libconcept{delta_quantity}@ auto operator%(const Value& lhs, const quantity<R1, Rep1>& rhs);
\end{itemdecl}

\begin{itemdescr}
\pnum
\effects
Equivalent to:
\tcode{return std::quantity\{lhs\} @ rhs;}
\end{itemdescr}

\indexlibrarymember{operator*}{quantity}
\indexlibrarymember{operator/}{quantity}
\begin{itemdecl}
template<auto R1, typename Rep1, auto R2, typename Rep2>
  requires @\exposconceptnc{InvocableQuantities}@<std::multiplies<>, quantity<R1, Rep1>, quantity<R2, Rep2>>
friend constexpr @\libconcept{delta_quantity}@ auto operator*(const quantity<R1, Rep1>& lhs,
                                               const quantity<R2, Rep2>& rhs);
template<auto R1, typename Rep1, auto R2, typename Rep2>
  requires @\exposconceptnc{InvocableQuantities}@<std::divides<>, quantity<R1, Rep1>, quantity<R2, Rep2>>
friend constexpr @\libconcept{delta_quantity}@ auto operator/(const quantity<R1, Rep1>& lhs,
                                               const quantity<R2, Rep2>& rhs);
\end{itemdecl}

\begin{itemdescr}
\pnum
\expects
If \tcode{@} is \tcode{/}, then \tcode{is_neq_zero(rhs)} is \tcode{true}.

\pnum
\effects
Equivalent to:
\begin{codeblock}
return std::quantity{
  lhs.numerical_value_ref_in(lhs.unit) @\atsign@ rhs.numerical_value_ref_in(rhs.unit), R1 @\atsign@ R2};
\end{codeblock}
\end{itemdescr}

\indexlibrarymember{operator*}{quantity}
\indexlibrarymember{operator/}{quantity}
\begin{itemdecl}
template<auto R1, typename Rep1, typename Value>
  requires(!@\exposconceptnc{PointReference}@<decltype(R1)>) &&
          (!@\libconcept{delta_quantity}@<Value>) && (!@\libconcept{quantity_reference}@<Value>) &&
          @\exposconceptnc{InvokeResultOf}@<quantity<R1, Rep1>::quantity_spec, std::multiplies<>, Rep1, const Value&>
friend constexpr @\libconcept{quantity_of}@<quantity<R1, Rep1>::quantity_spec> auto operator*(const quantity<R1, Rep1>& lhs,
                                                                               const Value& rhs);
template<auto R1, typename Rep1, typename Value>
  requires(!@\exposconceptnc{PointReference}@<decltype(R1)>) &&
          (!@\libconcept{delta_quantity}@<Value>) && (!@\libconcept{quantity_reference}@<Value>) &&
          @\exposconceptnc{InvokeResultOf}@<quantity<R1, Rep1>::quantity_spec, std::divides<>, Rep1, const Value&>
friend constexpr @\libconcept{quantity_of}@<quantity<R1, Rep1>::quantity_spec> auto operator/(const quantity<R1, Rep1>& lhs,
                                                                               const Value& rhs);
\end{itemdecl}

\begin{itemdescr}
\pnum
\expects
If \tcode{@} is \tcode{/}, then \tcode{rhs != representation_values<Value>::zero()} is \tcode{true}.

\pnum
\effects
Equivalent to:
\begin{codeblock}
return std::quantity{lhs.numerical_value_ref_in(lhs.unit) @\atsign@ rhs, R1};
\end{codeblock}
\end{itemdescr}

\indexlibrarymember{operator*}{quantity}
\indexlibrarymember{operator/}{quantity}
\begin{itemdecl}
template<typename Value, auto R1, typename Rep1>
  requires(!@\exposconceptnc{PointReference}@<decltype(R1)>) &&
          (!@\libconcept{delta_quantity}@<Value>) && (!@\libconcept{quantity_reference}@<Value>) &&
          @\exposconceptnc{InvokeResultOf}@<quantity<R1, Rep1>::quantity_spec, std::multiplies<>, const Value&, Rep1>
friend constexpr @\libconcept{quantity_of}@<quantity<R1, Rep1>::quantity_spec> auto operator*(const Value& lhs,
                                                                               const quantity<R1, Rep1>& rhs);
template<typename Value, auto R1, typename Rep1>
  requires(!@\exposconceptnc{PointReference}@<decltype(R1)>) &&
          (!@\libconcept{delta_quantity}@<Value>) && (!@\libconcept{quantity_reference}@<Value>) &&
          @\exposconceptnc{ScalarOperand}@<decltype(auto(quantity<R1, Rep1>::quantity_spec))> &&
          @\exposconceptnc{InvokeResultOf}@<quantity<R1, Rep1>::quantity_spec, std::divides<>, const Value&, Rep1>
friend constexpr @\libconcept{delta_quantity}@ auto operator/(const Value& lhs, const quantity<R1, Rep1>& rhs);
\end{itemdecl}

\begin{itemdescr}
\pnum
\expects
If \tcode{@} is \tcode{/}, then \tcode{is_neq_zero(rhs)} is \tcode{true}.

\pnum
\effects
Equivalent to:
\begin{codeblock}
return std::quantity{lhs @\atsign@ rhs.numerical_value_ref_in(rhs.unit), std::one @\atsign@ R1};
\end{codeblock}
\end{itemdescr}

\indexlibrarymember{operator+}{quantity}
\indexlibrarymember{operator-}{quantity}
\begin{itemdecl}
template<auto R1, typename Rep1, @\libconcept{delta_quantity}@ Q2>
  requires @\seebelownc@
friend constexpr @\libconcept{point_quantity}@ auto operator+(const quantity<R1, Rep1>& qp, const Q2& q);
template<auto R1, typename Rep1, @\libconcept{delta_quantity}@ Q2>
  requires @\seebelownc@
friend constexpr @\libconcept{point_quantity}@ auto operator+(const Q2& q, const quantity<R1, Rep1>& qp);
template<auto R1, typename Rep1, @\libconcept{delta_quantity}@ Q2>
  requires @\seebelownc@
friend constexpr @\libconcept{point_quantity}@ auto operator-(const quantity<R1, Rep1>& qp, const Q2& q);
\end{itemdecl}

\begin{itemdescr}
\pnum
Let \tcode{\placeholder{delta-expr}} be
\tcode{qp.quantity_from(qp.point_origin) @ q} and \tcode{D} be \tcode{decltype(\placeholder{delta-expr})}.

\pnum
\effects
Equivalent to:
\begin{codeblock}
if constexpr (@\exposidnc{is-derived-from-specialization-of}@<decltype(auto(quantity<R1, Rep1>::point_origin)),
                                                     natural_point_origin_>() &&
              D::quantity_spec != quantity<R1, Rep1>::quantity_type::quantity_spec)
  return std::quantity{@\placeholdernc{delta-expr}@, natural_point_origin<D::quantity_spec>};
else
  return std::quantity{@\placeholdernc{delta-expr}@, qp.point_origin};
\end{codeblock}

\pnum
\remarks
The expression in the \fakegrammarterm{requires-clause} is equivalent to:
\begin{codeblock}
@\exposconceptnc{PointReference}@<decltype(R1)> &&
  @\libconcept{point_origin_for}@<decltype(auto(quantity<R1, Rep1>::point_origin)), Q2::quantity_spec> &&
  requires(const quantity<R1, Rep1>& qp, const Q2& q) { @\placeholdernc{delta-expr}@; }
\end{codeblock}

\pnum
\begin{note}
A natural point origin is re-expressed for the quantity specification of the result
when that differs from the one of the delta of \tcode{qp},
as it does when a point of a kind meets a committed delta:
\tcode{point<m>(1) + isq::height(2 * m)} is measured from \tcode{natural_point_origin<isq::height>}.
\end{note}
\end{itemdescr}

\indexlibrarymember{operator-}{quantity}
\begin{itemdecl}
template<auto R1, typename Rep1, @\libconcept{point_quantity}@ QP2>
  requires @\seebelownc@
friend constexpr @\libconcept{delta_quantity}@ auto operator-(const quantity<R1, Rep1>& lhs, const QP2& rhs);
\end{itemdecl}

\begin{itemdescr}
\pnum
\effects
Equivalent to:
\begin{codeblock}
if constexpr (quantity<R1, Rep1>::point_origin == QP2::point_origin)
  return lhs.quantity_from(lhs.point_origin) - rhs.quantity_from(QP2::point_origin);
else
  return lhs.quantity_from(lhs.point_origin) - rhs.quantity_from(QP2::point_origin) +
         (lhs.point_origin - QP2::point_origin);
\end{codeblock}

\pnum
\remarks
The expression in the \fakegrammarterm{requires-clause} is equivalent to:
\begin{codeblock}
@\exposconceptnc{PointReference}@<decltype(R1)> &&
  @\libconcept{quantity_of}@<QP2, quantity<R1, Rep1>::absolute_point_origin> &&
  requires(const quantity<R1, Rep1>& lhs, const QP2& rhs) {
    lhs.quantity_from(lhs.point_origin) - rhs.quantity_from(QP2::point_origin);
  }
\end{codeblock}
\end{itemdescr}

\rSec2[qty.cmp]{Comparison}

\pnum
In the following descriptions,
let \tcode{@} be the \fakegrammarterm{operator}.

\indexlibrarymember{operator==}{quantity}
\indexlibrarymember{operator<=>}{quantity}
\begin{itemdecl}
template<auto R1, typename Rep1, auto R2, typename Rep2>
  requires @\seebelownc@
friend constexpr bool operator==(const quantity<R1, Rep1>& lhs,
                                 const quantity<R2, Rep2>& rhs);
template<auto R1, typename Rep1, auto R2, typename Rep2>
  requires @\seebelownc@
friend constexpr auto operator<=>(const quantity<R1, Rep1>& lhs,
                                  const quantity<R2, Rep2>& rhs);
\end{itemdecl}

\begin{itemdescr}
\pnum
Let \placeholder{C} be
\tcode{std::\stdconcept{equality_comparable}} if \tcode{@} is \tcode{==}, and
\tcode{std::\stdconcept{three_way_comparable}} if \tcode{@} is \tcode{<=>}.

\pnum
\effects
Equivalent to:
\begin{codeblock}
using ct = std::common_type_t<quantity<R1, Rep1>, quantity<R2, Rep2>>;
const auto& ct_lhs = @\exposidnc{to-common-quantity}@<ct>(lhs);
const auto& ct_rhs = @\exposidnc{to-common-quantity}@<ct>(rhs);
return ct_lhs.numerical_value_ref_in(ct::unit) @\atsign@ ct_rhs.numerical_value_ref_in(ct::unit);
\end{codeblock}

\pnum
\remarks
The expression in the \fakegrammarterm{requires-clause} is equivalent to:
\begin{codeblock}
@\libconcept{delta_quantity}@<quantity<R1, Rep1>> && @\libconcept{delta_quantity}@<quantity<R2, Rep2>> && requires {
  typename std::common_type_t<quantity<R1, Rep1>, quantity<R2, Rep2>>;
} && @\placeholder{C}@<typename std::common_type_t<quantity<R1, Rep1>, quantity<R2, Rep2>>::rep>
\end{codeblock}
\end{itemdescr}

\indexlibrarymember{operator==}{quantity}
\indexlibrarymember{operator<=>}{quantity}
\begin{itemdecl}
template<auto R1, typename Rep1, @\libconcept{Representation}@ Value>
  requires @\seebelownc@
friend constexpr bool operator==(const quantity<R1, Rep1>& lhs, const Value& rhs);
template<auto R1, typename Rep1, @\libconcept{Representation}@ Value>
  requires @\seebelownc@
friend constexpr auto operator<=>(const quantity<R1, Rep1>& lhs, const Value& rhs);
\end{itemdecl}

\begin{itemdescr}
\pnum
Let \placeholder{C} be
\tcode{std::\stdconcept{equality_comparable_with}} if \tcode{@} is \tcode{==}, and
\tcode{std::\stdconcept{three_way_comparable_with}} if \tcode{@} is \tcode{<=>}.

\pnum
\returns
\tcode{lhs.numerical_value_ref_in(lhs.unit) @ rhs}.

\pnum
\remarks
The expression in the \fakegrammarterm{requires-clause} is equivalent to:
\begin{codeblock}
(!@\exposconceptnc{PointReference}@<decltype(R1)>) && (quantity<R1, Rep1>::unit == std::one) &&
  @\placeholder{C}@<Rep1, Value>
\end{codeblock}
\end{itemdescr}

\indexlibrarymember{operator==}{quantity}
\indexlibrarymember{operator<=>}{quantity}
\begin{itemdecl}
template<auto R1, typename Rep1, @\libconcept{point_quantity}@ QP2>
  requires @\seebelownc@
friend constexpr bool operator==(const quantity<R1, Rep1>& lhs, const QP2& rhs);
template<auto R1, typename Rep1, @\libconcept{point_quantity}@ QP2>
  requires @\seebelownc@
friend constexpr auto operator<=>(const quantity<R1, Rep1>& lhs, const QP2& rhs);
\end{itemdecl}

\begin{itemdescr}
\pnum
Let \placeholder{C} be
\tcode{std::\stdconcept{equality_comparable_with}} if \tcode{@} is \tcode{==}, and
\tcode{std::\stdconcept{three_way_comparable_with}} if \tcode{@} is \tcode{<=>}.

\pnum
\effects
Equivalent to:
\begin{codeblock}
if constexpr (quantity<R1, Rep1>::point_origin == QP2::point_origin)
  return lhs.quantity_from(lhs.point_origin) @\atsign@ rhs.quantity_from(QP2::point_origin);
else
  return lhs - lhs.absolute_point_origin @\atsign@ rhs - rhs.absolute_point_origin;
\end{codeblock}

\pnum
\remarks
The expression in the \fakegrammarterm{requires-clause} is equivalent to:
\begin{codeblock}
@\exposconceptnc{PointReference}@<decltype(R1)> &&
  @\libconcept{quantity_of}@<QP2, quantity<R1, Rep1>::absolute_point_origin> &&
  @\placeholder{C}@<typename quantity<R1, Rep1>::quantity_type, typename QP2::quantity_type> &&
  ((quantity<R1, Rep1>::point_origin == QP2::point_origin) ||
   requires(const quantity<R1, Rep1>& lhs, const QP2& rhs) {
     lhs - lhs.absolute_point_origin @\atsign@ rhs - rhs.absolute_point_origin;
   })
\end{codeblock}
\end{itemdescr}

\rSec2[qty.val.cmp]{Value comparison}

\indexlibrarymember{is_eq_zero}{quantity}
\indexlibrarymember{is_neq_zero}{quantity}
\indexlibrarymember{is_lt_zero}{quantity}
\indexlibrarymember{is_gt_zero}{quantity}
\indexlibrarymember{is_lteq_zero}{quantity}
\indexlibrarymember{is_gteq_zero}{quantity}
\begin{itemdecl}
friend constexpr bool is_eq_zero(const quantity& q) requires @\seebelownc@;
friend constexpr bool is_neq_zero(const quantity& q) requires @\seebelownc@;
friend constexpr bool is_lt_zero(const quantity& q) requires @\seebelownc@;
friend constexpr bool is_gt_zero(const quantity& q) requires @\seebelownc@;
friend constexpr bool is_lteq_zero(const quantity& q) requires @\seebelownc@;
friend constexpr bool is_gteq_zero(const quantity& q) requires @\seebelownc@;
\end{itemdecl}

\begin{itemdescr}
\pnum
Let \tcode{is_\placeholdernc{F}_zero} be the function name.

\pnum
\returns
\begin{itemize}
\item
If \tcode{\placeholdernc{F}} is \tcode{eq}, returns \tcode{q == zero()}.
\item
Otherwise, if \tcode{\placeholdernc{F}} is \tcode{neq}, returns \tcode{q != zero()}.
\item
Otherwise, returns \tcode{is_\placeholdernc{F}(q <=> zero())}\irefcpp{compare.syn}.
\end{itemize}

\pnum
\remarks
Let \placeholder{C} be
\tcode{std::\stdconcept{equality_comparable_with}} if \placeholder{F} is \tcode{eq} or \tcode{neq}, and
\tcode{std::\stdconcept{three_way_comparable_with}} otherwise.
The expression in the \fakegrammarterm{requires-clause} is equivalent to:
\begin{codeblock}
(!@\exposconceptnc{PointReference}@<decltype(R)>) && requires {
  { T::zero() } -> @\placeholder{C}@<quantity>;
}
\end{codeblock}
\end{itemdescr}

\ednote{
P3045R10 no longer proposes \tcode{is_eq_zero}, \tcode{is_neq_zero}, \tcode{is_lt_zero},
\tcode{is_gt_zero}, \tcode{is_lteq_zero}, and \tcode{is_gteq_zero}.
They are to be removed, together with their uses elsewhere in this wording,
in favor of comparing a \tcode{quantity} with \tcode{zero},
a named constant of a tag type that doesn't convert to arithmetic types,
as P4406R0 (Comparing Quantities Against Zero) recommends.
The wording of the tag type follows the outcome of the polls in P4406R0.
}

\rSec2[qty.delta]{Variable template \tcode{delta}}

\indexlibraryglobal{delta}
\begin{itemdecl}
template<auto Arg>
  requires @\seebelownc@
constexpr auto delta = @\seebelownc@;
\end{itemdecl}

\begin{itemdescr}
\pnum
The value is
\tcode{\exposid{delta-spec-t}<decltype(Arg)>\{\}}, the quantity specification of a delta\iref{delta.qty},
if \tcode{decltype(Arg)} models \libconcept{quantity_spec},
and \tcode{delta_<decltype(Arg)>\{\}}, the construction helper of a delta in the unit \tcode{Arg}, otherwise.

\pnum
\remarks
The expression in the \fakegrammarterm{requires-clause} is equivalent to:
\begin{codeblock}
(@\libconcept{quantity_spec}@<decltype(Arg)> && (!@\exposconceptnc{PointQuantitySpec}@<decltype(Arg)>) &&
 (!@\exposconceptnc{DeltaQuantitySpec}@<decltype(Arg)>)) ||
@\libconcept{unit}@<decltype(Arg)>
\end{codeblock}

\pnum
\begin{note}
\tcode{delta<>} takes a unit or a quantity specification, and nothing else:
\tcode{delta<isq::height[m]>}, \tcode{delta<delta<isq::height>>}, and \tcode{delta<point<isq::height>>}
are ill-formed.
A delta in a unit that has a point origin is created by \tcode{delta<U>(v)},
as \tcode{delta<deg_C>(5)}, and given a quantity specification by its call operator\iref{qty.spec.ops},
as \tcode{isq::thermodynamic_temperature(delta<deg_C>(5))}.
\end{note}
\end{itemdescr}

\indexlibraryglobal{delta_}
\begin{codeblock}
namespace std {

template<@\exposconceptnc{BareReference}@ R>
struct delta_ {
  template<typename FwdRep,
           @\libconcept{representation_of}@<get_quantity_spec(R{})> Rep = std::remove_cvref_t<FwdRep>>
  constexpr quantity<R{}, Rep> operator()(FwdRep&& lhs) const;
};

}
\end{codeblock}

\indexlibrarymember{operator()}{delta_}
\begin{itemdecl}
template<typename FwdRep,
         @\libconcept{representation_of}@<get_quantity_spec(R{})> Rep = std::remove_cvref_t<FwdRep>>
constexpr quantity<R{}, Rep> operator()(FwdRep&& lhs) const;
\end{itemdecl}

\begin{itemdescr}
\pnum
\effects
Equivalent to:
\tcode{return quantity\{std::forward<FwdRep>(lhs), R\{\}\};}
\end{itemdescr}

\rSec2[qty.non.mem.conv]{Non-member conversions}

\indexlibrarymemberexpos{quantity}{sudo-cast}
\begin{itemdecl}
template<@\libconcept{delta_quantity}@ To, typename FwdFrom, @\libconcept{delta_quantity}@ From = std::remove_cvref_t<FwdFrom>>
  requires @\seebelownc@
constexpr To @\exposidnc{sudo-cast}@(FwdFrom&& q);  // \expos
template<@\libconcept{point_quantity}@ To, typename FwdFrom, @\libconcept{point_quantity}@ From = std::remove_cvref_t<FwdFrom>>
  requires @\seebelownc@
constexpr To @\exposidnc{sudo-cast}@(FwdFrom&& qp);  // \expos
\end{itemdecl}

\begin{itemdescr}
\pnum
\returns
TBD.
\end{itemdescr}

\pnum
\tcode{value_cast} is an explicit cast that allows truncation.

\indexlibrarymember{value_cast}{quantity}
\begin{itemdecl}
template<@\libconcept{unit}@ auto ToU, typename FwdQ, @\libconcept{delta_quantity}@ Q = std::remove_cvref_t<FwdQ>>
  requires(convertible(Q::reference, ToU))
constexpr @\libconcept{delta_quantity}@ auto value_cast(FwdQ&& q);
\end{itemdecl}

\begin{itemdescr}
\pnum
\effects
Equivalent to:
\begin{codeblock}
return @\exposidnc{sudo-cast}@<quantity<@\exposidnc{make-reference}@(Q::quantity_spec, ToU), typename Q::rep>>(
  std::forward<FwdQ>(q));
\end{codeblock}
\end{itemdescr}

\indexlibrarymember{value_cast}{quantity}
\begin{itemdecl}
template<@\libconcept{Representation}@ ToRep, typename FwdQ, @\libconcept{delta_quantity}@ Q = std::remove_cvref_t<FwdQ>>
  requires @\libconcept{representation_of}@<ToRep, Q::quantity_spec> &&
           std::@\stdconcept{constructible_from}@<ToRep, typename Q::rep>
constexpr quantity<Q::reference, ToRep> value_cast(FwdQ&& q);
\end{itemdecl}

\begin{itemdescr}
\pnum
\effects
Equivalent to:
\begin{codeblock}
return @\exposidnc{sudo-cast}@<quantity<Q::reference, ToRep>>(std::forward<FwdQ>(q));
\end{codeblock}
\end{itemdescr}

\indexlibrarymember{value_cast}{quantity}
\begin{itemdecl}
template<@\libconcept{unit}@ auto ToU, @\libconcept{Representation}@ ToRep, typename FwdQ,
         @\libconcept{delta_quantity}@ Q = std::remove_cvref_t<FwdQ>>
  requires @\seebelownc@
constexpr @\libconcept{delta_quantity}@ auto value_cast(FwdQ&& q);
template<@\libconcept{Representation}@ ToRep, @\libconcept{unit}@ auto ToU, typename FwdQ,
         @\libconcept{delta_quantity}@ Q = std::remove_cvref_t<FwdQ>>
  requires @\seebelownc@
constexpr @\libconcept{delta_quantity}@ auto value_cast(FwdQ&& q);
\end{itemdecl}

\begin{itemdescr}
\pnum
\effects
Equivalent to:
\begin{codeblock}
return @\exposidnc{sudo-cast}@<quantity<@\exposidnc{make-reference}@(Q::quantity_spec, ToU), ToRep>>(
  std::forward<FwdQ>(q));
\end{codeblock}

\pnum
\remarks
The expression in the \fakegrammarterm{requires-clause} is equivalent to:
\begin{codeblock}
(convertible(Q::reference, ToU)) && @\libconcept{representation_of}@<ToRep, Q::quantity_spec> &&
std::@\stdconcept{constructible_from}@<ToRep, typename Q::rep>
\end{codeblock}
\end{itemdescr}

\indexlibrarymember{value_cast}{quantity}
\begin{itemdecl}
template<@\libconcept{delta_quantity}@ ToQ, typename FwdQ, @\libconcept{delta_quantity}@ Q = std::remove_cvref_t<FwdQ>>
  requires(convertible(Q::reference, ToQ::unit)) && (ToQ::quantity_spec == Q::quantity_spec) &&
          std::@\stdconcept{constructible_from}@<typename ToQ::rep, typename Q::rep>
constexpr @\libconcept{delta_quantity}@ auto value_cast(FwdQ&& q);
\end{itemdecl}

\begin{itemdescr}
\pnum
\effects
Equivalent to: \tcode{return \exposidnc{sudo-cast}<ToQ>(std::forward<FwdQ>(q));}
\end{itemdescr}

\indexlibrarymember{value_cast}{quantity}
\begin{itemdecl}
template<@\libconcept{unit}@ auto ToU, typename FwdQP, @\libconcept{point_quantity}@ QP = std::remove_cvref_t<FwdQP>>
  requires(convertible(QP::quantity_type::reference, ToU))
constexpr @\libconcept{point_quantity}@ auto value_cast(FwdQP&& qp);
\end{itemdecl}

\begin{itemdescr}
\pnum
\effects
Equivalent to:
\begin{codeblock}
return quantity{value_cast<ToU>(std::forward<FwdQP>(qp).quantity_from(QP::point_origin)),
                QP::point_origin};
\end{codeblock}
\end{itemdescr}

\indexlibrarymember{value_cast}{quantity}
\begin{itemdecl}
template<@\libconcept{Representation}@ ToRep, typename FwdQP, @\libconcept{point_quantity}@ QP = std::remove_cvref_t<FwdQP>>
  requires @\libconcept{representation_of}@<ToRep, QP::quantity_spec> &&
           std::@\stdconcept{constructible_from}@<ToRep, typename QP::rep>
constexpr quantity<QP::reference, ToRep> value_cast(FwdQP&& qp);
\end{itemdecl}

\begin{itemdescr}
\pnum
\effects
Equivalent to:
\begin{codeblock}
return {value_cast<ToRep>(std::forward<FwdQP>(qp).quantity_from(QP::point_origin)), QP::point_origin};
\end{codeblock}
\end{itemdescr}

\indexlibrarymember{value_cast}{quantity}
\begin{itemdecl}
template<@\libconcept{unit}@ auto ToU, @\libconcept{Representation}@ ToRep, typename FwdQP,
         @\libconcept{point_quantity}@ QP = std::remove_cvref_t<FwdQP>>
  requires @\seebelownc@
constexpr @\libconcept{point_quantity}@ auto value_cast(FwdQP&& qp);
template<@\libconcept{Representation}@ ToRep, @\libconcept{unit}@ auto ToU, typename FwdQP,
         @\libconcept{point_quantity}@ QP = std::remove_cvref_t<FwdQP>>
  requires @\seebelownc@
constexpr @\libconcept{point_quantity}@ auto value_cast(FwdQP&& qp);
\end{itemdecl}

\begin{itemdescr}
\pnum
\effects
Equivalent to:
\begin{codeblock}
return quantity{value_cast<ToU, ToRep>(std::forward<FwdQP>(qp).quantity_from(QP::point_origin)),
                QP::point_origin};
\end{codeblock}

\pnum
\remarks
The expression in the \fakegrammarterm{requires-clause} is equivalent to:
\begin{codeblock}
(convertible(QP::quantity_type::reference, ToU)) && @\libconcept{representation_of}@<ToRep, QP::quantity_spec> &&
std::@\stdconcept{constructible_from}@<ToRep, typename QP::rep>
\end{codeblock}
\end{itemdescr}

\indexlibrarymember{value_cast}{quantity}
\begin{itemdecl}
template<@\libconcept{delta_quantity}@ ToQ, typename FwdQP, @\libconcept{point_quantity}@ QP = std::remove_cvref_t<FwdQP>>
  requires(convertible(QP::quantity_type::reference, ToQ::unit)) &&
          (ToQ::quantity_spec == QP::quantity_type::quantity_spec) &&
          std::@\stdconcept{constructible_from}@<typename ToQ::rep, typename QP::rep>
constexpr @\libconcept{point_quantity}@ auto value_cast(FwdQP&& qp);
\end{itemdecl}

\begin{itemdescr}
\pnum
\effects
Equivalent to:
\begin{codeblock}
return quantity{value_cast<ToQ>(std::forward<FwdQP>(qp).quantity_from(QP::point_origin)),
                QP::point_origin};
\end{codeblock}
\end{itemdescr}

\indexlibrarymember{value_cast}{quantity}
\begin{itemdecl}
template<@\libconcept{point_quantity}@ ToQP, typename FwdQP, @\libconcept{point_quantity}@ QP = std::remove_cvref_t<FwdQP>>
  requires(convertible(QP::quantity_type::reference, ToQP::unit)) &&
          (@\exposidnc{remove-point}@(ToQP::quantity_spec) == @\exposidnc{remove-point}@(QP::quantity_spec)) &&
          (@\exposidnc{same-absolute-point-origins}@(ToQP::point_origin, QP::point_origin)) &&
          std::@\stdconcept{constructible_from}@<typename ToQP::rep, typename QP::rep>
constexpr @\libconcept{point_quantity}@ auto value_cast(FwdQP&& qp);
\end{itemdecl}

\begin{itemdescr}
\pnum
\effects
Equivalent to:
\tcode{return \exposidnc{sudo-cast}<ToQP>(std::forward<FwdQP>(qp));}
\end{itemdescr}

\pnum
\tcode{quantity_cast} is an explicit cast that allows converting to more specific quantities.
\begin{example}
\begin{codeblock}
auto length = isq::length(42 * m);
auto distance = quantity_cast<isq::distance>(length);
auto altitude = quantity_cast<point<isq::altitude>>(point<isq::height>[m](42));
\end{codeblock}
\end{example}

\indexlibrarymember{quantity_cast}{quantity}
\begin{itemdecl}
template<@\libconcept{quantity_spec}@ auto ToQS, typename FwdQ, @\libconcept{delta_quantity}@ Q = std::remove_cvref_t<FwdQ>>
  requires @\exposconceptnc{QuantitySpecCastableTo}@<Q::quantity_spec, ToQS>
constexpr @\libconcept{delta_quantity}@ auto quantity_cast(FwdQ&& q);
\end{itemdecl}

\begin{itemdescr}
\pnum
\effects
Equivalent to:
\begin{codeblock}
return quantity{std::forward<FwdQ>(q).@\exposidnc{numerical-value}@, @\exposidnc{make-reference}@(ToQS, Q::unit)};
\end{codeblock}
\end{itemdescr}

\indexlibrarymember{quantity_cast}{quantity}
\begin{itemdecl}
template<@\libconcept{quantity_spec}@ auto ToQS, typename FwdQP, @\libconcept{point_quantity}@ QP = std::remove_cvref_t<FwdQP>>
  requires @\seebelownc@
constexpr @\libconcept{point_quantity}@ auto quantity_cast(FwdQP&& qp);
\end{itemdecl}

\begin{itemdescr}
\pnum
Let \tcode{po} be \tcode{\exposid{conversion-origin}(ToQS, QP::point_origin)}\iref{qty.spec.ops}.

\pnum
\effects
Equivalent to:
\begin{codeblock}
return quantity{quantity_cast<@\exposidnc{remove-point}@(ToQS)>(std::forward<FwdQP>(qp).quantity_from(po)), po};
\end{codeblock}

\pnum
\remarks
The expression in the \fakegrammarterm{requires-clause} is equivalent to:
\begin{codeblock}
@\exposconceptnc{PointQuantitySpec}@<decltype(ToQS)> &&
  @\exposconceptnc{QuantitySpecCastableTo}@<QP::quantity_spec, @\exposidnc{point-target}@(ToQS, QP::point_origin)> &&
  requires(const QP& qp) { qp.quantity_from(@\exposidnc{conversion-origin}@(ToQS, QP::point_origin)); }
\end{codeblock}

\pnum
\begin{note}
A point is cast to a point quantity specification only, and a delta never to one.
A point is cast along one branch of the hierarchy only and only from a compatible origin\iref{qty.spec.conv},
so a point of \tcode{isq::height} becomes a point of \tcode{isq::width} only through a delta from its origin.
The result keeps the origin of \tcode{qp} when that is a point origin for the scale of \tcode{ToQS}
and \tcode{ToQS} names no other origin, and is measured from the origin of \tcode{ToQS} otherwise.
\end{note}
\end{itemdescr}

\rSec2[qty.common.type]{\tcode{std::common_type} specializations}

\begin{codeblock}
template<std::@\libconcept{delta_quantity}@ Q1, std::@\libconcept{delta_quantity}@ Q2>
  requires requires {
    { std::get_common_reference(Q1::reference, Q2::reference) } -> std::@\libconcept{quantity_reference}@;
    typename std::common_type_t<typename Q1::rep, typename Q2::rep>;
    requires std::@\libconcept{representation_of}@<std::common_type_t<typename Q1::rep, typename Q2::rep>,
                                         std::get_common_quantity_spec(Q1::quantity_spec,
                                                                            Q2::quantity_spec)>;
  }
struct @\libspec{std::common_type}{quantity}@<Q1, Q2> {
  using type = std::quantity<std::get_common_reference(Q1::reference, Q2::reference),
                                  std::common_type_t<typename Q1::rep, typename Q2::rep>>;
};

template<std::@\libconcept{delta_quantity}@ Q, std::@\libconcept{Representation}@ Value>
  requires(Q::unit == std::one) && requires {
    typename std::quantity<Q::reference, std::common_type_t<typename Q::rep, Value>>;
  }
struct @\libspec{std::common_type}{quantity}@<Q, Value> {
  using type = std::quantity<Q::reference, std::common_type_t<typename Q::rep, Value>>;
};
\end{codeblock}

\ednote{
P3045 also proposes a \tcode{std::numeric_limits} specialization for \tcode{quantity}.
Its wording is not written yet and belongs in a new subclause [qty.numeric.limits] after this one.
\tcode{quantity} is a non-arithmetic standard type, and [numeric.limits.general] paragraph 6
says: ``Non-arithmetic standard types, such as \tcode{complex<T>}, shall not have specializations.''
The specialization needs that paragraph to be amended at the same time, for example by inserting
``, except as specified in [qty.numeric.limits]'' before its final period.
}

\rSec1[qty.systems]{Systems}

\rSec1[qty.chrono]{\tcode{std::chrono} interoperability}

\begin{codeblock}
namespace std {

template<typename Period>
consteval auto @\exposidnc{time-unit-from-chrono-period}@()
{
  using namespace si;

  if constexpr (is_same_v<Period, std::chrono::nanoseconds::period>)
    return nano<second>;
  else if constexpr (is_same_v<Period, std::chrono::microseconds::period>)
    return micro<second>;
  else if constexpr (is_same_v<Period, std::chrono::milliseconds::period>)
    return milli<second>;
  else if constexpr (is_same_v<Period, std::chrono::seconds::period>)
    return second;
  else if constexpr (is_same_v<Period, std::chrono::minutes::period>)
    return minute;
  else if constexpr (is_same_v<Period, std::chrono::hours::period>)
    return hour;
  else if constexpr (is_same_v<Period, std::chrono::days::period>)
    return day;
  else if constexpr (is_same_v<Period, std::chrono::weeks::period>)
    return mag<7> * day;
  else
    return mag_ratio<Period::num, Period::den> * second;
}

template<typename Rep, typename Period>
struct @\libspec{quantity_like_traits}{std::chrono::duration}@<std::chrono::duration<Rep, Period>> {
  static constexpr auto reference = @\exposidnc{time-unit-from-chrono-period}@<Period>();
  static constexpr bool explicit_import = false;
  static constexpr bool explicit_export = false;
  using rep = Rep;
  using @\exposidnc{T}@ = std::chrono::duration<Rep, Period>;

  static constexpr rep to_numerical_value(const @\exposid{T}@& q) noexcept(
    std::is_nothrow_copy_constructible_v<rep>)
  {
    return q.count();
  }

  static constexpr @\exposidnc{T}@ from_numerical_value(const rep& v) noexcept(
    std::is_nothrow_copy_constructible_v<rep>)
  {
    return @\exposidnc{T}@(v);
  }
};

template<typename Clock>
struct @\libglobal{chrono_point_origin_}@ final : absolute_point_origin<isq::time> {
  using clock = Clock;
};

template<typename Clock, typename Rep, typename Period>
struct @\libspec{quantity_like_traits}{std::chrono::time_point}@<
  std::chrono::time_point<Clock, std::chrono::duration<Rep, Period>>> {
  static constexpr auto reference =
    point<@\exposidnc{time-unit-from-chrono-period}@<Period>(), chrono_point_origin<Clock>>;
  static constexpr bool explicit_import = false;
  static constexpr bool explicit_export = false;
  using rep = Rep;
  using @\exposidnc{T}@ = std::chrono::time_point<Clock, std::chrono::duration<Rep, Period>>;

  static constexpr rep to_numerical_value(const @\exposid{T}@& tp) noexcept(
    std::is_nothrow_copy_constructible_v<rep>)
  {
    return tp.time_since_epoch().count();
  }

  static constexpr @\exposidnc{T}@ from_numerical_value(const rep& v) noexcept(
    std::is_nothrow_copy_constructible_v<rep>)
  {
    return @\exposidnc{T}@(std::chrono::duration<Rep, Period>(v));
  }
};

}
\end{codeblock}
